% ======================================================================
%  Constrained-worldline brachistochrones in non-stationary spacetimes
%  IOP iopart class -- Classical and Quantum Gravity
% ======================================================================
\documentclass[12pt]{iopart}

% iopart pre-defines \equation*, which clashes with amsmath; free it first
\makeatletter
\expandafter\let\csname equation*\endcsname\relax
\expandafter\let\csname endequation*\endcsname\relax
\makeatother
\usepackage{amsmath}
\usepackage{iopams}
\usepackage{graphicx}
\usepackage{bm}
\usepackage{xcolor}
\usepackage{hyperref}
\hypersetup{colorlinks=true, linkcolor=blue, citecolor=blue, urlcolor=blue}

% all figures collected here as PNG
\graphicspath{{Immagini/}}

% prose placeholder helper
\newcommand{\todo}[1]{{\color{red}\small[[#1]]}}

% shortcuts
\newcommand{\Ehat}{\hat E}
\newcommand{\rp}{r_{+}}
\newcommand{\re}{r_{e}}
\newcommand{\Jc}{J_{c}}
\newcommand{\Af}{A_{\mathrm{freeze}}}
\newcommand{\dd}{\mathrm{d}}

\begin{document}

\title{Constrained-worldline brachistochrones in non-stationary
spacetimes: Kodama rails, closed forms, and plunge inversion}

\author{Iman Rosignoli}
\address{University of Pavia, Pavia, Italy}
\ead{iman.rosignoli01@universitadipavia.it}

\begin{abstract}
We extend the relativistic brachistochrone (fastest-descent constrained rail
curves, Perlick 1991) from stationary to non-stationary spacetimes. The rail
constraint $-u\cdot W=\Ehat$ is an actively controlled invariant, legitimate via
Pontryagin's principle, with the selector $W$ following a hierarchy Killing $\to$
conformal Killing $\to$ Kodama vector; we treat FLRW, Vaidya ($W=\partial_v$, the
Kodama vector) and Thakurta--Kerr. Main results: Kodama energy conservation along
the rail; a plunge-radius law and a branch-intrinsic depth ordering (rotation and
the conformal factor invert the plunge under a same-launch comparison, with no
evaporative inversion at the frozen-clock level); closed-form equatorial
separatrices in Weierstrass functions, continued through the ergosphere, with a
trichotomy and an analytic cusp theorem; an existence theorem for the conformal
plunge inversion; and first-order adiabatic corrections whose \emph{on-shell}
component reduces in closed form to length-two iterated Abelian integrals on the
frozen spectral curve, unified across the three geometries by a breathing-indicatrix
structure. The \emph{off-shell} costate piece, of the same class, is assembled in
closed form on all four branches and verified against the true non-autonomous flow
to $O(\varepsilon^2)$; its source admits a universal form
($\lambda+\int E H_E\,\dd\lambda$ conformal, $[\,rp_r\,]-\lambda$ Vaidya).
\emph{Status.} The controlled-rail formulation, the closed-form separatrices and
cusp theorem, and the on-shell reduction are proved; the fixed-endpoint no-inversion
is proved for $r_0\le R_\ast(E,a)$, machine-certified at representative $r_0$,
conjectural in general; interior-ergosphere optimality, the outgoing-Vaidya problem,
and a rigorous higher-genus polylogarithm framework remain open. Only the canonical
\emph{naming} of the genus-two weight-two letters (single-valued completion,
irreducibility, basis minimality) is deferred; block coefficients are explicit
symbolic, the iterated integrals evaluated by a convergent theta-nome ($q$-)series.
Algebraic identities are independently confirmed in Mathematica; genus-two
theta/nome evaluations use Sage with \texttt{abelfunctions}.
\end{abstract}

\maketitle

% widen ToC number box so the wide "Appendix A" label does not overlap titles
\makeatletter
\renewcommand{\numberline}[1]{#1\enspace}
\makeatother
{\hypersetup{linkcolor=black}\tableofcontents}


% ======================================================================
\section{Introduction}
\label{sec:intro}
% ----------------------------------------------------------------------
The relativistic brachistochrone---the fastest constrained ``rail'' worldline from a
given launch event to a fixed spatial target, at a \emph{free} value of the arrival
clock---was formulated by Perlick for stationary spacetimes
\cite{Perlick1991,PerlickRayOptics}, where a timelike Killing vector provides
the conserved rail energy. It connects naturally to Fermat's principle and to
optical (Randers--Zermelo--Finsler) geometry
\cite{Kovner1990,Randers1941,GibbonsWerner2008,GibbonsRanders2009,BaoRoblesShen}.
Building on our earlier relativistic-brachistochrone studies---the stationary
Kerr geometry \cite{RosignoliKerr}, the interior Schwarzschild metric
\cite{RosignoliIntSchw}, and Tolman--Oppenheimer--Volkoff stellar interiors
\cite{RosignoliTOV}---we extend the construction to \emph{non-stationary}
spacetimes: FLRW, Vaidya \cite{Vaidya1951}, and Thakurta--Kerr
\cite{Thakurta1981}.
These extensions raise three questions that organize the paper.
\emph{(i) Is the rail constraint legitimate without a Killing vector?} We show it
is: $-u\cdot W=\Ehat$ is an actively controlled invariant, and the fastest
constrained worldline solves a well-posed time-optimal control problem
(Sec.~\ref{sec:principle}) whose control cost measures precisely the missing
symmetry. \emph{(ii) What replaces the conserved energy?} In a spherically
symmetric dynamical spacetime the selector $W$ is the Kodama vector, whose flux
charge is the Misner--Sharp mass; the rail conserves the Kodama energy
identically even though it is \emph{not} Noether-conserved for free fall
(Sec.~\ref{sec:vaidya}). \emph{(iii) Which phenomena are genuinely
non-stationary?} We isolate three: capture by cosmological expansion, the
endpoint-sensitive (future-anticipating) turning point in Vaidya, and the breathing
freezing surface that eats the radius of convergence of the quasi-constants in
Thakurta--Kerr. The three spacetimes form a ladder of increasing structure---%
FLRW (degenerate base), Vaidya (dynamical horizon), Thakurta--Kerr
(rotation $+$ conformal expansion)---treated in
Secs.~\ref{sec:flrw}--\ref{sec:thakurta}; equatorial closed forms and the
plunge-inversion analysis follow in Secs.~\ref{sec:closed}--\ref{sec:bvp}.

To keep the reader oriented about what is \emph{proved} versus \emph{verified}
versus \emph{conjectural}, Table~\ref{tab:claims} states each principal claim with
its status and the supporting evidence; the adiabatic first-order result in
particular is deliberately split into three separate claims (on-shell closed form;
exact off-shell quadrature; and the explicit off-shell assembly as depth-two
iterated Abelian integrals plus elementary pieces on all four branches---with only
the canonical higher-genus-polylogarithm \emph{naming} (single-valued completion,
irreducibility, and basis minimality) left open).

\begin{table}[htbp]
    \centering
    \caption{Principal claims, their status, and supporting evidence. ``Proved''
    means an analytic proof in the text; ``machine-certified'' means a
    computer-assisted proof over a stated domain (interval arithmetic);
    ``verified'' means checked numerically to the quoted residual; ``conjectural''
    is stated as such.}
    \label{tab:claims}
    \small
    \begin{tabular}{p{0.40\textwidth}p{0.20\textwidth}p{0.30\textwidth}}
        \hline\hline
        Claim & Status & Evidence\\
        \hline
        Controlled-rail (Kodama) formulation; $-u\cdot W=\Ehat$ from Pontryagin
            & proved & Sec.~\ref{sec:principle}, Eq.~\eqref{eq:transversality}\\
        Kodama energy conserved along the Vaidya rail
            & proved & Sec.~\ref{sec:vaidya}, Fig.~\ref{fig:kodama}\\
        Closed-form equatorial separatrices; ergosphere cusp theorem
            & proved & Sec.~\ref{sec:closed}\\
        On-shell first-order correction $=$ length-two iterated Abelian integrals
            & proved (reduction) & Sec.~\ref{sec:thakurta}, Tab.~\ref{tab:script-map}\\
        Exact \emph{off-shell} first-order term (extended-Hamiltonian quadrature)
            & derived $+$ verified & \ref{app:firstorder};
              Figs.~\ref{fig:phi-true-dynamic},~\ref{fig:vaidya-offshell}
              (slope $\to2$)\\
        Explicit off-shell assembly $=$ depth-two iterated Abelian integrals
        $+$ elementary (all four branches)
            & assembled $+$ verified & \ref{app:firstorder},
              Tab.~\ref{tab:script-map}; canonical polylog naming deferred
              (row below)\\
        Fixed-endpoint no-inversion
            & tiered: proved / machine-certified / conjectural
            & Sec.~\ref{sec:bvp}, Tab.~\ref{tab:protocol}\\
        Genus-two irreducibility \& minimal basis dimension
            & conjectural & \S\ref{sec:thakurta}\\
        Interior-ergosphere PMP optimality; outgoing-Vaidya BVP
            & open & analytic continuation only\\
        \hline\hline
    \end{tabular}
\end{table}


% ======================================================================
\section{The controlled-rail principle}
\label{sec:principle}
% ----------------------------------------------------------------------

\subsection{Rail constraint and the indicatrix}
\label{sec:indicatrix}
The rail is a constrained worldline with the invariant
\begin{equation}
    -\,u\cdot W \;=\; \Ehat ,
    \label{eq:rail}
\end{equation}
actively maintained by an ideal rail force $f_\mu=Du_\mu/\dd\tau$ with
$f\cdot u=0$ (automatic for a rest-mass--preserving force), which creates no rest
mass but is free to exert power against the selector (see below).

\begin{table}[htbp]
    \centering\small
    \caption{Notation used throughout. Rail charges $(\Ehat,J)$ are conserved control
    costates; the ``eff'' quantities are their optical/conformal normalizations.}
    \label{tab:notation}
    \begin{tabular}{@{}ll@{}}
        \hline\hline
        symbol & meaning\\
        \hline
        $\Ehat=-u\cdot W$ & rail energy invariant (conserved control charge)\\
        $J=p_\varphi$ & axial control costate (conserved; \emph{not} $g(u,\partial_\varphi)$)\\
        $L_{\rm mech}=g(u,\partial_\varphi)$ & mechanical angular momentum (\emph{not} conserved on the rail)\\
        $A(\eta)$ & Thakurta conformal factor, $g=A^2 g_{\rm Kerr}$\\
        $m(v)$ & Vaidya Misner--Sharp mass\\
        $a$ & Kerr spin\\
        $E_{\rm eff}=\Ehat/A$, $J_{\rm eff}=J/A$ & conformally normalized charges (weight $A^{-1}$)\\
        $P_r=p_r/A$ & normalized radial momentum (conformally symplectic, not canonical)\\
        $t,\tau,\eta,v$ & coordinate, proper, conformal, advanced-null clocks\\
        $s$ & the selected evolution clock / slow parameter\\
        $\lambda$ & elapsed clock from the free-arrival endpoint\\
        $\varepsilon$ & adiabatic rate ($A'/A$ for Thakurta, $\dot m$ for Vaidya)\\
        \hline\hline
    \end{tabular}
\end{table}

\paragraph{The ellipse.} Fix an event and a time function adapted to $W$. The
linear constraint~\eqref{eq:rail} fixes the $W$-component of $u$; the mass shell
$u\cdot u=-1$ is a quadric; their intersection, projected onto the physical
velocity plane $(\dd r,\,r\,\dd\varphi)$ per unit selector-time, is an ellipse
\begin{equation}
    \dot x(\theta)=c+R(\theta),\qquad \theta\in[0,2\pi),
    \label{eq:ellipse}
\end{equation}
whose center $c$ is the Zermelo wind (zero for FLRW, ingoing radial for Vaidya,
azimuthal---frame dragging---for Thakurta--Kerr) and whose axes $R$ shrink with
the local speed as the freezing surface is approached (Fig.~\ref{fig:indicatrici}).
Every cost clock---arrival/coordinate time $t$, proper time $\tau$, conformal time
$\eta$---shares this \emph{same} indicatrix; the branches differ only in the cost
one-form, i.e.\ in which clock measures the elapsed travel time. Two of these are
\emph{not} independent spatial families: since $\dd t=A(\eta)\,\dd\eta$ with $A>0$ is
a monotone reparametrization, the coordinate-time and conformal-time functionals are
extremized by the \emph{same} spatial path ($\arg\min t_f=\arg\min\eta_f$), so $\eta$
is the canonical conformal-time \emph{representation} of the arrival branch rather
than a third independent brachistochrone. The genuinely distinct branch is the proper
time $\tau$, whose cost carries the extra $\sqrt{1-v^2}$ weight; for Vaidya the native
arrival clock is the advanced null time $v$, again distinct from $\tau$.

\paragraph{Domain of the indicatrix.} The intersection~\eqref{eq:ellipse} is a
compact, strictly convex oval---the hypothesis under which the Pontryagin reduction
below applies---precisely where the selector $W$ is \emph{timelike} and the quadric
intersection is nondegenerate. Three loci require separate treatment. (i) At the
freezing surface the oval collapses to a point, $R\to0$ (the strict convexity
degenerates; this is the turning-point/separatrix locus). (ii) At the horizon the
radial metric coefficient degenerates. (iii) For Thakurta--Kerr inside the
\emph{ergosphere}, $\partial_\eta$---hence $W$---becomes spacelike, so
$-u\cdot W=\Ehat$ fixes a \emph{spatial} component of $u$ and the intersection with
the mass shell is no longer automatically a compact oval in the PMP sense. The local
theorems of this section are therefore stated in the timelike-$W$, nondegenerate
domain. The interior expressions are the \emph{analytic continuation} of these
closed forms (Sec.~\ref{sec:trichotomy}, where the penetrating $t$-branch is analyzed
explicitly); we do \emph{not} claim they are proved optimal-control solutions inside
the ergosphere. A genuinely different, and inequivalent, construction---re-anchoring
the constraint $-u\cdot W=\Ehat$ to a timelike congruence (Doran/ZAMO)---would change
the selector $W$, the indicatrix, the cost, and the Hamiltonian, and is left to
future work; a coordinate change alone (e.g.\ to Doran variables) does not alter the
geometric selector and hence cannot substitute for it.

\paragraph{The optimal-control problem.} Precisely, the brachistochrone is the
following time-optimal control problem. The state is the spatial position
$x=(r,\varphi)$ carried along the worldline; the evolution parameter is the selected
clock $s\in\{t,\tau,\eta,v\}$; the control is the unit spatial direction $\theta\in
S^1$ of the velocity on the rail indicatrix~\eqref{eq:ellipse}, with dynamics
$\dd x/\dd s=\dot x(\theta;x,s)$ read off that indicatrix. The boundary data are a
\emph{fixed initial event} $(s_0,x_0)$ and a \emph{fixed spatial target} $x_f$ (a
point, or a target set/worldline), with the arrival clock $s_f$ left \emph{free}; the
cost is the elapsed clock $s_f-s_0$. Because $s_f$ is free, the Pontryagin
transversality condition is $H|_{s=s_f}=0$ [Eq.~\eqref{eq:transversality}]---a free
final-clock condition, \emph{not} a two-event condition. We accordingly reserve
``fixed endpoints'' for the two-point boundary-value problem of Sec.~\ref{sec:bvp}, in
which both spatial endpoints \emph{and} the arrival are prescribed; the default problem
fixes the spatial target at a free arrival clock.

\paragraph{Pontryagin reduction.} Minimizing travel time subject
to~\eqref{eq:ellipse} is a time-optimal control problem with control $\theta$.
The velocity indicatrix~\eqref{eq:ellipse} is a compact, strictly convex
\emph{closed curve} (the boundary of a convex region), and the control is its
direction $\theta\in S^1$, itself a compact one-dimensional set. Maximizing the
support function $p\cdot\dot x(\theta)$ over this curve has a \emph{unique} maximizer
by strict convexity, so \emph{no} Filippov relaxation is needed---the control set
$S^1$ is already compact and the maximizer single-valued; we do \emph{not} claim the
velocity locus is a convex \emph{subset} of the plane, nor convexify it (which would
admit averaged velocities violating the unit-mass shell). Excluding $p=0$ by the
nontriviality condition, the optimal control $\theta^\star(x,p)$ is single-valued and
smooth (away from the freezing and horizon degeneracies, treated separately) and
Pontryagin's
maximum principle \cite{Pontryagin,Liberzon} delivers the branch Hamiltonian
$H(x,p)=\max_\theta p\cdot\dot x(\theta)$ in closed form
(Secs.~\ref{sec:vaidya}--\ref{sec:thakurta}). This is a genuine time-optimal
control problem, for which the Pontryagin (vakonomic-type) reduction is the correct
prescription; we do not invoke the nonholonomic d'Alembert principle, whose
equations differ from the vakonomic ones unless the constraint distribution is
integrable---a coincidence that convexity does not entail. Two structural facts hold throughout: $p_\varphi=J$ is conserved (axial symmetry),
and the free-arrival-time transversality fixes the branch Hamiltonian to vanish
\emph{at the free endpoint},
\begin{equation}
    H\big|_{s=s_f}=0 \qquad\text{(free arrival time).}
    \label{eq:transversality}
\end{equation}
For a non-autonomous system $H$ is \emph{not} a global first integral: along the
extremal $\dd H/\dd s=\partial_s H$, proportional to $m'(v)$ (Vaidya) or $A'(\eta)$
(Thakurta--Kerr), so $H\neq0$ in the interior---consistently, $p_v(r_1)=0$ at the
free arrival radius yet $p_v(r)\neq0$ along the path (Fig.~\ref{fig:vaidya-traj}).
The globally conserved object is the extended Hamiltonian $\mathcal H=p_s+H=0$ on
the enlarged phase space, with $p_s$ conjugate to the evolution parameter $s$; only
in the stationary limit is $H$ autonomous and $H\equiv0$ a genuine first integral.

\paragraph{Cost of control.} For an accelerated worldline with rail four-force
$a^\mu=Du^\mu/\dd\tau$, differentiating the invariant gives the \emph{exact}
identity
\begin{equation}
    \frac{\dd(-u\cdot W)}{\dd\tau}=-\,a\cdot W-u^a u^b\,\nabla_{(a}W_{b)},
    \label{eq:cost}
\end{equation}
whose two terms play distinct roles. The constraint $a\cdot u=0$ holds
\emph{automatically} for any four-force that preserves the mass shell
$u\cdot u=-1$ (it is the derivative of the shell condition); it is \emph{not} a
special ``magnetic'' property and does not by itself imply that the rail performs no
work relative to a physical observer. It means only that no rest mass is created;
the rail is nonetheless free to exert power \emph{against the selector} $W$ (the term
$a\cdot W$). One should distinguish (i) preservation of rest mass ($a\cdot u=0$),
(ii) power measured by a chosen observer congruence, (iii) the constraint force
itself, and (iv) the additional power needed to compensate non-stationarity.
A free geodesic ($a=0$) would let $-u\cdot W$ drift at $-u^au^b\nabla_{(a}W_{b)}$,
which vanishes for all $u$ iff $W$ satisfies the Killing equation; to hold
$-u\cdot W=\Ehat$ fixed the rail supplies exactly $a\cdot W=-u^au^b\nabla_{(a}W_{b)}$,
so that the left-hand side of~\eqref{eq:cost} is zero. This control power equals the
failure of the Killing equation---the missing symmetry---which is why the
non-stationary rail is legitimate but not free: $\Ehat$ is actively controlled, not
conserved by a symmetry. The optical reductions of $H$
reproduce the Randers/Zermelo structure
\cite{GibbonsRanders2009,CaponioWind}.

\paragraph{Operational reading: thrust and fuel.} Since $a^\mu=Du^\mu/\dd\tau$ is the
\emph{proper} acceleration---the covariant derivative already contains gravity, so a
free-fall geodesic reads $a=0$ on an accelerometer---the rail force is exactly the
thrust per unit mass a rocket must supply to fly the brachistochrone. Its magnitude
$|a|=\sqrt{a\cdot a}\ge0$ ($a$ is spacelike since $a\cdot u=0$) integrates to a
characteristic velocity $\Delta v=\int|a|\,\dd\tau$ and hence, \emph{heuristically}
and under a chosen constant-exhaust rocket model (not needed for the mathematical
argument below), to a fuel fraction $1-\exp(-\Delta v/v_{\rm e})$ for exhaust
speed $v_{\rm e}$. For a representative Schwarzschild arc ($E=1.4$, $r:10M\!\to\!5.7M$)
one finds $|a|\sim0.03$--$0.2/M$ and $\Delta v\approx0.6\,c$ (a photon rocket burns
$\sim\!45\%$ of its mass; a chemical one cannot); the proper acceleration diverges as
the freezing/horizon surface is approached, so the fuel budget diverges there
(\texttt{brachistochrone\_thrust\_fuel.py}). The brachistochrone is the
\emph{fastest} rail curve, not the fuel-minimal one---an unpowered free-fall geodesic
(after launch) requires no rail thrust---so this is the price of speed at fixed $\Ehat$; in the non-stationary case the
$a\cdot W$ term adds the thrust that fights the accretion/evaporation drift.

\paragraph{Geometric dictionary.} The indicatrix~\eqref{eq:ellipse} carries two
independent data---its \emph{center} $c=(c_r,c_\varphi)$ (the Zermelo wind) and
the \emph{eccentricity} of its semi-axes (the anisotropy of the local speed,
$a_r\neq a_\varphi$)---and each governs a distinct feature of the
brachistochrones, the same for every branch:
\begin{itemize}\itemsep2pt
\item eccentricity $\neq0$ (a spatial gradient $\partial_r f\neq0$)
$\Rightarrow$ the $t$ and $\tau$ branches split in depth, $n_t/n_\tau=E/f\neq1$;
\item azimuthal offset $c_\varphi\neq0$ (frame dragging)
$\Rightarrow$ rotational plunge inversion becomes possible;
\item radial offset $c_r\neq0$ (a non-stationary wind, $m'\neq0$)
$\Rightarrow$ endpoint-sensitive memory, $p_v\neq0$ along the path;
\item shrinking axes $R\to0$ (local speed $v\to0$) $\Rightarrow$ freezing and
capture.
\end{itemize}
Eccentricity and center offset are \emph{separate} levers: the first switches on
the $t/\tau$ split, the second the inversions. In particular the three branches
\emph{degenerate}---the $t$, $\tau$ and $\eta$ brachistochrones coincide---iff the
indicatrix is, up to a position-independent conformal scale, a rigid centered
circle ($c=0$, $a_r=a_\varphi$, radius depending on the flow parameter alone),
equivalently iff the optical cost ratios $n_t/n_\tau,\,n_t/n_\eta$ are
position-independent. This is exactly spatial homogeneity: FLRW is the unique
such case in our ladder (Sec.~\ref{sec:flrw}), whereas any gradient or wind
breaks it---even the static, non-rotating Schwarzschild ellipse, which is
centered ($c=0$) yet eccentric, already splits $t$ from $\tau$. The indicatrix
shape is thus the organizing thread of the paper: FLRW (circle) $\to$ Vaidya
(radial wind) $\to$ Thakurta--Kerr (angular wind) is a progression of ellipse
deformations, each switching on one further phenomenon.

\begin{figure}[tbp]
    \centering
    \includegraphics[width=\textwidth,height=0.68\textheight,keepaspectratio]{Immagini/fig_indicatrici.png}
    \caption{Rail indicatrices (linear constraint $\cap$ normalization
    $=$ ellipse): the unifying formalism. FLRW gives shrinking circles (no
    wind, no eccentricity $\Rightarrow$ degenerate branches); Vaidya an ingoing
    radial wind ($f-E^2<0$, driven by $m(v)$ $\Rightarrow$ endpoint-sensitive memory);
    Thakurta--Kerr an angular wind (frame dragging $\Rightarrow$ rotational
    inversion), the ellipse shrinking with $A(\eta)$. Eccentricity switches on
    the $t/\tau$ split, the center offset the inversions.}
    \label{fig:indicatrici}
\end{figure}

\subsection{Existence, normality, and global optimality}
\label{sec:control-foundations}
The Pontryagin reduction above supplies only \emph{necessary} conditions. We
record here the three complementary facts that make the constrained problem well
posed for the non-stationary rail: an existence theorem for the minimizer, the
normality of every regular extremal (so the reduction $-u\cdot W=\Ehat$ is never
vacuous), and a Hamilton--Jacobi certificate that upgrades a saturated extremal to
a \emph{global} optimum. All three are elementary under the standing hypothesis
that the indicatrix is a compact, uniformly convex oval [Eq.~\eqref{eq:ellipse}].

\paragraph{Existence (direct method).} Let the spatial arc range over a compact,
regular domain $\Omega$ on which the rail one-form and the metric are smooth and
the arrival-time indicatrix obeys the uniform two-sided bound
\begin{equation}
  \rho\,|v| \le F(x,v) \le R\,|v|, \qquad 0<\rho\le R<\infty,
  \label{eq:coercive}
\end{equation}
with $\rho,R$ the inradius and circumradius of the oval. Then the arrival time
$T[\gamma]=\int F(\gamma,\dot\gamma)\,\dd\sigma$ attains its infimum in the class
of Lipschitz arcs joining the fixed endpoints. The left inequality gives
coercivity and a uniform Lipschitz---hence, by Arzel\`a--Ascoli, uniformly
convergent---minimizing sequence; convexity of $F$ in $v$ gives sequential lower
semicontinuity (Tonelli); and the non-autonomous coupling to the evolution
parameter $s$ (through $m(v)$ or $A(\eta)$) is controlled by ODE stability and a
Gr\"onwall estimate, so the limit arc is admissible and minimizing. The two-sided
bound~\eqref{eq:coercive} fails exactly at the freezing/horizon surfaces, where
$\rho\to0$ and the fuel budget of Sec.~\ref{sec:indicatrix} diverges; the theorem
therefore applies on any compact subdomain bounded away from them.

\paragraph{Normality.} Every regular minimizer is a \emph{normal} extremal---the
abnormal multiplier $p_0$ cannot vanish. If $p_0=0$ the maximized Hamiltonian
would read $\max_v p\cdot v=0$ with $p\neq0$; but the support function of the oval
obeys $h(p)=\max_{v\in\partial B}p\cdot v\ge\rho\,|p|>0$, a contradiction. Hence
$p_0$ normalizes to $-1$, the costate $p$ never vanishes, and the reduction to the
branch Hamiltonian $H=\max_\theta p\cdot\dot x(\theta)$ together with the rail
invariant $-u\cdot W=\Ehat$ is legitimate along the entire extremal. This is the
control-theoretic content of ``excluding $p=0$'' used in the Pontryagin reduction
of Sec.~\ref{sec:indicatrix}: nontriviality is not an assumption but a consequence
of $\rho>0$.

\paragraph{Global optimality (Hamilton--Jacobi certificate).} The maximum
principle is necessary but not sufficient: a PMP extremal is a \emph{candidate}
minimizer, guaranteed only \emph{locally} optimal. A sufficient global certificate
comes from dynamic programming. In the \emph{stationary} (frozen) problem the value
function solves the eikonal Hamilton--Jacobi--Bellman equation
$F^\ast(x,\nabla V)=1$ ($F^\ast$ the dual/optical norm); but the rail is
\emph{non-autonomous} (Vaidya $m(v)$, Thakurta--Kerr $A(\eta)$), so the value function
depends on the evolution clock $s$ and obeys the time-dependent equation
\begin{equation}
  \partial_s V+\min_{u}\big\{\,1+\nabla_x V\cdot f(x,u,s)\,\big\}=0
  \label{eq:hjb}
\end{equation}
(equivalently, an autonomous HJB on the state extended by $s$). This is a
\emph{verification theorem}: \emph{if} a Lipschitz viscosity subsolution $V$ exists
and a candidate extremal \emph{saturates} it ($\nabla V\cdot\dot x=F$ pointwise), then
that extremal realizes the value function and is a \emph{global} minimizer on the
single-valued region of $V$. We do \emph{not} construct such a global $V$ for generic
non-stationary Kerr/Vaidya trajectories; in the stationary limit the closed-form
optical (Randers) metric of Secs.~\ref{sec:vaidya}--\ref{sec:thakurta} supplies one on
any region carrying a single-valued $V$, and elsewhere our statements are those of PMP
extremals with local checks. Where $V$ turns multivalued the certificate localizes, as
follows.

\paragraph{Conjugate points, Maxwell set, cut locus.} Two mechanisms end
optimality. A \emph{conjugate point} (Jacobi) ends \emph{local} minimality when
the second variation degenerates; up to the first conjugate point the extremal is
a strict local minimum. A \emph{Maxwell point}---two distinct extremals of equal
arrival time---ends \emph{global} minimality earlier and is exactly a shock
(gradient discontinuity) of $V$, i.e.\ a point of the cut locus. Strict convexity
of the indicatrix removes Filippov relaxation but does \emph{not} exclude Maxwell
points: distinct windings past the rotation axis, or the two $\pm J$ arcs reaching
a common target, can tie. Our closed forms are certified between the launch point
and the first such singularity of $V$; beyond it the global optimum is selected by
comparing arrival times across the competing extremals.

\paragraph{What is \emph{not} claimed inside the ergosphere.} Where the Killing
selector $\xi$ becomes spacelike---inside the ergosphere---the invariant
$-u\cdot\xi=\Ehat$ no longer bounds a timelike-future cost and the indicatrix oval
degenerates: the closed-form $\varphi(r)$ continued through $\re$
(Sec.~\ref{sec:trichotomy}) is a genuine \emph{analytic continuation} of the
extremal, not a certified physical minimum. The existence and Hamilton--Jacobi
statements above are asserted only outside the ergosphere, on the region where $F$
is a bona fide convex indicatrix; inside, a spacelike (Kodama-type) selector would
be needed to restore a timelike cost, which we do not develop. Freezing (where
$\rho\to0$), turning (where a semi-axis vanishes), and the separatrix (a double
root of the radical) are distinct degeneracies and are not conflated.

\subsection{Selecting $W$: the Killing--CKV--Kodama hierarchy}
\label{sec:whierarchy}
The selector $W$ is fixed by a hierarchy of decreasing symmetry, each level
recovering the previous one in the appropriate limit.

\emph{(a) Killing.} If a timelike Killing vector $\xi$ exists, $W=\xi$ and the
right-hand side of Eq.~\eqref{eq:cost} vanishes identically: $-u\cdot\xi$ is
Noether-conserved and the rail is free. This is Perlick's stationary case
\cite{Perlick1991}.

\paragraph{Recovery of Perlick's stationary optics.} In the Killing limit
$W=\xi=\partial_t$ the controlled-rail arrival-time Finsler function is exactly a
Randers metric, so the construction reduces to Perlick's fixed-energy optics rather
than merely resembling it. Solving the rail $-u\cdot\xi=\Ehat$ jointly with
$g(u,u)=-1$ for the coordinate-time rate $t'=\dd t/\dd\sigma$ gives a quadratic whose
future-directed root is
\begin{equation}
F_{\rm rail}=\beta_a v^a+\sqrt{a_{ab}\,v^a v^b},\quad
\beta_a=\frac{g_{ta}}{f},\quad
a_{ab}=\frac{\Ehat^2\bigl(g_{ta}g_{tb}+f\,g_{ab}\bigr)}{f^2(\Ehat^2-f)},
\label{eq:perlick-randers}
\end{equation}
with $f=-g_{tt}$ and $v^a=\dd x^a/\dd\sigma$. In the threading form
$g=-f(\dd t-\omega_a\dd x^a)^2+h_{ab}\dd x^a\dd x^b$ one has
$g_{ta}g_{tb}+f\,g_{ab}=f\,h_{ab}$, so $\beta_a=\omega_a$ and
$a_{ab}=\Ehat^2 h_{ab}/[f(\Ehat^2-f)]$---precisely Perlick's fixed-energy stationary
optical metric \cite{Perlick1991}. The null/high-energy limit $\Ehat\to\infty$ gives
$a_{ab}\to h_{ab}/f$, the standard Fermat--Randers optical metric of Kerr
\cite{GibbonsWerner2008}. This is the functional equality $F_{\rm rail}=F_{\rm Perlick}$
(stronger than coincidence of the Hamiltonian flows on $H=0$), verified against the
directly solved $F_{\rm rail}$ on equatorial Kerr to machine precision.

\emph{(b) Conformal Killing.} If the spacetime is conformally stationary,
$g=\Omega^2\hat g$ with $\hat g$ stationary, the timelike conformal Killing
vector $\partial_\eta$ selects the conformal-time branch; $-u\cdot W$ is exactly
conserved for null rays and controlled up to the conformal factor for massive
rails. FLRW and constant-$A$ Thakurta--Kerr live at this level.

\emph{(c) Kodama.} In a spherically symmetric \emph{dynamical} spacetime with no
Killing or conformal Killing vector, the Kodama vector
$K^a=\varepsilon^{ab}\nabla_b r$ \cite{Kodama1980,Hayward1996} is the canonical
replacement: it is divergence-free (Eq.~\eqref{eq:kodama}), reduces to $\xi$ in
the static limit, and its flux $G^{ab}K_b$ is conserved with charge the
Misner--Sharp mass. Vaidya is the paradigm (Sec.~\ref{sec:kodama-energy}).

\emph{(d) Convention.} Absent even spherical symmetry the Kodama construction
fails and $W$ becomes a declared convention, with no distinguished normalization
of the cost one-form. Levels (a)--(c) are what make the non-stationary rail
canonical rather than arbitrary.


% ======================================================================
\section{FLRW: the degenerate base case}
\label{sec:flrw}
% ----------------------------------------------------------------------
For spatially flat FLRW, $g=a(\eta)^2\eta_{\rm Mink}$, the rail uses the
conformal Killing vector $\partial_\eta$; the local speed is
\begin{equation}
    v(\eta)=\sqrt{1-a(\eta)^2/\Ehat^{2}},
    \label{eq:flrwspeed}
\end{equation}
vanishing at the freezing barrier $a=\Ehat$.
\paragraph{Branch degeneracy.} Because $a=a(\eta)$ depends on conformal time
alone (spatial homogeneity), the optical index built from the conformal Killing
vector,
\begin{equation}
    n(\eta,r)=\frac1v=\frac{\Ehat}{\sqrt{\Ehat^2-a(\eta)^2}},
    \label{eq:flrwindex}
\end{equation}
is \emph{position-independent} at fixed $\eta$. The conformal-time functional
$T_\eta=\int n\,\dd\ell_{\rm flat}$ is then minimized by a flat straight line
(Fermat with a purely time-dependent index). The same curve minimizes $T_t$ and
$T_\tau$: with $\dd t=a\,\dd\eta$ and $\dd\tau=a\sqrt{1-\dot{\vec x}^2}\,\dd\eta$,
the monotone reparametrization $t(\eta)=\int a\,\dd\eta$ maps critical points to
critical points, and the $\tau$-cost differs from the $\eta$-cost only by the
same $\eta$-dependent weight. Hence
\begin{equation}
\begin{split}
    \arg\min T_t&=\arg\min T_\tau=\arg\min T_\eta\\
    &=\text{comoving straight line}.
\end{split}
    \label{eq:flrwdegen}
\end{equation}
The $t/\tau$ splitting that drives every later result is a gravitational-%
\emph{gradient} effect ($\partial_r f\neq0$); it is absent whenever the metric
functions depend on time only. FLRW is thus the degenerate base of the ladder:
it exhibits freezing at $a=\Ehat$ (Fig.~\ref{fig:flrw-kin})
but no branch structure (Fig.~\ref{fig:flrw-var}).

\begin{figure}[tbp]
    \centering
    \includegraphics[width=\textwidth,height=0.68\textheight,keepaspectratio]{Immagini/fig_flrw_cinematica.png}
    \caption{de Sitter FLRW rail kinematics: local speed $v$ and rail redshift
    $E_{\rm phys}=\Ehat/a$; freezing at $a=\Ehat$.}
    \label{fig:flrw-kin}
\end{figure}


\begin{figure}[tbp]
    \centering
    \includegraphics[width=\textwidth,height=0.68\textheight,keepaspectratio]{Immagini/fig_flrw_variazionale.png}
    \caption{Variational check: the comoving straight line minimizes
    \emph{both} $T_t$ and $T_\tau$ (branch degeneracy).}
    \label{fig:flrw-var}
\end{figure}


% ======================================================================
\section{Vaidya: dynamical black hole and the Kodama rail}
\label{sec:vaidya}
% ----------------------------------------------------------------------
The ingoing Vaidya metric $\dd s^2=-f\,\dd v^2+2\,\dd v\,\dd r+r^2\dd\Omega^2$,
$f=1-2m(v)/r$ \cite{Vaidya1951,LindquistSchwartzMisner1965}, has no timelike
Killing vector; its trapping structure is that of a dynamical horizon
\cite{AshtekarKrishnan2003,Booth2005,FaraoniBook}. The Kodama vector
\cite{Kodama1980,Hayward1996,AbreuVisser2010}
\begin{equation}
    K^a=\varepsilon^{ab}\nabla_b r = \partial_v ,\qquad
    \nabla_a K^a=0,
    \label{eq:kodama}
\end{equation}
is divergence-free without any Killing symmetry, and its charge is the
Misner--Sharp mass $m(v)$ \cite{MisnerSharp1964,HaywardMukohyama1999}.

\subsection{Kodama energy is the rail invariant}
\label{sec:kodama-energy}
From the rail constraint alone one finds the exact identity
$-g(T,T)=(f\,u^v-1)^2/\Ehat^2$, hence
\begin{equation}
    -\,u\cdot K \;=\; \Ehat \quad\text{identically along the brachistochrone.}
    \label{eq:kodama-conserved}
\end{equation}
This is \emph{controlled} conservation, not Noether: for a geodesic
$-u\cdot K$ would drift at rate $-m'(u^v)^2/r\propto m'$, exactly the
right-hand side of Eq.~\eqref{eq:cost} for $W=K$. Only the unnormalized
$K=\partial_v$ has zero control cost in the static limit; the normalized
$\hat K=K/\sqrt f$ does not (discriminating test), which is why the Kodama
vector, not a unit selector, is the canonical choice.

\paragraph{Indicatrix and Hamiltonians.} With $w\equiv\Ehat^2-f$ the constraint
$-u_v=\Ehat$ turns the admissible velocities at fixed $(v,r)$ into an ellipse,
\begin{equation}
    r'=(f-\Ehat^2)+\Ehat\sqrt w\,\cos\theta,\qquad
    \varphi'=\frac{\sqrt w}{r}\sin\theta,
    \label{eq:vaidya-ind}
\end{equation}
a Zermelo problem with ingoing radial wind $f-\Ehat^2<0$. Maximizing
$p\cdot\dot x$ over $\theta$ at fixed $p_\varphi=J$ gives the closed-form branch
Hamiltonians
\begin{align}
    H_v&=p_r(f-\Ehat^2)-1+\sqrt w\,\sqrt{\Ehat^2 p_r^2+J^2/r^2},
    \label{eq:Hv}\\
    H_\tau&=p_r(f-\Ehat^2)-\Ehat
        +\sqrt w\,\sqrt{(\Ehat p_r+1)^2+J^2/r^2}.
    \label{eq:Htau-vaidya}
\end{align}
The $\tau$-branch is the $v$-branch under $p_r\to p_r+1/\Ehat$ with unit cost
replaced by $\Ehat$. Non-stationarity is the explicit dependence
$\dd H/\dd v=\partial_v H\propto m'(v)$; the \emph{teleological} (here and
throughout: \emph{endpoint-sensitive}, arising from the two-point boundary-value
problem with a free future arrival---not a causal effect) momentum $p_v=-H$
satisfies $p_v(r_1)=0$ at the free arrival radius yet $p_v(r)\neq0$ along the
path, so the turning point \emph{anticipates future accretion}
(Fig.~\ref{fig:vaidya-traj}).

\paragraph{Branch costate vs mechanical angular momentum.} The conserved $J=p_\varphi$
is the Noether \emph{costate} of the reduced control Hamiltonian---the canonical
momentum conjugate to $\varphi$ in the arrival-time problem---\emph{not} the spacetime
angular momentum $L_{\rm mech}=g(u,\partial_\varphi)$. Although $\partial_\varphi$ is a
Killing vector, Noether's theorem for the \emph{control} symmetry returns the control
costate, not the mechanical momentum: along a forced (rail-constrained, accelerated)
worldline $L_{\rm mech}$ is \emph{not} conserved, since
$\dd L_{\rm mech}/\dd\tau=a\cdot\partial_\varphi\neq0$. The two are therefore distinct
objects. On the native Vaidya $v$-branch shell the Legendre map gives the exact
relation
\begin{equation}
    L_{\rm mech}=r^2 u^\varphi=\frac{\Ehat^2-f(r)}{\Ehat}\,J,
    \label{eq:costate-mech}
\end{equation}
which varies with $r$ while $J$ stays constant (they agree only at the isolated radius
$f=\Ehat^2-\Ehat$); verified from the indicatrix~\eqref{eq:vaidya-ind}. The Legendre
map is moreover \emph{branch-specific}---the $t$-, $\tau$- and $\eta$-costs carry
different cost one-forms---so the conserved costate differs by branch: the native
$v/t$ value is $J$, while the proper-time branch carries
$\tilde p_\varphi=J-A^2b/\Ehat$ (the frame-drag shift of the second-kind clock). We
therefore reserve ``angular momentum'' for the mechanical $g(u,\partial_\varphi)$
throughout and call $J$ the branch costate. The optical rescaling
$J_{\rm eff}=J/A=P_\varphi=p_\varphi/A$ is a normalization of the \emph{costate}
(conformal weight $A^{-1}$), \emph{not} the Kerr-frame mechanical momentum; the
conserved datum is $J$, and $L_{\rm mech}$ is a derived, non-conserved quantity.

\begin{figure}[ht]
    \centering
    \includegraphics[width=\textwidth,height=0.68\textheight,keepaspectratio]{Immagini/fig_kodama_conservazione.png}
    \caption{Kodama energy along the $\tau$-brachistochrone. Top: the
    \emph{residual} $|{-}u\cdot K-\Ehat|$ on a \emph{log} scale, held below
    $6.7\times10^{-16}$---an \emph{actively controlled} (rail-maintained) invariant,
    not a Noether constant. Bottom: the free-particle Kodama drift $\propto m'$ that
    the rail cancels (the missing symmetry).}
    \label{fig:kodama}
\end{figure}

\paragraph{Sanity check: rail energy along the Thakurta--Kerr rail.} The same
test applies to Thakurta--Kerr, where---the geometry being only axisymmetric---no
Kodama vector exists and the rail invariant is instead tied to the conformal
Killing vector $\partial_\eta$ (with factor $A'/A$). Along a \emph{geodesic} the
associated energy $\Ehat=-g^{\rm TK}(u,\partial_\eta)$ drifts at exactly the
conformal rate $\mathrm{d}\Ehat/\mathrm{d}\tau=A'/A$ (verified numerically to
$\sim\!10^{-11}$), whereas along the \emph{brachistochrone}---instantaneously an
on-shell Kerr orbit at $E_{\rm eff}=\Ehat/A(\eta)$---the rail holds $\Ehat$
constant to $8.9\times10^{-16}$, the exact analogue of the Vaidya value
(\texttt{ThakurtaMetric/rail\_conservation.py}). This confirms that the controlled charges
$E_{\rm eff}=\Ehat/A$ (and, by the same conformal weight, $J_{\rm eff}=J/A$)
track along the orbit. Conservation of $\Ehat$ determines how the parameters
\emph{drift}; it does \emph{not} force the non-autonomous trajectory to remain on
the instantaneous $H=0$ shell---the costate equation is independent information.
The frozen shape relation $\mathrm{d}\varphi/\mathrm{d}r=F(r;E_{\rm eff},J_{\rm eff})$
and the parameter-drift correction
$\delta\varphi=-\int(\Ehat\,\partial_E F+J\,\partial_J F)\,\eta\,\mathrm{d}r$
[Eq.~\eqref{eq:deltaphi-EJ}] are therefore the \emph{on-shell component} of the
first-order term; the complete first-order correction is
Eq.~\eqref{eq:adiab-exact}, which adds the off-shell (costate $p_\eta$) piece
[Sec.~\ref{sec:vaidya-adiabatic}, Fig.~\ref{fig:phi-true-dynamic}].


\subsection{Phenomenology: penetration, timing, bounce}
\label{sec:vaidya-pheno}
Three features distinguish the dynamical hole from its frozen-mass
Schwarzschild limit.

\emph{Penetration.} The $\tau$-brachistochrone reaches the apparent horizon
$r=2m(v)$ only if $|J|$ lies below a threshold $J_c(v_0)$; accretion ($m'>0$)
opens a \emph{finite} window $J_c>0$, whereas the negative-rate continuation
($m'<0$ in the \emph{ingoing} metric---a formal continuation; physical evaporation is
the \emph{outgoing} metric, Sec.~\ref{sec:inversion}) closes it to zero
measure---an accretion/(negative-rate) asymmetry with no static analogue. The
finite capture interval $J\in(0,\Jc]$ is the analogue of the Kerr $t$-branch
dichotomy, generated here by the \emph{dynamics}: grazing orbits linger near
periapsis while the growing horizon closes the gap.

\emph{Bounce.} The $r$-parametrization degenerates at periapsis
($\dd v/\dd r\to\infty$); reparametrizing by $v$ makes the Euler--Lagrange
problem regular at the turning point $r'=0$, so the trajectory is integrated
straight \emph{through} the bounce---the same extremal, via the exact change of
variable $(f v'-1)\dd r=(f-r')\dd v$. The $v$-flow reaches the periapsis exactly
where the $r$-flow diverges ($r_{\min}=3.0105$ at $v=7.29$ for $\mu\equiv m'=0.01$),
and the endpoint-sensitive memory $p_v(r)\propto m'$ measures the departure from the
frozen-mass Schwarzschild curve (Fig.~\ref{fig:vaidya-bounce}).

\emph{Timing and the orbit--horizon race.} Because $m$ grows along the path,
$r_{\min}$ depends on \emph{when} the particle arrives, the epoch $v_1$---a
genuinely non-stationary effect absent in Kerr. Strikingly, the absolute and
relative penetration have \emph{opposite} winners
(Fig.~\ref{fig:vaidya-timing}): under the negative-rate branch ($m'<0$) the absolute
periapsis descends but the horizon $2m(v)$ retreats faster, so the relative penetration
$r_{\min}/2m(v_{\rm peri})$ \emph{worsens} for later arrivals; under accretion the
periapsis rises yet the horizon grows faster, so late arrivals graze the relative
horizon. The orbit--horizon chase has reversed outcomes measured in $r$ versus in
horizon units.

\begin{figure}[ht]
    \centering
    \includegraphics[width=\textwidth,height=0.68\textheight,keepaspectratio]{Immagini/fig_vaidya_traiettorie.png}
    \caption{$\tau$-brachistochrones Schwarzschild vs Vaidya, and $p_v(r)$:
    transversality $p_v(r_1)=0$, non-autonomous memory $\propto m'$.}
    \label{fig:vaidya-traj}
\end{figure}


\subsection{Plunge-radius law and the (absent) evaporative inversion}
\label{sec:vaidya-plunge}
The turning-point law has the static Schwarzschild form evaluated at the mass
at periapsis, $f_*=1-2m(v_{\rm peri})/r_{\min}$:
\begin{align}
    \tau:&\quad (\Ehat^2-f_*)J^2 = f_* r_{\min}^2 ,\\
    t:&\quad (\Ehat^2-f_*)J^2 = \Ehat^2 r_{\min}^2/f_* ,
    \label{eq:vaidya-turn}
\end{align}
plus a genuine $O(m')$ correction from the endpoint-sensitive costate memory $p_v$. The
static ratio $J^2_t/J^2_\tau=\Ehat^2/f^2>1$ gives $r_{\min}^t<r_{\min}^\tau$
(``$t$ deeper''), a \emph{branch-intrinsic} ordering set by $\Ehat/f>1$ alone---it
does not depend on the sign of $m'$ nor on the ingoing/outgoing character of the
flux. We stress that \emph{physical} evaporation is the \emph{outgoing} Vaidya metric
$\dd s^2=-f\dd u^2-2\dd u\,\dd r$ (retarded time $u$), not the ingoing metric with
$m'<0$. The depth ordering $r_{\min}^t<r_{\min}^\tau$ is \emph{branch-intrinsic}: it
follows from $\Ehat/f>1$ alone and is therefore independent of the sign of $m'$ and of
the ingoing/outgoing character of the flux, so to this (static, frozen-clock) order it
is unchanged under evaporation. The apparent linear-model ``inversion'' from flipping
$m'<0$ in the \emph{ingoing} model is an unphysical extrapolation to $m\le0$. We do
\emph{not} here solve the full outgoing optimal-control problem with its own
Hamiltonians, costates and transversality: the accretion--evaporation asymmetry
[Eq.~\eqref{eq:vaidya-asymmetry}] is a clock decomposition ($v\!\leftrightarrow\!u$,
the tortoise $r_*$ changing sign) of the frozen first-order correction, not a complete
outgoing boundary-value solution. Accordingly we state the absence of an evaporative
inversion as a property of the branch-intrinsic static ordering together with this
frozen clock decomposition, and leave the full outgoing BVP---and a genuine outgoing
trajectory to accompany Fig.~\ref{fig:vaidya-noinv}---to future work.

\begin{figure}[ht]
    \centering
    \includegraphics[width=\textwidth,height=0.68\textheight,keepaspectratio]{Immagini/fig_vaidya_plunge_t_tau.png}
    \caption{Plunge radii $r_{\min}^{t,\tau}$ and gap vs $\mu=m'$; $t$ deeper,
    accretion widens the gap.}
    \label{fig:vaidya-plunge}
\end{figure}

\begin{figure}[ht]
    \centering
    \includegraphics[width=\textwidth,height=0.68\textheight,keepaspectratio]{Immagini/fig_vaidya_no_inversione_evaporazione.png}
    \caption{This figure establishes \emph{only} the absence of inversion within
    the ingoing, positive-mass model and the frozen-clock branch decomposition; it
    does \emph{not} exhibit a solution of the outgoing-Vaidya boundary-value
    problem. The linear-model sign change lies entirely in the unphysical $m<0$
    region; with $m=e^{-\lambda v}>0$ the branch gap shrinks but does not invert.
    The static/frozen branch ordering does not invert under the outgoing-clock
    replacement; the full outgoing-Vaidya optimal-control boundary-value problem
    remains open.}
    \label{fig:vaidya-noinv}
\end{figure}

\subsection{Closed-form adiabatic correction and the horizon dilogarithm}
\label{sec:vaidya-adiabatic}
When $m(v)$ varies slowly the same frozen-coefficient WKB scheme as in
Sec.~\ref{sec:adiabatic} applies, with the mass function $m(v)$ playing the role of
the conformal factor and $\dot m=\dd m/\dd v$ the small parameter. The frozen family
is the \emph{Schwarzschild} brachistochrone---the $a\to0$ limit of the
Thakurta--Kerr forms---on the sextic
$S=r(r-2m)DE\,[r^2(r-2m)-J^2 DE]$, $DE=(E^2-1)r+2m$. Notably this sextic keeps six
distinct roots: the constrained brachistochrone is genus \emph{two} even in
Schwarzschild (only the geodesics are elliptic; the separatrix $|J|=\Jc$ is the
genus drop). The mass derivative reduces exactly, the horizon factor cancelling,
\begin{equation}
    \partial_m F=\frac{N_m}{S^{3/2}},\quad N_m=S\,\partial_m K-\tfrac12 K\,\partial_m S,
    \qquad
    \partial_m\varphi_0=\frac{\mathcal{A}^m}{\sqrt S}+\sum_{k=0}^4 c_k^m U_k,
    \label{eq:dm-reduction}
\end{equation}
with $U_k=\int r^k/\sqrt S\,\dd r$ and \emph{algebraic} constants $c_k^m$ (verified
to $10^{-15}$). The clock is now the \emph{advanced time}: from $v=t+r_*$,
\begin{equation}
    \frac{\dd v}{\dd r}=\frac{E}{f}\frac{\dd\tau}{\dd r}+\frac1f
    =\frac{E\,r^3}{\sqrt S}+\frac{r}{r-2m}
    \;\Rightarrow\;
    v=E\,U_3+r+2m\ln(r-2m),
    \label{eq:vaidya-clock}
\end{equation}
the tortoise term injecting a \emph{horizon logarithm}. This is the advanced-time
(native) branch; the \emph{proper-time} branch instead uses the second-kind clock
$\dd\tau/\dd r=r^2(r-2m)/\sqrt S$, i.e.\ $\eta=U_3-2m\,U_2$, $b=(0,0,-2m,1,0)$ (no
horizon log), and closes by the very same
assembly~\eqref{eq:psi-split}---$\delta\varphi_\tau/\dot m=\tfrac12\sum_{k<j}Q_{kj}
W_{kj}-\eta(\mathcal{A}^m/\sqrt S+\tfrac12\sum c_k^m U_k)+\sum_j b_j\mathcal{I}_j$ with
$Q_{kj}=c_k^m b_j-c_j^m b_k$, all coefficients algebraic (verified to
$10^{-14}$). The two canonical clocks of the ingoing problem are therefore the
\emph{advanced null time} $v$ (adapted to the Kodama selector $K=\partial_v$, the
native evolution parameter) and the \emph{proper time} $\tau$; the coordinate time
$t$ that enters the static depth comparison~\eqref{eq:vaidya-turn} is the
\emph{frozen} Schwarzschild-like time, a comparison device valid at frozen order and
\emph{not} a canonical non-autonomous clock---once $m=m(v)$ the tortoise
$r_*=r+2m\ln(r-2m)$ acquires $\dd m$ terms and no Killing time survives. The retarded
time $u=t-r_*$ of the outgoing metric is a \emph{separate} outgoing-Vaidya problem
(its own Hamiltonians, costates, and transversality), not a third clock of the ingoing
family. The advanced-time first-order correction
$\delta\varphi=\dot m\!\int\!\partial_m F\,v\,\dd r$ then integrates by parts into a
fully explicit closed form (verified to $9\times10^{-14}$, coefficients given),
\begin{equation}
    \frac{\delta\varphi}{\dot m}
    =A_m\,v-E\!\int\! A_m\,\frac{r^3}{\sqrt S}\,\dd r
    -\int\! A_m\,\frac{r}{r-2m}\,\dd r,\qquad A_m\equiv\partial_m\varphi_0,
    \label{eq:vaidya-deltaphi}
\end{equation}
This is the \emph{on-shell} component of the Vaidya first-order correction (the
mass-drift of the frozen family). As in Thakurta--Kerr, the interior costate does not
vanish---$p_v=-H_v=O(\dot m)$ with $p_v(r_1)=0$ only at the free arrival---so the
\emph{complete} first-order term adds the off-shell piece given by the Vaidya analogue
of Eq.~\eqref{eq:adiab-exact}, with the single-parameter operator $\Theta=m\partial_m$
and $S=\int_0^\lambda\Theta H_v\,\dd\lambda'$ terminally anchored; the extra term is
$-\!\int(\partial_{p_r}G/H_{p_r})\,S\,\dd r$, a length-two iterated integral of the
same horizon-dilogarithm class. Applying the identical canonical perturbation theory
to $H_v$ closes the Vaidya first order to $O(\dot m^2)$ (residual slope $\approx2$;
\texttt{VaidyaMetric/vaidya\_first\_order\_offshell.py}). We display the on-shell
component in closed form below;
its two integrals expand into named building blocks with algebraic coefficients.
Writing $A_m=\mathcal{A}^m/\sqrt S-\alpha_0+\sum_k c_k^m U_k$ with the constant
$\alpha_0=\mathcal{A}^m(r_0)/\sqrt{S(r_0)}$, integration by parts in each factor gives
the \emph{fully assembled} closed form (verified to $2\times10^{-13}$)
\begin{equation}
\begin{aligned}
    \frac{\delta\varphi}{\dot m}
    ={}&A_m\,v
    -E\Big[\mathcal{I}_{\rm poly}-\alpha_0 U_3
        +\tfrac12\sum_{k=0}^4 c_k^m\big(U_kU_3-W_{3k}\big)\Big]\\
    &{}-\Big[\mathcal{L}_{2m}-\alpha_0\,(v-E\,U_3)
        +\sum_{k=0}^4 c_k^m\big(r\,U_k-U_{k+1}+2m\,\mathcal{D}_k\big)\Big],
\end{aligned}
    \label{eq:vaidya-full}
\end{equation}
with the named building blocks: the abelian integrals $U_k=\int r^k/\sqrt S\,\dd r$;
the elementary rational-log primitive
$\mathcal{I}_{\rm poly}=\int\mathcal{A}^m r^3/S\,\dd r$; the genus-two
polylogarithms $W_{3k}=\int(U_3\dd U_k-U_k\dd U_3)$ \emph{at infinity} (as in Kerr);
the \emph{second}-kind horizon integral
$\mathcal{L}_{2m}=\int\mathcal{A}^m r/[(r-2m)\sqrt S]\,\dd r$---since
$\mathcal{A}^m(2m)\ne0$ its integrand behaves as $(r-2m)^{-3/2}$, a double pole with
\emph{no} residue at the horizon Weierstrass point, closed by the Kleinian
$\zeta$/quasi-periods~\cite{BuchstaberEnolskii1997,BakerAbelian1897}, not by Fay; the
genuinely \emph{third}-kind object is the horizon letter $\dd r/(r-2m)$ (the tortoise
log, a simple pole of residue two at that Weierstrass point), which closes as a Fay
$\theta$-ratio~\cite{Fay1973} exactly as $L$ in Eq.~\eqref{eq:thirdkind-theta} but
centred at $r=2m$ instead of infinity; and
the \emph{horizon dilogarithms} $\mathcal{D}_k=\int U_k/(r-2m)\,\dd r$. Since $r=2m$
is a branch point of $S$, in the local uniformiser $t=\sqrt{r-2m}$ one has
$\dd r/(r-2m)=2\,\dd t/t$: a \emph{simple} pole of residue two, i.e.\ a genuine
third-kind letter $\dd\log$ anchored \emph{at the horizon}. On the compact
genus-two curve this pole cannot stand alone: by the residue theorem the
differential $\dd r/(r-2m)$ carries residue $+2$ at the horizon Weierstrass point
and compensating residues $-1$ at each of the two points at infinity $\infty_\pm$ of
the sextic ($\dd r/(r-2m)\sim-\dd s/s$ in $s=1/r$ on each sheet), so the residues sum
to zero. ``$\dd\log$ at the horizon'' names the horizon part of this divisor; the
Fay/prime-form ($\theta$-ratio) representation of the third-kind differential carries
the full divisor $2\,(r{=}2m)-\infty_+-\infty_-$. Its elementary primitive
is $\int\dd r/(r-2m)=\ln(r-2m)$, so one integration by parts closes $\mathcal{D}_k$
onto a \emph{single} weight-two letter (verified to $10^{-12}$). Throughout, displayed
logarithms are to be read with a fixed reference length, $\ln(r-2m)\equiv\ln[(r-2m)/L]$
(likewise $\ln[(r-2M)/L]$); the scale $L$ is dimensionally required but shifts only the
additive constants $\alpha_0$ and cancels in every assembled closed form, so we suppress
it in the displayed primitives.
\begin{equation}
    \mathcal{D}_k=U_k\,\ln(r-2m)-\mathcal{G}_k,\qquad
    \mathcal{G}_k=\int\ln(r-2m)\,\frac{r^k}{\sqrt S}\,\dd r,
    \label{eq:horizon-dilog}
\end{equation}
where $\mathcal{G}_k$ is the horizon-anchored genus-two dilogarithm---the exact
counterpart of the Kerr weight-two piece $W_{jk}$, in the same
hyperelliptic-dilogarithm class (conjecturally irreducible to weight one, the
theta-divisor obstruction $\Theta=W_{g-1}$ preventing the genus-one factorization),
but with its $\dd\log$ letter sourced by the
advanced-time divergence at the horizon $r=2m$ rather than by a third-kind dipole at
infinity. On the separatrix $|J|=\Jc$ the sextic develops a double root at $r=r_d$
and the genus drops to one: $\sqrt S=(r-r_d)\sqrt{Q_4}$ with $Q_4$ quartic, so the
$U_k$ reduce to elliptic integrals (verified to machine precision against the direct
genus-two quadrature) and $\mathcal{G}_k$ collapses to the tabulated \emph{elliptic
(Bloch--Wigner) dilogarithm} on $E:\,w^2=Q_4(r)$---the maximally explicit, directly
reusable form. For the representative $m=1,E=1.4$ one finds $J_{\rm deg}=7.0266$,
$r_d=-3.3637$, branch points $\{-2.083,0,2m,8.727\}$, elliptic modulus $k^2=0.6067$
and rectangular nome $\tau=0.90597\,i$ (agreeing to fifteen digits between the
Legendre form $\tau=iK'/K$ and the Sage period lattice); the horizon $r=2m$, being a
branch point of $Q_4$, sits at a \emph{two-torsion} point (half the imaginary
period), so it carries the $\dd\log$ letter of the elliptic dilogarithm. At these
parameters $r_d=-3.36<0$: this locus is the \emph{algebraic} genus-degeneration
$J_{\rm deg}$ (a double root of the sextic where the closed form drops to genus one),
not a physical external separatrix---the Schwarzschild $\tau$-brachistochrone at
$E=1.4$ has \emph{no} double root at $r>2m$, hence no external capture$/$scatter
unstable circular orbit, and $J_{\rm deg}$ must not be conflated with the dynamical
penetration threshold $J_c(v_0)$. The genus-one Weierstrass reduction that follows is
purely algebraic and valid regardless of the sign of $r_d$
(\texttt{separatrix\_classification.py}). Concretely,
with the Weierstrass $\sigma,\zeta$ of $E$ (built from $\tau$ with base at the
Weierstrass point, so the involution is $z\to-z$) the weight-one integrals close in
closed form as third-kind elliptic integrals,
$U_0=\rho\,[\ln\sigma(z-z_d)-\ln\sigma(z+z_d)]+Cz$ with the residue
$\rho=1/\sqrt{Q_4(r_d)}$ and $C$ both \emph{algebraic} (derived, not fitted; verified
to $10^{-14}$ against direct quadrature), and one further integration against the
horizon third-kind letter turns $\mathcal{G}_k$ into the elliptic
dilogarithm~\cite{Bloch2000,BeilinsonLevin1994}
$\int\ln\sigma\,\zeta\,\dd z$, whose single-valued projection is the tabulated
Zagier series~\cite{Zagier1990}
$D^E(z)=\sum_{n\ge0}D(q^n\zeta)-\sum_{n\ge1}D(q^n/\zeta)$
($q=e^{2\pi i\tau}$, $D$ the Bloch--Wigner function). As in Kerr, every coefficient is
algebraic (rational in $m,E,J$); the transcendence lives entirely in
$W_{jk},\mathcal{G}_k$, so the correction is reusable by direct evaluation once $S$
and the $c_k^m$ are fixed.

\paragraph{Fully explicit separatrix correction and its minimal dilogarithm basis.}
The separatrix reduction can be pushed to a fully explicit form. Writing
$\partial_m F\,\dd r=R(z)\,\dd z$ with $R$ the elliptic function
$R=N_m(r(z))/[(r(z)-r_d)^3\,Q_4(r(z))]$ on $E$ (where $r(z)$ is itself explicit,
$r(z)=c_r-a_4^{-1/2}[\zeta(z-z_\infty)-\zeta(z+z_\infty)]$), its Mittag--Leffler
decomposition on the torus---principal parts extracted by contour residues, in
closed algebraic form---integrates term by term to
\begin{equation}
    G(z)=\int\partial_m F=C_0\,z+\sum_a\Big[b_1^a\ln\sigma(z-a)-b_2^a\zeta(z-a)
        -\tfrac12 b_3^a\wp(z-a)\Big],
    \label{eq:vaidya-sep-G}
\end{equation}
verified $\dd G/\dd z=R$ to $10^{-13}$, the poles $a\in\{\pm z_d,\,0,\,i\omega_{\rm im}\}$
carrying respectively the third-, second- and higher-kind letters, and each residue
$b_n^a$ given in closed form (the triple pole at $z_d$ from a Hermite--Ostrogradsky
Laurent expansion, the double poles at the two-torsion turning points from
$b_2=4N_m(e_i)/[(e_i-r_d)^3 Q_4'(e_i)^2]$; the additive constant closes at the
regular point $z_\infty$, where $R$ \emph{vanishes} as $1/r$, giving
$C_0=-\sum_a[b_1^a\zeta(z_\infty-a)+b_2^a\wp(z_\infty-a)-\tfrac12 b_3^a
\wp'(z_\infty-a)]$, and the analogous constant of the clock differential
$\eta'$ from its value at the regular turning image $z=0$). The weight-two blocks
$W_{kj}=2\int U_k\,\dd U_j-U_kU_j$ close on the Brown--Levin elliptic
polylogarithm, whose transcendental core is the length-two Kronecker--Eisenstein
iterated integral~\cite{BrownLevin2011}
$\tilde\Gamma(1,1;\hat a,\hat b;y)=\int_0^y g^{(1)}(t-\hat a)\big[\int_0^t
g^{(1)}(s-\hat b)\,\dd s\big]\dd t$ with the \emph{tabulatable} $q$-series
$g^{(1)}(x)=\pi\cot\pi x+4\pi\sum_{n\ge1}\frac{q^{2n}}{1-q^{2n}}\sin 2\pi n x$
(agreeing with the direct nested quadrature, and with the GiNaC elliptic-polylog
evaluation, to $10^{-18}$). This building-block form---$G$ from
Eq.~\eqref{eq:vaidya-sep-G} plus the $W_{kj}=\tilde\Gamma(1,1)$---is the
\emph{reusable} closed form: \emph{every} coefficient in it (the residues $b_n^a$,
the mass constants $c_k^m$, $\Jc$ itself) is given in closed form as an algebraic
function of $m,E$, so a reader reconstructs $\delta\varphi_\tau|_{\rm sep}$ for any
$(m,E)$ by evaluating those formulas on the corresponding curve---checked to
$10^{-9}$ against the direct integral and, term by term for the closed-form
coefficients, across $E\in[1.3,1.6]$ (each a distinct curve).

Beyond this reusable form there is a purely \emph{structural} statement about the
transcendental weight: the parity involution $z\to-z$ groups the raw terms into
$24$ blocks, the shuffle identity~\cite{Chen1977} recasts these as $21$
antisymmetric elliptic dilogarithms, and the Weierstrass-addition (Fay) relations
collapse them further, so that the entire weight-two content of
$\delta\varphi_\tau|_{\rm sep}$ appears to span a space of dimension five: a
high-precision numerical reduction gives a stable rank five (the discarded singular
values fall to $10^{-79}$ in arbitrary precision). We state this as strong
\emph{numerical evidence}, not a proof---an exact bound on the dimension would
require a motivic/Hodge basis-count or an explicit basis with a proven
independence, which we do not carry out. The coefficients $\beta_b$ of any chosen
five-element dilogarithm basis are \emph{curve-specific numbers}---obtained per curve by a
mechanical linear reduction (here, at $m=1,E=1.4$, to eighty digits), not a
universal formula---and a closed symbolic form for them, which would require the
full second-kind reduction of $E$, is left open. What makes the algebraic reduction
possible is that every $e_i(z)$ entering the assembly is a \emph{rational} function
of $\wp(z)$---e.g.\ the third-kind combination
$\zeta(z-a)-\zeta(z+a)=-2\zeta(a)+\wp'(a)/[\wp(z)-\wp(a)]$.

\paragraph{Accretion--evaporation asymmetry.} At first order $\delta\varphi$ is
linear in $\dot m$, so na\"ively $\dot m\to-\dot m$ merely flips its sign. The
physical asymmetry is subtler: accretion ($\dot m>0$) is \emph{ingoing} Vaidya,
clocked by the \emph{advanced} time $v=t+r_*$, whereas evaporation ($\dot m<0$) is
\emph{outgoing} Vaidya, clocked by the \emph{retarded} time $u=t-r_*$
($\dd s^2=-f\dd u^2-2\dd u\,\dd r$). The frozen reduction $\partial_m F$ is identical
(it depends only on $S$); only the clock changes, and $v\leftrightarrow u$ flips the
sign of the tortoise $r_*$. Splitting the clock into its bulk and tortoise parts,
$v=E\,U_3+r_*$ with $r_*=(r-r_0)+2m\ln\!\frac{r-2m}{r_0-2m}$,
\begin{equation}
    \frac{\delta\varphi}{\dot m}\Big|_{\rm accr}=\mathcal{A}_\infty+\mathcal{B}_{\rm hor},
    \quad
    \frac{\delta\varphi}{\dot m}\Big|_{\rm evap}=\mathcal{A}_\infty-\mathcal{B}_{\rm hor},
    \quad
    \mathcal{A}_\infty=\!\int\!\partial_m F\,E\,U_3\,\dd r,\;\;
    \mathcal{B}_{\rm hor}=\!\int\!\partial_m F\,r_*\,\dd r,
    \label{eq:vaidya-asymmetry}
\end{equation}
so that, since $\dot m<0$ in evaporation flips the sign of
$\delta\varphi_{\rm evap}=\dot m(\mathcal{A}_\infty-\mathcal{B}_{\rm hor})$, at equal
rate magnitude the physical \emph{difference} is
$\delta\varphi_{\rm accr}-\delta\varphi_{\rm evap}=2|\dot m|\,\mathcal{A}_\infty$
(carried by the infinity bulk), while the \emph{deviation from exact antisymmetry} is
$\delta\varphi_{\rm accr}+\delta\varphi_{\rm evap}=2|\dot m|\,\mathcal{B}_{\rm hor}$,
carried \emph{entirely by the tortoise/horizon}: the horizon term $\mathcal{B}_{\rm hor}$
is exactly what breaks the na\"ive $\dot m\to-\dot m$ antisymmetry. Within
$\mathcal{B}_{\rm hor}$ the transcendental, direction-sensitive content is precisely
the horizon term $2m\!\int\!\partial_m F\ln(r-2m)\,\dd r$, which by parts is the
horizon dilogarithm $\mathcal{D}_k$ of Eq.~\eqref{eq:horizon-dilog}. Thus the
accretion/evaporation asymmetry is the physical face of the horizon weight-two piece,
just as the ergosphere-penetration window (finite for one sign of $J$, of measure zero
for the other) is the phenomenological face of the same horizon geometry. For
$m=1,E=1.4,J=2.5$: $\mathcal{A}_\infty=8.7276$, $\mathcal{B}_{\rm hor}=6.1127$ (of which
$1.5301$ is the horizon-log dilogarithm sector), giving
$\delta\varphi/\dot m|_{\rm accr}=14.840$ against
$\delta\varphi/\dot m|_{\rm evap}=2.615$ (both verified to $10^{-15}$); the physical
difference is then $2\mathcal{A}_\infty=17.455$ and the deviation from antisymmetry
$2\mathcal{B}_{\rm hor}=12.225$ (each $\times|\dot m|$).


% ======================================================================
\section{Thakurta--Kerr: conformal rotating compact object}
\label{sec:thakurta}
% ----------------------------------------------------------------------
Thakurta--Kerr is conformal Kerr, $g=A(\eta)^2 g_{\rm Kerr}$
\cite{Thakurta1981,Kerr1963}, one of a family of cosmological geometries
\cite{McVittie1933,SultanaDyer2005,Kaloper2010,NolanMcVittie1998,MelloMaciel2017,FaraoniBook}.
For constant $A$
the transfer theorem maps the problem to Kerr with $E_{\rm eff}=\Ehat/A$,
$J_{\rm eff}=J/A$. Throughout we use $\rp$ (where $\Delta=0$) and $\re=2M$ as
labels for the \emph{seed} Kerr null surface and conformal-stationary-limit surface:
a conformal factor preserves the null \emph{character} of $\Delta=0$, so $\rp$ is a
conformal Killing horizon candidate, but for time-dependent $A(\eta)$ this does
\emph{not} by itself fix the apparent/trapping or event horizon (which require the
null expansions and the global causal structure, not computed here). With that
caveat, $\rp$ is rigid, $\re=2M$ is crossed regularly, and the freezing surface
$A^2 f=\Ehat^2$ breathes.

\paragraph{Indicatrix and Hamiltonians.} Equatorially, write
$\Delta=r^2-2Mr+a^2$, $\bar v^2=1-A^2f/\Ehat^2$, and
$\bar P=P+A^2(2Ma/r)^2/\Ehat^2$ with $P=r^2+a^2+2Ma^2/r$. The rail
$\Ehat=-u_\eta$ makes the indicatrix an ellipse with an \emph{azimuthal} wind
(frame dragging),
\begin{equation}
    \varphi'_0=\frac{2Ma}{r}\frac{\bar v^2}{\bar P},\qquad
    R^2=f\bar v^2+\bar P\varphi_0'^2=\frac{\Delta\,\bar v^2}{\bar P},
    \label{eq:tk-ind}
\end{equation}
the last equality using the Kerr identity $fP+(2Ma/r)^2=\Delta$. The
conformal-time, coordinate-time and proper-time Hamiltonians are
\begin{align}
    H_\eta&=p_\varphi\varphi'_0
        +R\sqrt{(\Delta/r^2)p_r^2+p_\varphi^2/\bar P}-1,
    \label{eq:Heta}\\
    H_t&=p_\varphi\varphi'_0
        +R\sqrt{(\Delta/r^2)p_r^2+p_\varphi^2/\bar P}-A,
    \label{eq:Ht-tk}\\
    H_\tau&=\tilde p_\varphi\varphi'_0
        +R\sqrt{(\Delta/r^2)p_r^2+\tilde p_\varphi^2/\bar P}-\frac{A^2 f}{\Ehat},
    \label{eq:Htau-tk}
\end{align}
with the conformal gravitomagnetic shift
$\tilde p_\varphi=p_\varphi-2MaA^2/(r\Ehat)$. The arrival-time branch $H_t$ differs
from $H_\eta$ only in the unit cost ($-1\to-A$) and reduces to it for $A=1$; it is
the branch compared with $\tau$ in the plunge inversion
(Sec.~\ref{sec:inversion}), where its turning potential
$H_{t0}=H_t|_{p_r=0}$ enters. All three are regular through the
ergosphere ($f=0\Rightarrow R^2=\bar v^2\Delta/\bar P>0$) and degenerate only at
the horizon $\Delta=0$, which is therefore conformally invariant---independent of
$A$ and $\Ehat$. The three branches share this ellipse: $t$ follows the same
curves as $\eta$ through $t=\int A\,\dd\eta$, while $\tau$ is Riemannian-pure, the
wind cancelling the gravitomagnetic term, so it survives the expansion
(Fig.~\ref{fig:th-branches}).

\paragraph{The conformal-time ($\eta$) brachistochrone.} The relation
$\dd t=A(\eta)\,\dd\eta$ makes $\eta$ the canonical conformal clock and shows that the
arrival branch carries no new content beyond $t$: the two are the \emph{same} spatial
family. Frozen, the arrival-time Finsler functions obey $F_t=A\,F_\eta$ exactly, so at
fixed $A$ the brachistochrone curves coincide (a constant Finsler rescaling leaves
geodesics invariant), the conserved charge merely relabelled by the conformal weight
$J_t=A\,J_\eta$: the $\eta$-orbit at $J_\eta$ and the $t$-orbit at $A\,J_\eta$ trace the
same $\varphi(r)$ to $7\times10^{-9}$
(\texttt{ThakurtaMetric/eta\_brachistochrone\_check.py}), whereas the naive same-$J$
comparison does \emph{not}. Under a drift the equivalence survives by monotonicity: for
any fixed-endpoint path the arrival time $t_f=\int_0^{\eta_f}A\,\dd\eta$ is a strictly
increasing function of $\eta_f$ ($A>0$), so $\arg\min\eta_f=\arg\min t_f$ for
\emph{every} $A(\eta)$. Hence the frozen, adiabatic and separatrix $\eta$-solutions are
exactly the $t$-solutions of Secs.~\ref{sec:closed}--\ref{sec:thakurta}; there are
\emph{two} distinct optimal families, not three---the arrival branch $t\equiv\eta$ and
the proper-time branch $\tau$. The $\eta$ gauge is nonetheless the natural one: $H_\eta$
carries unit cost, the extended Hamiltonian $\mathcal H=p_\eta+H_\eta=0$ is the cleanest
non-autonomous form, the drift enters as the single scalar $\varepsilon=A'/A$, and the
off-shell source takes the universal form $S_{\mathcal D}=\eta+\int E_{\rm eff}
H_{E_{\rm eff}}\,\dd\eta$ [Eq.~\eqref{eq:SD-conformal}].

\paragraph{Non-autonomous optical reduction.} Solving the indicatrix for
$\dd\eta$ at fixed spatial displacement casts the $\eta$-branch as a closed-form
Randers metric, non-autonomous through the conformal factor,
\begin{align}
    \dd\eta &= n(\eta,r)\,\alpha_K+\beta_K,\qquad
    n=\frac{\Ehat}{\sqrt{\Ehat^2-A^2 f}},
    \label{eq:randers}\\
    \alpha_K^2 &= \frac{r^2\,\dd r^2}{f\Delta}+\frac{\Delta\,\dd\varphi^2}{f^2},
    \qquad
    \beta_K = -\frac{2Ma}{rf}\,\dd\varphi.
    \label{eq:randers-data}
\end{align}
The Riemannian data $\alpha_K$ and the one-form $\beta_K$ are \emph{exactly} those
of the \emph{null} Fermat (Randers) metric of Kerr \cite{GibbonsWerner2008}: all
of the mass, rail, and expansion physics is carried by the scalar index $n$,
which multiplies only the Riemannian part, while the frame-dragging one-form
$\beta_K$ is rigid---conformally invariant and independent of $\Ehat$. The
reduction is thus a Zermelo problem with a time-dependent landscape in which only
the metric term breathes, not the wind. It degenerates at the ergosphere ($f=0$,
where $1/F$ is intrinsic, as for stationary Kerr) even though the
Hamiltonian~\eqref{eq:Heta} remains regular there; the limits $\Ehat\to\infty$
(null Kerr optics), $A=1$ (stationary Perlick on Kerr), $a=0$
(Thakurta--Schwarzschild) and $M=a=0$ ($n=1/v$, FLRW) are all recovered.



\begin{figure}[ht]
    \centering
    \includegraphics[width=\textwidth,height=0.68\textheight,keepaspectratio]{Immagini/fig_thakurta_kerr_rami.png}
    \caption{The three branches: $\tau$ bounces (mirror pair, no wind), $\eta$
    spirals (gravitomagnetic wind); breathing conformal scalars $n,k_\tau$.}
    \label{fig:th-branches}
\end{figure}

\paragraph{Branch optical metrics across the ladder.} The same Zermelo inversion
casts \emph{every} branch of \emph{every} metric in the paper as a Randers form
$\dd(\text{cost})=n\,\alpha+\beta$ (Table~\ref{tab:optical}), with the Riemannian
Fermat data $\alpha$ and wind one-form $\beta$ shared by the branches of a given
metric and only the scalar index $n$ set by the clock. Three structural facts
hold throughout and are verified to machine precision: (i) the proper-time branch
is always \emph{Riemannian-pure}, $\beta_\tau=0$---the wind (ingoing for Vaidya,
gravitomagnetic for Thakurta--Kerr) cancels in $\tau$; (ii) the arrival/native
branches carry the wind but share $\alpha$, and the arrival--proper split is the
universal ratio $n_t/n_\tau=E/f$ ($E=\Ehat$ for Vaidya, $E_{\rm eff}=\Ehat/A$ for
Thakurta--Kerr); (iii) non-stationarity makes only $\alpha,n$ \emph{breathe}
(through $m(v)$ or $A(\eta)$), while $\beta$ is rigid. FLRW is the degenerate top
of the ladder: $n$ is position-independent so all three branches share the same
(straight) geodesic. The $a=0$ limit of the Thakurta $\tau$-row reproduces the
static Vaidya $\tau$-row, and $M=a=0$ (or $m=0$) gives $n=1/v$, FLRW.

\begin{table}[ht]
    \centering
    \small
    \caption{Branch optical (Randers) metrics $\dd(\text{cost})=n\,\alpha+\beta$.
    Here $f=1-2m(v)/r$ (Vaidya) or $1-2M/r$ (Thakurta, with spin $a$,
    $\Delta=r^2-2Mr+a^2$), $w=\Ehat^2-f$, $w_A=\Ehat^2-A^2f$. The wind $\beta$
    acts on the native/arrival branch ($v$ for Vaidya, $\eta,t$ for Thakurta);
    $\tau$ is always Riemannian-pure ($\beta_\tau=0$). Native indices:
    $n_v=n_t^{\rm V}$ (Vaidya), $n_\eta=\Ehat/\sqrt{w_A}$ (Thakurta).}
    \label{tab:optical}
    \begin{tabular}{lcccc}
    \hline
    metric & $\alpha^2$ & $\beta$ & $n_\tau$ & $n_t$ \\
    \hline
    FLRW & $\dd r^2+r^2\dd\varphi^2$ & $0$
      & \multicolumn{2}{c}{$\Ehat/\sqrt{\Ehat^2-a^2}$ \ (degenerate)} \\[2pt]
    Vaidya & $\dd r^2+f\,r^2\dd\varphi^2$ & $\dd r/f$
      & $1/\sqrt{w}$ & $\Ehat/(f\sqrt{w})$ \\[2pt]
    Thakurta--Kerr
      & $\dfrac{r^2\dd r^2}{f\Delta}+\dfrac{\Delta\,\dd\varphi^2}{f^2}$
      & $-\dfrac{2Ma}{rf}\dd\varphi$
      & $\dfrac{A^2 f}{\sqrt{w_A}}$ & $\dfrac{A\Ehat}{\sqrt{w_A}}$ \\[3pt]
    \hline
    \end{tabular}
\end{table}

\subsection{Quasi-constants and drift}
\label{sec:quasiconst}
For constant $A$ the transfer theorem is exact at the level of conserved
quantities: $u_{\rm TK}=u_{\rm K}/A$ maps the $\Ehat$-rail onto the Kerr rail with
$E_{\rm eff}=\Ehat/A$, so every Kerr quasi-constant carries over by substitution,
\begin{equation}
    K^{\rm TK}=K^{\rm Kerr}(E_{\rm eff}),\qquad
    K_\tau=Q_{\rm std}+a^2 f_2\,p_\theta^2+O(a^4),
    \label{eq:Ktransfer}
\end{equation}
with $Q_{\rm std}$ the Carter constant and $f_2$ the next-to-leading polar
coefficient \cite{Carter1968,RosignoliKerr}. The polar libration that carries the
quasi-constant lives off the equatorial plane, governed by the full Hamiltonian
on a generic slice $\Sigma=r^2+a^2\cos^2\theta$,
\begin{equation}
    H_\tau=\tilde p_\varphi\varphi'_0
      +R\sqrt{\frac{\Delta}{\Sigma}p_r^2+\frac{p_\theta^2}{\Sigma}
      +\frac{\tilde p_\varphi^2}{\bar G}}-\frac{A^2 f_\Sigma}{\Ehat},
    \label{eq:H3d}
\end{equation}
with $f_\Sigma=1-2Mr/\Sigma$, $b=2Mar\sin^2\theta/\Sigma$,
$\bar G=G+A^2 b^2/\Ehat^2$, $\tilde p_\varphi=J-A^2 b/\Ehat$,
$\varphi'_0=b\bar v_\Sigma^2/\bar G$, $R^2=\bar v_\Sigma^2\Delta\sin^2\theta/\bar G$,
and the block identity $f_\Sigma G+b^2=\Delta\sin^2\theta$; it reduces to
Eqs.~\eqref{eq:Heta}--\eqref{eq:Htau-tk} on the equator. For $A=A(\eta)$ the
substitution~\eqref{eq:Ktransfer} is only the zeroth order: the transport equation
$\dd K/\dd\lambda=\{K,H\}+\partial_\eta K$ acquires an explicit source at order
$a^2\varepsilon$, $\varepsilon=A'/A$,
\begin{equation}
    \frac{\dd K}{\dd\eta}\bigg|_{\rm expl}=\varepsilon\,a^2\,S_1,\qquad
    S_1=2E_{\rm eff}^2\!\left[\cos^2\theta
        +\frac{M\cos2\theta\,p_\theta^2}{r\,DE_0^2}\right],
    \label{eq:S1}
\end{equation}
where $DE_0=r^2 w_{\rm eff}$. The counterterm $h$ solving $\{h,H_0\}=-S_1$
inherits the double pole at the \emph{freezing wall}
$r_w=2M/(1-E_{\rm eff}^2)$, so the dynamical quasi-constant is
\begin{equation}
    K^{\rm TK}(\eta)=K^{\rm Kerr}\!\big(E_{\rm eff}(\eta)\big)
        -a^2\!\int\!\varepsilon\,S_1\,\dd\eta,
    \label{eq:Kdyn}
\end{equation}
with closed toroidal averages
$\langle S_1\rangle=E_{\rm eff}^2(1-J^2/L^2)
+2ME_{\rm eff}^2|J|(L-|J|)/(r\,DE_0^2)$, $L^2=p_\theta^2+J^2/\sin^2\theta$. In
Kerr ($E>1$) the wall is absent; in Thakurta--Kerr with $E_{\rm eff}<1$ it enters
from the far field and advances inward as $A$ grows, a \emph{moving kinematic
boundary} (not a proved dynamical attractor): past it the bound $\tau$-libration
ceases to exist, so the dynamical quasi-constants live on scattering arcs
($E_{\rm eff}>1$) or on the $t$-branch.

\subsection{Semi-analytic first-order adiabatic orbit shape}
\label{sec:adiabatic}
The transport equation~\eqref{eq:Kdyn} controls the quasi-constant; the
\emph{orbit shape} $\varphi(r,A)$ itself admits a companion first-order adiabatic
form~\cite{BornFock1928,BenderOrszag1999} that is explicit up to a single higher
transcendental. Writing the frozen integrand as
$F(r;E_{\rm eff},J_{\rm eff})=\dd\varphi/\dd r$ (Eq.~\eqref{eq:t-K} for the
$t$-branch, Eq.~\eqref{eq:tau-sextic} for $\tau$) and linearizing the
\emph{conformal weights} of the Kerr charges---both $E_{\rm eff}=\Ehat/A$ and
$J_{\rm eff}=J/A$ carry weight $A^{-1}$, since for $u^{\rm TK}=u^{\rm Kerr}/A$ one
has $-g^{\rm TK}(u,\partial_\eta)=A\,E_{\rm eff}$ and
$g^{\rm TK}(u,\partial_\varphi)=A\,J_{\rm eff}$---in
$A(\eta)=A_0(1+\varepsilon\eta+\dots)$, $\varepsilon=A'/A$, gives
\begin{equation}
    \varphi(r,A)=\underbrace{\varphi_0(r;E_{\rm eff},J_{\rm eff})}_{\text{frozen, Kleinian}}
    +\varepsilon\big[\mathcal{C}(r)+\psi(r)\big]+O(\varepsilon^2),
    \label{eq:phi-adiab}
\end{equation}
where the first-order coefficient is the elapsed-clock--weighted \emph{scaling}
derivative
\begin{equation}
    \delta\varphi=-\!\int\!\big(\Ehat\,\partial_E F+J\,\partial_J F\big)\,\eta\,\dd r
    =-\!\int\!\big(E\,\partial_E+J\,\partial_J\big)F\,\eta\,\dd r,
    \label{eq:deltaphi-EJ}
\end{equation}
i.e.\ the Euler scaling operator $(E\partial_E+J\partial_J)F$ acting through both
drifting charges (the $J\,\partial_J F$ term is required by axisymmetric momentum
conservation and was absent in an earlier form; Vaidya, whose parameter is the
mass rather than a conformal factor, keeps only its $\partial_m F$ source). Here
$\eta(r)$
is the clock the orbit accumulates along the frozen slice: coordinate time
$\eta=t(r)$ on the $t$-branch, proper time $\eta=\tau(r)=\int\!L_\tau\,\dd r$ with
$L_\tau=\sqrt{Q/w}$, $Q=\Delta\varphi'^2-2a\beta\varphi'+r^2/(r^2+a^2)$, on the
$\tau$-branch (the extra $A^{-2}$ of the proper-time functional affects the
\emph{timing}, not the shape, so Eq.~\eqref{eq:phi-adiab} is common to both
branches).

\paragraph{Closed piece.} The scaling derivative reduces to the same depth-one
class as $\varphi_0$ itself, by the classical reduction of abelian differentials to
first/second kind~\cite{BakerAbelian1897,GriffithsHarris1978}. Both pieces
$\partial_E F$ and $\partial_J F$ share the identical basis: solving
\begin{equation}
    \big(E\partial_E+J\partial_J\big)F=\dd_r\!\left[\frac{\mathcal{A}(r)}{\mathcal{R}}\right]
    +\sum_{k=0}^{4}c_k\,r^k/\mathcal{R}
    \label{eq:dEphi0}
\end{equation}
for a degree-five numerator $\mathcal{A}(r)=\sum_{i=0}^{5}a_i r^i$ and constants
$c_k=\Ehat\,c_k^{E}+J\,c_k^{J}$ (here $\mathcal{R}$ is the branch sextic radical,
$\sqrt{R_6}$ or $\sqrt{S}$) closes in the second kind with \emph{no} third-kind
(horizon-pole) contribution for both branches (verified to machine precision); the
$J\,\partial_J F$ term merely adds its own coefficients $c_k^{J}$ to the existing
$c_k^{E}$, requiring no new building block
(\texttt{ThakurtaMetric/conformal\_source\_Jterm.py}). Hence
$\mathcal{C}(r)=-\tfrac12\big(E\partial_E+J\partial_J\big)\varphi_0\,\eta(r)$ is
a product of depth-one objects and is closed in the Kleinian/abelian class.

\paragraph{Clock differentials.} The clock $\eta$ enters $\mathcal{C}$ through its
differential, which we obtain in closed form on the branch curve. On the
$\tau$-branch the proper-time density reduces to a \emph{rational} second-kind
differential,
\begin{equation}
    \frac{\dd\tau}{\dd r}=\frac{r^2(r-2M)}{\sqrt{S}},
    \label{eq:clock-tau}
\end{equation}
the frame-dragging term cancelling identically (verified to $10^{-16}$). The
$t$-branch coordinate-time density is \emph{also entirely on-curve}: using the
Randers form $\dd t=n_t\alpha+\beta$ and the identity
$(r-2M)Q_2+DE\,(J(r-2M)+2Ma)^2=E^2r^3\Delta$ (so that $n_t\alpha=E^2r^3/(f\sqrt{R_6})$
exactly, $E^2-f=DE/r$), the gravitomagnetic wind collapses and
\begin{equation}
    \frac{\dd t}{\dd r}=\frac{\rho_t(r)}{\sqrt{R_6}},\qquad
    \rho_t=\frac{E^2r^3-(2Ma/r)\,\mathcal K}{f}
          =P_3(r)+\frac{R_\Delta(r)}{\Delta},
    \label{eq:clock-t}
\end{equation}
with $\mathcal K$ the frozen numerator [Eq.~\eqref{eq:t-K}], the ergosphere pole
$r=2M$ \emph{removable}, and third-kind poles only at the horizons ($\Delta=0$). The
second-kind polynomial part is
$P_3=E^2r^3+2E^2Mr^2+4E^2M^2r+8E^2M^3-2(E^2-1)JMa$ (giving the clock vector
$b=(8E^2M^3-2(E^2-1)JMa,\,4E^2M^2,\,2E^2M,\,E^2,\,0)$, the $t$-analogue of the
$\tau$-branch $b=(0,0,-2M,1,0)$), and $R_\Delta$ is the degree-one horizon numerator
(explicit in \texttt{reproduce\_reductions.py}). Thus \emph{neither} branch leaves
the genus-two curve---the $t$-branch merely adds a third-kind horizon letter to the
second-kind clock (verified against the Hamiltonian flow to $10^{-10}$); the apparent
Doran double cover of a first analysis was spurious. Table~\ref{tab:branch-adiab}
collects the branch data: the two solutions share the structure~\eqref{eq:phi-adiab}
and differ \emph{only} in the branch curve and the clock.
\begin{table}[htbp]
    \centering
    \caption{The two branches of the adiabatic solution~\eqref{eq:phi-adiab}. They
    share $\varphi=\varphi_0+\varepsilon[\mathcal{C}+\psi]$;
    only the branch curve $\mathcal{R}^2$, the frozen numerator $K$
    ($F=K/\mathcal{R}$), and the clock $\dd\eta/\dd r$ change. Here
    $\Delta=r^2-2Mr+a^2$, $DE=(E^2-1)r+2M$, all at $E=E_{\rm eff}=\Ehat/A$;
    $\varphi_0,\partial_E\varphi_0$ are Kleinian/abelian on $y^2=\mathcal{R}^2$, and
    $\psi=\tfrac12\Ehat\sum Q_{kj}W_{kj}$ has algebraic coefficients times the
    genus-two polylogarithms $W_{kj}$ [Eq.~\eqref{eq:psi-split}].}
    \label{tab:branch-adiab}
    \begin{tabular}{lll}
        \hline\hline
         & $t$-branch (coord.\ time) & $\tau$-branch (proper time)\\
        \hline
        curve $\mathcal{R}^2$ & $R_6=r\,Q_2\,DE$ & $S=r(r{-}2M)DE[r\Delta{-}J^2DE]$\\
        numerator $K$ & $r\,DE\,(J(r{-}2M){+}2Ma)/\Delta$ & $J\,r(r{-}2M)DE/\Delta$\\
        clock $\dd\eta/\dd r$ & $\dfrac{\rho_t}{\sqrt{R_6}}$ (2nd$+$3rd kind)
            & $\dfrac{r^2(r-2M)}{\sqrt{S}}$ (2nd kind)\\[4pt]
        off-curve cover & none & none\\
        timing prefactor & $1$ & $A^{-2}$\\
        \hline\hline
    \end{tabular}
\end{table}

\paragraph{The correction $\psi$: algebraic coefficients times genus-two
polylogarithms.} By parts and the shuffle identity~\cite{Chen1977} the symmetric
combination $\rho+\tilde\rho=\partial_E\varphi_0\,\eta$
($\rho=\int\partial_E\varphi_0\,\dd\eta$,
$\tilde\rho=\int\eta\,\partial_E F\,\dd r$) is closed, and the antisymmetric
remainder
\begin{equation}
    \psi(r)=\tfrac12\Ehat\,(\rho-\tilde\rho)
    \label{eq:psi-irr}
\end{equation}
is a depth-two iterated integral. Expand $\partial_E\varphi_0$ and $\eta$ in the
cohomology basis $U_k=\int r^k/\mathcal{R}\,\dd r$: the reduction
$(E\partial_E+J\partial_J)F=N/\mathcal{R}^3$ with $N=\Ehat\,N_E+J\,N_J$ a degree-seven
polynomial ($N_E=EJ\,r^4(r-2M)^2\,DE$, and $N_J$ its angular analogue; the horizon
factor $\Delta$ cancels identically in both) gives
$(E\partial_E+J\partial_J)\varphi_0=\mathcal{A}/\mathcal{R}+\sum_{k=0}^{4} C_k U_k$,
$C_k=\Ehat\,c_k^{E}+J\,c_k^{J}$, and on the
$\tau$-branch $\eta=U_3-2M\,U_2$, i.e.\ $b=(0,0,-2M,1,0)$. The \emph{complete}
first-order coefficient $\delta\varphi=\mathcal{C}+\psi=-\!\int(\Ehat\,\partial_E F
+J\,\partial_J F)\,\eta\,\dd r$ then reorganizes, with \emph{no term discarded}, into
\begin{equation}
    \delta\varphi=
    \underbrace{\tfrac12\!\sum_{k<j}Q_{kj}\,W_{kj}}_{\text{weight two}}
    \;-\;\underbrace{\eta\Big(\tfrac{\mathcal{A}}{\mathcal{R}}
    +\tfrac12\!\sum_k C_k U_k\Big)}_{\text{clock}\times\text{abelian}}
    \;+\;\underbrace{\sum_j b_j\,\mathcal{I}_j}_{\text{elementary}},
    \label{eq:psi-split}
\end{equation}
with $Q_{kj}=C_k b_j-C_j b_k$, $W_{kj}=\int(U_k\,\dd U_j-U_j\,\dd U_k)$ the
genus-two polylogarithms, and $\mathcal{I}_j=\int\mathcal{A}\,r^j/\mathcal{R}^2\,
\dd r$ \emph{elementary} (rational integrands, evaluating to logarithms). Every
piece is explicit: the second group is a product of the elapsed clock with the
reduced abelian integrals (weight one), the third an elementary logarithmic
remainder---none is spurious or dropped. The coefficients are \emph{algebraic}
(rational in $M,a,E,J$ through the $C_k=\Ehat c_k^E+J c_k^J$); we
verified~\eqref{eq:psi-split} to $4\times10^{-14}$ for the energy part and to
$1\times10^{-7}$ for the full Euler source with all coefficients \emph{given}, not
fitted. The
holomorphic$\times$holomorphic coefficient vanishes, $Q_{01}=0$ (the clock has no
first-kind part). On the $t$-branch the same reorganization holds with
$\mathcal{R}^2=R_6$ and $K$ of Table~\ref{tab:branch-adiab}: its clock is likewise
on-curve, with the second-kind vector $b$ of Eq.~\eqref{eq:clock-t} feeding the same
$W_{kj}$ of $R_6$, plus the horizon third-kind part $R_\Delta/\Delta$ contributing
$\ln\sigma$-quotients at $r_\pm$ (the coordinate-time analogue of the $\tau$-branch's
dipole at infinity)---no double cover. The full $t$-branch assembly (the
$W_{kj}$-reorganization plus the horizon piece) reproduces the direct double integral
$-\Ehat\!\int\!\partial_E F\,t\,\dd r$ to $9\times10^{-13}$.

The building blocks $W_{kj}$ are length-two iterated Abelian
integrals~\cite{Chen1977} of $y^2=\mathcal{R}^2$, of \emph{weight two}. Their elliptic
($g{=}1$) analogues are the polylogarithms of Beilinson--Levin and
Brown--Levin~\cite{BeilinsonLevin1994,BrownLevin2011}; a rigorous \emph{higher-genus}
(hyperelliptic) polylogarithm framework---base points, paths, tangential
regularization, monodromy and single-valued completion---is the subject of recent work
\cite{BauneEtal2024,DHokerSchlotterer2024,DHokerHiddingSchlotterer2025,DHokerEnriquezSchlottererZerbini2026},
and would be required before ``genus-two
dilogarithm'' can be used as a theorem-level term; we therefore adopt the neutral name
above and keep irreducibility and basis dimension conjectural. As iterated integrals
they do \emph{not} reduce to a weight-one Kleinian $\zeta,\sigma$
combination~\cite{BakerAbelian1897,BuchstaberEnolskii1997}, because for
genus two the Abel image of single points lies on the theta divisor
$\Theta=W_{g-1}$~\cite{MumfordTata}, where $\zeta,\wp$ are singular; the single-point
second-kind integrals are instead Fay theta-derivatives at shifted (off-divisor)
arguments. The one third-kind piece---the clock's dipole at infinity, residue
$\rho_0=M/(E^2-1)^{3/2}$---integrates to
\begin{equation}
    L(u)=\log\frac{\theta[\delta](u-e_+)}{\theta[\delta](u-e_-)}
    \label{eq:thirdkind-theta}
\end{equation}
(Fay prime form~\cite{Fay1973}, odd half-characteristic $\delta$, $e_\pm$ the Abel
images of the two points at infinity; differentiating $L(u(r))$ back reproduces the
third-kind differential to $5\times10^{-8}$). The $W_{kj}$ are length-two iterated Abelian integrals of the frozen curve; their
elliptic ($g{=}1$) analogues admit Kronecker--Eisenstein
representations~\cite{BeilinsonLevin1994}, but that reference does not by itself
establish a higher-genus construction (a genus-two framework needs base points, paths,
tangential regularization, monodromy and single-valued completion). The
weight-two content is genuine and \emph{conjecturally} irreducible to weight one, while
all coefficients remain
algebraic.

\paragraph{Numerical validation.} Table~\ref{tab:adiab-valid} compares the
hybrid form~\eqref{eq:phi-adiab}--\eqref{eq:psi-irr} (closed
$\mathcal{C}$ plus numerical $\psi$) against the directly integrated non-autonomous
coefficient $-\Ehat\!\int\partial_E F\,\eta\,\dd r$, for both branches
($M=1$, $a=0.9$, $\Ehat=1.4$; $t$-branch $J=6$, $\tau$-branch scattering $J=2.5$).
The agreement is at the discretization floor and scales linearly with
$\varepsilon$, confirming that the split is an exact reorganization (integration by
parts) rather than an approximation. Figure~\ref{fig:phi-closed} shows the same
comparison as orbit shapes in the plane, closed form against direct integration,
for both branches on compact sub-arcs bounded away from the turning point. The frozen family
itself already exhibits nontrivial ergosphere dynamics: for penetrating $t$-branch
angular momenta the orbit dips inside $r_e$, co-rotates, and turns just outside the
horizon, and a growing conformal factor eventually expels it from the ergosphere at
an adiabatic threshold (Fig.~\ref{fig:phi-penetrating}).

\paragraph{Exact first-order term and the true dynamic test.} Table~\ref{tab:adiab-valid}
and Fig.~\ref{fig:phi-closed} certify only the \emph{algebraic} identity (closed
$\mathcal{C}+\psi$ equals the direct integral of the same coefficient); they do not by
themselves test the adiabatic form against the actual non-autonomous dynamics. Integrating
the true non-autonomous optical-metric geodesic---the branch-Hamiltonian flow with
$E_{\rm eff}=\Ehat/A(\eta)$, $J_{\rm eff}=J/A(\eta)$---and comparing to
$\varphi_0+\varepsilon\,\delta\varphi$ shows the \emph{physical} accuracy is $\sim2\%$
(linear in $\varepsilon$), far above the algebraic floor
(Fig.~\ref{fig:phi-true-dynamic}). Linearizing the extended Hamiltonian
$\mathcal H=p_s+H=0$ about the frozen extremal---with the slow variable
$\lambda$ the elapsed clock measured \emph{from the free-arrival endpoint} (where
$H=0$ by the transversality~\eqref{eq:transversality}), oriented inward, and
$A(\eta)=A_0(1+\varepsilon\lambda+\cdots)$---gives, in the \emph{normalized} momentum
$P_r=p_r/A$ (\ref{app:firstorder}), the first-order shift of the momentum shell
$\delta P_r=(\lambda\,\Theta H-S_{\mathcal D})/H_{P_r}$ and
\begin{equation}
    \begin{aligned}
    &\frac{\delta\varphi}{\varepsilon}=\int\!\Big[\frac{\partial_{P_r}G}{H_{P_r}}\,
    \big(\lambda\,\Theta H-S_{\mathcal D}\big)-\lambda\,\Theta G\Big]\dd r,\qquad
    G=\frac{H_J}{H_{P_r}}=\varphi',\\[2pt]
    &S_{\mathcal D}(\lambda)=\!\int_0^{\lambda}\!\mathcal D H\,\dd\lambda',\qquad
    \mathcal D\equiv\underbrace{E_{\rm eff}\partial_{E_{\rm eff}}+J_{\rm eff}\partial_{J_{\rm eff}}}_{\Theta}
    +\,P_r\partial_{P_r},
    \end{aligned}
    \label{eq:adiab-exact}
\end{equation}
which closes the correction to $O(\varepsilon^2)$. The source uses the \emph{full} Euler
operator $\mathcal D$: the normalization $P_r=p_r/A$ is a time-dependent (conformally
symplectic) rescaling---\emph{not} canonical---so the canonical flow in the original
$(r,p_r)$ carries, after passing to $P_r$, a dilation term
$-\!\left(A'/A\right)P_r$ (\ref{app:firstorder}), and the shell-departure rate is $\mathcal D H=\Theta
H+P_r H_{P_r}$, \emph{not} $\Theta H$ alone (omitting $P_r H_{P_r}$ leaves an
$O(\varepsilon)$ error against the true flow---this is the correction of
\ref{app:firstorder}). Here $\Theta$ differentiates only the \emph{dynamical}
charges $E_{\rm eff}=\Ehat/A$, $J_{\rm eff}=J/A$ (the rail invariants $\Ehat,J_{\rm TK}$
are inert); for single-parameter Vaidya there is no rescaling, $\mathcal D=\Theta=m\partial_m$.
The source $S_{\mathcal D}$ is \emph{terminally anchored} ($S_{\mathcal D}=0$ at the
free-arrival endpoint); the reduction holds on compact sub-arcs bounded away from the
turning point ($H_{P_r}\to0$), and is verified to machine precision against the true
canonical flow (slope~$\approx2$;
\texttt{ThakurtaMetric/adiabatic\_first\_order\_exact.py}, \texttt{tests/}). This splits
into an \emph{on-shell} part (parameter drift, the closed form $\mathcal C+\psi$ of
Sec.~\ref{sec:thakurta}, weight one in Kleinian $\zeta,\sigma$) and an \emph{off-shell}
part driven by the dilation letter $\int P_r H_{P_r}\dd\lambda=\int p_{r0}\dd r$ (the
radial action, reduced in closed form in \ref{app:firstorder},
Eqs.~\eqref{eq:action-reduction}--\eqref{eq:ak}): a length-two iterated integral on the
same spectral curve, closing in $\zeta,\sigma$ \emph{plus} one genus-two dilogarithm
sourced by the third-kind letter at the seed Kerr null surface $\Delta=0$.

\paragraph{Universal form of the off-shell source.} The source $S_{\mathcal D}$ does
not, in fact, require the per-metric reduction to be written in closed form: it admits
a metric-independent structure fixed by the \emph{Randers} character of the
arrival-time Hamiltonian. Any such Hamiltonian is
$H=\beta(r,E)\,J+\sqrt{A_{rr}(r,E)\,p_r^2+A_{\varphi\varphi}(r,E)\,J^2}-1$, homogeneous
of degree one in $(p_r,J)$ apart from the normalization $-1$; Euler's theorem then
gives the identity
\begin{equation}
    \big(J\,\partial_J+P_r\,\partial_{P_r}\big)H=H+1,
    \label{eq:finsler-euler}
\end{equation}
valid for \emph{arbitrary} $\beta,A_{rr},A_{\varphi\varphi}$ (verified symbolically).
On the frozen shell $H=0$ this collapses the integrand of $S_{\mathcal D}$ to a
universal skeleton. For the \emph{conformal} family ($\mathcal D=\Theta+P_r\partial_{P_r}$),
$\mathcal D H=E H_E+(H+1)$, so
\begin{equation}
    S_{\mathcal D}(\lambda)=\lambda+\int_0^{\lambda}E_{\rm eff}H_{E_{\rm eff}}\,\dd\lambda',
    \label{eq:SD-conformal}
\end{equation}
the $\lambda$ term universal and $E_{\rm eff}H_{E_{\rm eff}}$ the single
metric-specific ``energy-Euler'' piece (fixed by $E\partial_E v^2=2(1-v^2)$). For the
\emph{mass-function} family (Vaidya, $\mathcal D=m\partial_m$), Schwarzschild
self-similarity $(r,m)\to\kappa(r,m)$ (with $E,p_r$ inert, $J\to\kappa J$) gives
$m H_m=-(rH_r+JH_J)$, and combining with~\eqref{eq:finsler-euler} and Hamilton's
equations ($p_rH_{p_r}-rH_r=\dd(rp_r)/\dd\lambda$) yields the still simpler
\begin{equation}
    S_{\mathcal D}(\lambda)=\big[\,r\,p_r\,\big]_0^{\lambda}-\lambda,
    \label{eq:SD-vaidya}
\end{equation}
a boundary term minus $\lambda$. The two deformation classes thus carry
\emph{different} universal sources---no single formula spans both, because the slow
parameter couples to the geometry differently (a conformal rescaling of all momenta
versus a self-similar mass shift)---yet both share the same universal $\pm\lambda$
skeleton from~\eqref{eq:finsler-euler}. Both~\eqref{eq:SD-conformal}
and~\eqref{eq:SD-vaidya} are verified to machine precision against the direct
$\int\mathcal D H\,\dd\lambda$ (\texttt{ThakurtaMetric/universal\_SD\_source\_check.py},
\texttt{VaidyaMetric/universal\_SD\_source\_vaidya.py}).

\paragraph{Fixed-charge versus fixed-endpoint.} The correction~\eqref{eq:adiab-exact}
is computed at \emph{fixed canonical charges} $(\Ehat,J)$: it is the same-launch,
fixed-charge response of the orbit to the slow drift, the object that the true-flow
test of \ref{app:firstorder} validates. A genuine two-point boundary-value
problem with \emph{fixed spatial endpoints} instead lets the conserved charge float,
$J=J_0+\varepsilon\,J_1+O(\varepsilon^2)$, with $J_1$ fixed by holding the arrival
azimuth,
\begin{equation}
    0=\delta\varphi_{\rm dyn}+J_1\,\partial_J\varphi_0
    \;\Longrightarrow\;
    J_1=-\,\frac{\delta\varphi_{\rm dyn}}{\partial_J\varphi_0},
    \label{eq:J1-fixedendpoint}
\end{equation}
where $\delta\varphi_{\rm dyn}$ is the fixed-charge displacement above. The
fixed-endpoint shape then adds $\varepsilon J_1\,\partial_J\varphi_0$, which merely
re-selects the launch angle and does not alter the transcendental class. We report
the fixed-charge form throughout and quote~\eqref{eq:J1-fixedendpoint} as the explicit
map to the fixed-endpoint problem.

\begin{table}[htbp]
    \centering
    \caption{Hybrid closed form versus the direct integral of the \emph{same}
    first-order coefficient, $\max_r|\varphi_{\rm hybrid}-\varphi_{\rm full}|$ (an
    integration-by-parts identity---an algebraic self-consistency, \emph{not} a test
    against the true dynamics; for the latter see Fig.~\ref{fig:phi-true-dynamic}).
    Both branches; the residual is the numerical-quadrature floor.}
    \label{tab:adiab-valid}
    \begin{tabular}{lcc}
        \hline\hline
        $\varepsilon=A'/A$ & $t$-branch ($\eta=t$) & $\tau$-branch ($\eta=\tau$)\\
        \hline
        $0.005$ & $1.4\times10^{-6}$ & $2.5\times10^{-8}$\\
        $0.02$  & $5.6\times10^{-6}$ & $1.0\times10^{-7}$\\
        $0.06$  & $1.7\times10^{-5}$ & $3.0\times10^{-7}$\\
        \hline\hline
    \end{tabular}
\end{table}

\begin{figure}[tbp]
    \centering
    \includegraphics[width=\textwidth]{Immagini/fig_phi_closed_validation.pdf}
    \caption{Closed-form adiabatic orbit shape $\varphi(r,A)$ against the
    \emph{direct quadrature of the same on-shell coefficient}, for both branches
    (this is an integration-by-parts self-consistency check, \emph{not} a test
    against the true non-autonomous dynamics; for the latter see
    Fig.~\ref{fig:phi-true-dynamic}). Panels (a,b): the closed form
    $\varphi_0+\varepsilon\,\delta\varphi$ (dashed), with $\delta\varphi$ assembled
    from Eq.~\eqref{eq:psi-split}, overlies the direct quadrature of the same
    first-order coefficient (solid) on compact sub-arcs bounded away from the
    turning point (marker); the orbits wind close to the ergosphere $r_e$. Panel
    (c): the residual $\max_r|\varphi_{\rm closed}-\varphi_{\rm quad}|$ grows
    approximately \emph{linearly} with $\varepsilon=A'/A$ over the tested range
    (residual$/\varepsilon\approx2.8\times10^{-4}$), for the $t$-branch
    ($J=6$) and $\tau$-branch ($J=2.5$). This linear growth measures the size of
    the retained first-order correction and a fixed numerical floor; it is
    \emph{not} an $O(\varepsilon^2)$ truncation test. The dedicated log--log
    convergence study against the full non-autonomous flow (slope two once the
    complete first-order, on-shell $+$ off-shell, term is used) is supplied in
    Fig.~\ref{fig:phi-true-dynamic}. All coefficients in $\delta\varphi$ are algebraic
    and given, not fitted.}
    \label{fig:phi-closed}
\end{figure}

\begin{figure}[tbp]
    \centering
    \includegraphics[width=\textwidth]{Immagini/fig_phi_validation_true_dynamic.pdf}
    \caption{\emph{Physical} validation against the \emph{true canonical flow}
    (Eq.~\eqref{eq:normflow}, with the dilation term $-\!\left(A'/A\right)P_r$). Left:
    the actual optical-metric geodesic ($t$-branch) versus the frozen shape plus the
    \emph{leading} on-shell correction (dashed) and plus the \emph{exact}
    Eq.~\eqref{eq:adiab-exact} correction ($S_{\mathcal D}$, dotted), at $A'/A=0.04$.
    Right: the residual on a compact sub-arc bounded away from the turning point, for
    \emph{both branches} and \emph{two parameter sets} ($t$: $a{=}0.9,\Ehat{=}1.4,J{=}6$;
    $\tau$: $a{=}0.9,\Ehat{=}1.4,J{=}2.5$ scattering; $t$: $a{=}0.5,\Ehat{=}1.3,J{=}5$).
    The on-shell term alone (dotted; source $\int\Theta H\dd\lambda$ without the
    $P_r H_{P_r}$ dilation) has an $O(\varepsilon)$ physical error (slope $\approx1.0$);
    the exact term (solid; source $S_{\mathcal D}$) closes it to $O(\varepsilon^2)$, with
    least-squares log--log slopes $2.05$ ($t$), $2.07$ ($\tau$) and $2.01$ (second set),
    far above the algebraic integration-by-parts floor of
    Fig.~\ref{fig:phi-closed}/Table~\ref{tab:adiab-valid} ($\sim10^{-6}$). The
    $\sim2\%$ is precisely the off-shell (costate) residual of the leading
    approximation.}
    \label{fig:phi-true-dynamic}
\end{figure}

\begin{figure}[tbp]
    \centering
    \includegraphics[width=\textwidth]{Immagini/fig_phi_penetrating.pdf}
    \caption{Ergosphere-penetrating $t$-branch orbits ($J=3<J_c^+$), frozen family.
    Panel (a): as the conformal factor $A$ grows ($E_{\rm eff}=\Ehat/A$ drops) the
    \emph{analytically continued extremal} (inside $r_e$ it is a continuation of the
    exterior extremal, not a proved PMP solution) \emph{enters the ergosphere}
    ($r<r_e$), is forced to co-rotate by frame dragging, and turns ($\circ$) in the
    annulus $r_+<r_{\rm t}<r_e$---\emph{outside} the horizon---before exiting; the
    rail turning point $R_6=0$ prevents it from ever reaching $r_+$. Panel (b): the
    penetration depth $r_{\rm t}(A)$ stays just above $r_+$ until a frozen-family
    threshold $A_c\!\simeq\!1.5$ (each curve is frozen at fixed $A$),
    beyond which the orbit no longer penetrates
    ($r_{\rm t}$ jumps outward): the breathing conformal factor expels the orbit
    from the ergosphere.}
    \label{fig:phi-penetrating}
\end{figure}

\paragraph{Penetration phase diagram.} The frozen ergosphere dynamics is organized
by two \emph{analytic} thresholds, both of which we verify against direct
root-tracking (Fig.~\ref{fig:penetration-phase}). Writing $DE=(E_{\rm eff}^2-1)r+2M$,
for $E_{\rm eff}<1$ a \emph{freezing wall} $r_w=2M/(1-E_{\rm eff}^2)$ (the zero of
$DE$, and of $R_6$) advances inward as $A$ grows; the launch radius $r_0$ is
swallowed at
\begin{equation}
    A_c^{\rm wall}=\frac{\Ehat}{\sqrt{1-2M/r_0}},
    \label{eq:Awall}
\end{equation}
beyond which no orbit issues from $r_0$. Below the wall the $t$-orbit crosses the
ergosphere iff its outer turning lies inside $r_e$; the prograde boundary is
\begin{equation}
    J_c^+(A)=\frac{2M^2}{a}+a+\frac{a\,A^2}{2\Ehat^2}
    \label{eq:Jcplus}
\end{equation}
(matching the numerics to $10^{-15}$), the retrograde one a saddle-node $J_c^-(A)$. The window $(J_c^-,J_c^+)$ splits into a
\emph{plunge} core (no turning, the orbit reaches $r_+$) and a thin
\emph{penetrate-and-bounce} rim just below $J_c^+$, where the orbit dips inside
$r_e$, co-rotates, and turns in the annulus $r_+<r_{\rm t}<r_e$ before exiting
(Fig.~\ref{fig:bounce}). The breathing conformal factor thus \emph{expels} orbits
from the ergosphere at $A_c^{\rm wall}$: a purely non-stationary effect.

\begin{figure}[tbp]
    \centering
    \includegraphics[width=0.82\columnwidth]{Immagini/fig_penetration_threshold.pdf}
    \caption{Ergosphere-penetration phase diagram of the frozen $t$-branch in the
    $(A,J)$ plane ($M=1$, $a=0.9$, $\Ehat=1.4$, $r_0=12$). Four regimes: plunge to
    the horizon, penetrate-and-bounce, scattering, and forbidden (beyond the
    freezing wall $A_c^{\rm wall}=1.534$, dashed). The prograde ergosphere boundary
    $J_c^+(A)$ [Eq.~\eqref{eq:Jcplus}] is analytic (solid, verified to
    $10^{-15}$); the penetrating window is $(J_c^-,J_c^+)$.}
    \label{fig:penetration-phase}
\end{figure}

\begin{figure}[tbp]
    \centering
    \includegraphics[width=0.9\columnwidth]{Immagini/fig_penetrate_bounce_orbit.pdf}
    \caption{A penetrate-and-bounce $t$-brachistochrone ($J=3.2$, frozen): it
    enters the ergosphere, is forced to co-rotate (red arc), turns at
    $r_{\rm t}=1.87$---inside $r_e$ but outside the horizon $r_+$---and exits. The
    rail turning $R_6=0$ guarantees it never reaches $r_+$; crossing the ergosphere
    and returning is causally allowed.}
    \label{fig:bounce}
\end{figure}




% ======================================================================
\section{Equatorial closed forms and the ergosphere}
\label{sec:closed}
% ----------------------------------------------------------------------

All closed forms in this section and in \ref{app:explicit} are obtained at
\emph{frozen conformal factor}: constant $A$, equivalently stationary Kerr with
the substitution $E\to E_{\rm eff}=\Ehat/A$ (and $J\to J/A$, the spin $a$ retained) in every
parameter (Sec.~\ref{sec:thakurta}). They are therefore the exact
\emph{instantaneous} orbit shapes on a Thakurta--Kerr slice of given $A$; the
running $A(\eta)$ makes the full problem non-autonomous---the quasi-constant
drift of Sec.~\ref{sec:quasiconst}, $K^{\rm TK}(\eta)=K^{\rm Kerr}(E_{\rm
eff}(\eta))+O(A'/A)$---so a genuinely time-dependent trajectory is recovered by
letting $E_{\rm eff}=E_{\rm eff}(\eta)$ run through this family, not by a single
autonomous curve.

\subsection{Separatrix in Weierstrass functions}
\label{sec:separatrix}
From the doranTau factorization $\Delta-\Jc^2 w=f(r^2+c^2)$, $c=a/E$
\cite{RosignoliKerr}, the
equatorial $\tau$-separatrix $\varphi(r)$ is genus~1, closed in the
Weierstrass functions $\wp,\sigma,\zeta$
\cite{WhittakerWatson,DLMF,Hackmann2008,Hackmann2010}: the angle is linear in the
uniformizing variable $z(r)=\wp^{-1}(A/r+B)$ plus two $\sigma$-ratios
[\ref{app:explicit}, Eq.~\eqref{eq:sep-phi}], regular \emph{through} the
ergosphere because the factor $r^2+c^2$ divides the turning quartic and removes
$\re$ from the radicand. The Boyer--Lindquist form winds logarithmically at the
horizon; the Doran shift---a second, quasi-lemniscatic elliptic
curve---cancels it, so the ergosphere-penetrating separatrix
$\varphi_D=\varphi_{\rm BL}-\varphi_{\rm shift}$ is closed on two elliptic curves.

What controls the reduction is the \emph{genus} of the turning radical, not the
rationality of $\mathcal K$. Being Riemannian-pure
(Sec.~\ref{sec:thakurta}, the gravitomagnetic wind cancels), the $\tau$-branch
factorizes as doranTau---the factor $r^2+c^2$ dividing the quartic---dropping a
generically genus-2 reduction to genus~1. Generic $t$-branch orbits carry the
wind and have the sextic radical
$\mathcal R(r)=r\,Q_2(r)\,[(E^2-1)r+2M]$ (genus~2, hyperelliptic), with no such
cancellation. The $t$-brachistochrone nonetheless has its own genus-1 closed-form
separatrix, by a \emph{different} mechanism: at the marginal-capture (saddle-node)
momentum $J_c^-$ [Eq.~\eqref{eq:Jc-minus}] the sextic develops a \emph{double
root}, $\mathcal R=(r-r_*)^2 Q_4$ with $Q_4$ quartic (whereas the prograde
ergosphere crossing $J_+^t$ has only a simple root and stays genus~2). The double
root drops the radical to $Q_4$, and $\varphi(r)$ closes in the same
Weierstrass $\wp,\sigma,\zeta$ form as the $\tau$-separatrix
(\ref{app:explicit}). Both separatrices are thus genus~1---the $\tau$ one
via the doranTau factorization, the $t$ one via a double root---while generic
(non-separatrix) $t$-orbits are genuinely genus~2, closed only in genus-2
Kleinian functions.
The Kerr geodesic structure and its constants are classical
\cite{Carter1968,Chandrasekhar}. Validated end-to-end to
$2.8\times10^{-10}$ through the ergosphere, and general across parameters.

\begin{figure}[ht]
    \centering
    \includegraphics[width=\textwidth,height=0.68\textheight,keepaspectratio]{Immagini/fig_separatrix_3traiettorie.png}
    \caption{Prograde $\tau$-branch separatrix $J=+\Jc$ (the single penetrating
    value): ODE, quadrature and evaluated Weierstrass agree to
    $2.8\times10^{-10}$; the PMP-$\eta$ terminates at $\re$.}
    \label{fig:sep3}
\end{figure}


\subsection{Trichotomy and the ergosphere cusp}
\label{sec:trichotomy}
For the equatorial $\tau$-branch, $\dd\varphi/\dd r\sim K\sqrt{r-\re}$ as
$r\to\re$ (because $\sqrt f\to0$ while the turning radicand stays finite),
giving
\begin{equation}
    \varphi-\varphi_e \sim \tfrac{2}{3}K\,(r-\re)^{3/2},
    \label{eq:cusp}
\end{equation}
a cusp with radial tangent. The trichotomy in $|J|$ (sign of $a^2-J^2E^2$):
smooth periapsis above $\re$ ($|J|>\Jc$); marginal contact with $\re$
($|J|=\Jc=a/E$, separatrix); cusp reflection at $\re$ ($|J|<\Jc$). Only the
\emph{prograde} marginal $J=+\Jc$ actually crosses: inside $\re$ frame dragging
forces co-rotation ($\dd\varphi/\dd\lambda>0$), so the retrograde marginal
$J=-\Jc$, which would counter-rotate, reflects with a cusp instead. The crossing
(prograde) branch is continued through $\re$ to the horizon in Doran
(horizon-penetrating) coordinates
\cite{Doran2000,Natario2009,MartelPoisson2001},
$\varphi_D=\varphi_{\rm BL}-\varphi_{\rm shift}$.

\paragraph{Retrograde branch and the $\tau/t$ asymmetry.} The trichotomy above is
a $\tau$-branch statement. The set that \emph{reaches} $\re$ is symmetric in the
sign of $J$---the doranTau factorization depends on $J$ only through $J^2$, so the
ingoing orbit reaches $\re$ for $|J|\le\Jc=a/E$, prograde and retrograde sharing
the same threshold despite frame dragging---but the set that \emph{penetrates} is
not: only the single prograde $+\Jc$ crosses, the co-rotation constraint breaking
the $\pm J$ symmetry at the wall. The $t$-branch is not protected. Its turning function (from
$H_\eta=0$ at $p_r=0$, with $b=2Ma/r$) is
\begin{equation}
    N(r)=\big(\bar P-Jb\bar v^2\big)^2-J^2\Delta\bar v^2,
    \label{eq:tbranch-turn}
\end{equation}
which is only linear-plus-quadratic in $J$. At the ergosphere
$N(\re)=\bar P_e\,(\bar P_e-2aJ)$ with
$\bar P_e=\bar P(\re)=4M^2+2a^2+a^2/E^2$, so the $t$-orbit crosses $\re$ at the
single \emph{prograde} value
\begin{equation}
    J_+^{t}=\frac{\bar P_e}{2a}=\frac{4M^2+2a^2+a^2/E^2}{2a};
    \label{eq:Jplus-t}
\end{equation}
there is no retrograde $r_e$-crossing. Instead the retrograde $t$-capture is
bounded by the \emph{marginal} (double-root) condition $N=N'=0$ at a radius
$r_*>\re$---the massive analogue of the retrograde photon sphere---which closes
to
\begin{equation}
    J_c^{-}=-\,\frac{\bar P}{\bar v\,(\sqrt\Delta-b\bar v)}\bigg|_{r_*}<0 .
    \label{eq:Jc-minus}
\end{equation}
For $a=0.9$, $E=1.2$ (and $A=1$) these give $J_+^{t}=3.43$ and $J_c^{-}=-8.05$
($r_*=3.51$): the $t$-branch captures the wide, asymmetric band
$J\in[J_c^{-},J_+^{t}]$, far more capture-prone on the retrograde side. A
moderately retrograde orbit ($J=-5$, say) is therefore captured on the
$t$-branch but scattered on the $\tau$-branch (Fig.~\ref{fig:master-pen}). For
constant $A$ the results carry
over by $E\to\Ehat/A$, $J\to J/A$ (spin $a$ retained); in particular
$\Jc^{\tau}=sA^2/\Ehat$.

\begin{figure}[ht]
    \centering
    \includegraphics[width=\textwidth,height=0.34\textheight,keepaspectratio]{Immagini/fig_master_penetration_taut.png}
    \caption{Ergosphere behaviour of the two branches versus $J$
    (Table~\ref{tab:penetration}). The $\tau$-branch (bottom bar) is a
    \emph{trichotomy}: smooth bounce for $|J|>J_c$, cusp bounce at $r_e$ for
    $|J|<J_c$ (red), and a \emph{single} penetrating value $J=+J_c$ (star)---the
    reality wall $\sqrt{wf}$ and forced co-rotation block all others. The
    $t$-branch (top bar) is a \emph{dichotomy}: it penetrates the whole interval
    $(J_c^-,J_+^t)$ (green), reflecting smoothly outside.}
    \label{fig:master-pen}
\end{figure}

\begin{figure}[ht]
    \centering
    \includegraphics[width=\textwidth,height=0.68\textheight,keepaspectratio]{Immagini/fig_thakurta_cuspide_ergosfera.png}
    \caption{Behaviour at the ergosphere: cusp reflection ($|J|<\Jc$), smooth
    periapsis ($|J|>\Jc$), and the \emph{prograde} marginal $J=+\Jc$ crossing to
    $\rp$ in Doran coordinates---only $+\Jc$ crosses, the retrograde $-\Jc$
    cusping instead. Power law $2/3$ (cusp) vs $2$ (periapsis).}
    \label{fig:th-cusp}
\end{figure}

\subsection{Plunge inversion: rotational and conformal}
\label{sec:inversion}
Two mechanisms invert the $t/\tau$ ordering. Rotationally, in the $(a,J)$ plane,
$\rho(r_*)^2=H_{t0}/H_{\tau0}$, extending the stationary Kerr analysis
\cite{RosignoliKerr}. Conformally, the inversion locus is
\begin{equation}
    \mathrm{cond}(r,A)=A^2 b\,Q-\bar P(\Ehat-A^2 f)=0,\quad
    Q=b\bar v^2+\sqrt{\bar v^2\Delta},
    \label{eq:condinv}
\end{equation}
with $\mathrm{cond}(1)<0$ and
$\mathrm{cond}(\Af)=\Ehat r\Delta(\Ehat-1)/(r-2M)$. Under the explicit hypotheses
$\Ehat>1$ (scattering energy), $r>2M$ and $\Delta>0$ (exterior, off-horizon), the
latter is $>0$, so by the intermediate value theorem a root
$A_{\rm inv}\in(1,\Af)$ exists; it is reachable when $A_{\rm inv}<\Af$, with far
field $A_{\rm inv}\to\sqrt{\Ehat}$. For $\Ehat\le1$ the sign is non-positive and no
such root is guaranteed: the conformal inversion is thus a \emph{conditional}
statement, holding in the stated energy/domain window rather than universally.

\begin{figure}[ht]
    \centering
    \includegraphics[width=\textwidth,height=0.68\textheight,keepaspectratio]{Immagini/fig_thakurta_kerr_plunge_t_tau.png}
    \caption{Rotational inversion in the $(a,J)$ plane; ODE locus vs analytic
    $\rho(r_*)^2=H_{t0}/H_{\tau0}$.}
    \label{fig:th-plungeAJ}
\end{figure}

\begin{figure}[tbp]
    \centering
    \includegraphics[width=\textwidth,height=0.68\textheight,keepaspectratio]{Immagini/fig_thakurta_kerr_inversione_AJ.png}
    \caption{Conformal inversion in the $(J,A)$ plane: $\Delta r$ changes sign;
    analytic $\mathrm{cond}=0$ matches the numerical contour; reachable
    ($A_{\rm inv}<\Af$).}
    \label{fig:th-invAJ}
\end{figure}


\subsection{Fixed-endpoint comparison and protocol dependence}
\label{sec:bvp}
The statement ``branch $X$ plunges deeper'' is protocol dependent. With
\emph{fixed initial data} (same $r_0,J,p_r$) the $t$-branch is deeper
($r_{\min}^t<r_{\min}^\tau$). With \emph{fixed endpoints} (the genuine
two-point brachistochrone), the two branches join the same $A,B$ but are
distinct paths, and the ordering reverses: $r_{\min}^\tau<r_{\min}^t$. The
optical-index ratio $n_t/n_\tau=E/f>1$ penalizes deep radii more heavily for
the $t$-branch, so the $t$-brachistochrone bows outward while the $\tau$-branch
dips. This must be specified when comparing depths.

\begin{figure}[ht]
    \centering
    \includegraphics[width=\textwidth,height=0.68\textheight,keepaspectratio]{Immagini/fig_bvp_estremi_fissi.png}
    \caption{Fixed-endpoint $t$ and $\tau$ brachistochrones (Schwarzschild,
    $E=1.2$): symmetric (top) and asymmetric (bottom) endpoints. Both branches
    join the same $A,B$; the $\tau$-branch dips deeper.}
    \label{fig:bvp}
\end{figure}

Crucially, neither driver of same-launch inversion survives the
fixed-endpoint constraint. Scanning the spin $a$ over the entire
Kerr flow with fixed symmetric endpoints $(r_0,\pm\Phi)$ leaves
$\Delta r=r_{\min}^t-r_{\min}^\tau$ strictly positive
($\Delta r\in[+0.45,+2.06]$ on $a\in[0.05,0.95]$); the map is nearly flat
in $a$, its gradient set almost entirely by the endpoint angle
$\Phi$ (encoded by $J$). The conformal scan (via the transfer
$E_{\rm eff}=\hat E/A$) is likewise everywhere positive
($\Delta r\in[+0.002,+3.711]$). The pointwise index bound $n_t/n_\tau=E/f>1$ does
not by itself order the minima of two distinct functionals---indeed at \emph{fixed}
$J$ the ordering is the opposite (the $t$-orbit turns deeper), and the fixed-endpoint
result is a reversal driven by the accumulated-angle map $J\mapsto\Phi$. It admits,
however, a clean reduction. Parametrize each branch by its turning radius $r_{\min}$
(which fixes $J$ through the turning condition $r\Delta=J^2 DE$ for $\tau$, $Q_2=0$
for $t$), and let
\begin{equation}
    \Phi_{\rm br}(r_{\min})=\int_{r_{\min}}^{r_0}\!F_{\rm br}\big(r;J_{\rm br}(r_{\min})\big)\,\dd r
\end{equation}
be the half-angle swept by the orbit turning at $r_{\min}$. Two lemmas---verified
across $a\in[0.3,0.9]$, $E\in[1.05,1.5]$---close the result. \emph{(A)} At matched
turning radius the $t$-integrand dominates \emph{pointwise},
$F_t(r;J_t)>F_\tau(r;J_\tau)$ for $r\in(r_{\min},r_0)$; since $J_{t,\tau}(r_{\min})$
are \emph{algebraic}, this is a polynomial (not transcendental) inequality, whence
$\Phi_t(r_{\min})>\Phi_\tau(r_{\min})$. \emph{(B)} $\Phi_{\rm br}$ is strictly
decreasing. Together: at equal angle
$\Phi_\tau(r_{\min}^\tau)=\Phi_t(r_{\min}^t)>\Phi_\tau(r_{\min}^t)$, so
$r_{\min}^\tau<r_{\min}^t$ by (B)---the fixed-endpoint no-inversion. Lemma~(A) is
\emph{proven in full}. Writing $F_t^2-F_\tau^2=(r\,DE/\Delta^2)\,N_G/(Q_2 P_\tau)$
with $N_G=(J_t(r-2M)+2Ma)^2 P_\tau-J_\tau^2(r-2M)Q_2$ and reducing modulo the turning
conditions ($J_\tau^2=r_{\min}\Delta_{\min}/DE_{\min}$, $Q_2(r_{\min},J_t)=0$), one
finds the key factorization
\begin{equation}
    N_G=(r-r_{\min})\,\frac{4r\,\Delta(r)}{r_{\min}-2M}\,W(r),\qquad W\ \text{\emph{linear} in }r .
\end{equation}
The prefactor is positive for $r>r_{\min}>2M$, so (A) reduces to $W(r)>0$; and a
linear $W$ with $W(r_{\min})>0$ and positive slope stays positive for all
$r>r_{\min}$. Both $W(r_{\min})=(r_{\min}-2M)\,a\,DE_{\min}(J^\ast-J_t)$ and the slope
$w_1=a\,DE_{\min}(J^{\ast\ast}-J_t)$ are positive because the prograde root $J_t$ of
$Q_2(r_{\min},\cdot)$ lies below both thresholds
$J^\ast=[E^2r_{\min}(a^2+r_{\min}^2)+2a^2]/(a\,DE_{\min})$,
$J^{\ast\ast}=[E^2r_{\min}^2(r_{\min}-M)+E^2a^2r_{\min}+a^2]/(a\,DE_{\min})$: the
quadratic opens downward ($A=-(r_{\min}-2M)DE_{\min}<0$) with negative vertex
$-2a/(r_{\min}-2M)$, while $J^\ast,J^{\ast\ast}>0$ and
\[
\begin{aligned}
Q_2(r_{\min},J^\ast)&=\frac{-E^2r_{\min}^3(E^2r_{\min}^2+a^2)\Delta_{\min}}{a^2DE_{\min}},\\[2pt]
Q_2(r_{\min},J^{\ast\ast})&=\frac{-(r_{\min}-2M)(E^2r_{\min}^2+a^2)\big(E^2r_{\min}^2[(r_{\min}-M)^2+a^2]+a^2\big)}{a^2DE_{\min}}
\end{aligned}
\]
are both manifestly negative, placing the thresholds beyond the larger root
(\texttt{no\_inversion\_reduction.py}). Hence (A) is a theorem for all $r_{\min}>2M$,
$E>1$, $a>0$. Lemma~(B), the monotonicity $\Phi_\tau'(r_{\min})<0$, is
\emph{conditional}: $\Phi_\tau$ has a single maximum at a deflection radius
$r_{\rm pk}\approx3M$ (the strong-deflection/winding threshold, where the orbit
approaches the unstable circular $\tau$-orbit) and is strictly decreasing for
$r_{\min}>r_{\rm pk}$. In the effective-potential form
$\Phi_\tau=\sqrt{V_{\min}}\int_{r_{\min}}^{r_0}K/\sqrt{V-V_{\min}}\,\dd r$ with
$V=r\Delta/DE$, $K=\sqrt{r(r-2M)}/\Delta$, $V_{\min}=V(r_{\min})$, the substitution
$V=V_{\min}+\sigma^2$ and one integration by parts give, for $W\equiv K/V'$
($W>0$, $W'<0$),
\begin{equation}
    \frac{\dd\Phi_\tau}{\dd V_{\min}}=
    \frac{(V_0-2V_{\min})\,W(V_{\min})}{\Sigma\sqrt{V_{\min}}}
    +\frac{2}{\sqrt{V_{\min}}}\!\int_0^{\Sigma}\!W'(V_{\min}+\sigma^2)\,g(\sigma)\,\dd\sigma,
\end{equation}
$\Sigma=\sqrt{V_0-V_{\min}}$, where $g(\sigma)=(V_0-2V_{\min})\sigma/\Sigma+V_{\min}-\sigma^2$
satisfies $g(0)=V_{\min}>0$, $g(\Sigma)=0$ with its other root negative, hence
$g>0$ on $[0,\Sigma)$; the second term is therefore \emph{manifestly negative}.
When $V(r_0)\le2V(r_{\min})$ (the grazing regime, $r_{\min}$ near $r_0$) the first
term is $\le0$ as well, and $\Phi_\tau'<0$ is \emph{proven in closed form}. The same
closed form---using the exact identity (integration by parts on the Abel integral)
\begin{equation}
    \sqrt{V_{\min}}\,\Phi_\tau'=\frac{W(V_0)(V_0-2V_{\min})}{\Sigma}
    +\int_{V_{\min}}^{V_0}\!\frac{W'(t)\,(2V_{\min}-t)}{\sqrt{t-V_{\min}}}\,\dd t ,
    \label{eq:phitau-closed}
\end{equation}
together with the convexity of $W$ ($W''>0$)---extends the proof to the quarter
regime $V(r_0)\le4V(r_{\min})$. We accordingly \emph{define}
\begin{equation}
    R_\ast(E,a)=\sup\big\{\,r_0:\ V(r_0)\le 4\,V(r_{\min}(r_0))\ \text{along the marginal
    turning family}\,\big\},
    \label{eq:Rstar}
\end{equation}
with $V=r\Delta/DE$ the turning potential---an explicit $(E,a)$-dependent radius (it is
\emph{not} a function of $E$ alone, since $V$ carries the spin through $\Delta$)---so the
closed-form proof covers all endpoints $r_0\le R_\ast(E,a)$. In the
static limit $a\to0$ (Vaidya $\dot m=0$ / frozen Thakurta) we can go further:
$W(V)=\tfrac1{2V}-\tfrac1{2b^{3/2}V^{3/2}}+O(V^{-2})$ with $b=E^2-1$, so the
large-$r_0$ angle is
$\Phi_\tau=\arctan\!\sqrt{V_0/V_{\min}-1}-b^{-3/2}\sqrt{1/V_{\min}-1/V_0}+\cdots$,
exhibiting a \emph{single, non-degenerate} maximum at $V_{\rm pk}\sim b^{-3/2}\sqrt{V_0}$
with $\Phi_\tau''(V_{\rm pk})=-\tfrac12 b^{9/4}V_0^{-5/4}<0$; hence (B) holds for all
sufficiently large $r_0$. For intermediate $r_0$ the single-crossing property
(at every $r_{\min}$, \emph{either} $\Phi_\tau'<0$ \emph{or} $\Phi_\tau''<0$) is
established by a validated interval-arithmetic certificate, completed for $r_0=10M$
over the entire continuum $r_{\min}\in(2.6,6.05)M$ (the elementary grazing bound
covering $r_{\min}\gtrsim6.05M$; \texttt{no\_inversion\_schwarzschild\_CAP\_r0\_10.py}).
A uniform symbolic proof for all $(a,E,r_0)$ remains open---the deciding inequality is
tight at the \emph{transcendental} deflection threshold $r_{\rm pk}$. We therefore state
(B), and hence the fixed-endpoint no-inversion, as a \emph{theorem for $r_0\le R_\ast(E,a)$
and asymptotically as $r_0\to\infty$ in the static case}, machine-certified at
representative $r_0$, and a numerically supported conjecture in full generality
(verified over $a\in[0.2,0.9]$, $E\in[1.05,1.5]$); the fixed-endpoint turning points of
Fig.~\ref{fig:bvp} ($r_{\min}^{\tau,t}\gtrsim4.7M$) lie in the certified scattering
regime.
Table~\ref{tab:protocol} summarizes the resulting protocol dependence:
inversion is a property of the \emph{comparison protocol}, not an invariant
of the two brachistochrones---it requires the freedom to let the endpoint
move.

\begin{figure}[tbp]
    \centering
    \includegraphics[width=\textwidth,height=0.68\textheight,keepaspectratio]{Immagini/fig_colormap_spin_estremi_fissi.png}
    \caption{Numerical spin scan at fixed symmetric endpoints (Kerr, Hamiltonian
    flow): $\Delta r$ over the sampled $(J,a)$ domain stays positive---spin
    does not invert once the endpoints are pinned (a finite scan, consistent with the
    tiered no-inversion status of \S\ref{sec:inversion}). The gradient is set by the
    endpoint angle ($J$), not by $a$.}
    \label{fig:cm-spin-fix}
\end{figure}


The verdict is independent of the endpoint symmetry: repeating both scans
with \emph{asymmetric} fixed endpoints ($r_A=10,r_B=6$) leaves
$\Delta r>0$ everywhere as well (spin $[+0.76,+2.64]$, conformal
$[+0.16,+3.24]$; Figs.~\ref{fig:asym-spin}--\ref{fig:asym-conf}),
consistent with the pointwise bound $n_t/n_\tau=E/f$.

\begin{table}[t]
    \centering
    \caption{Protocol dependence of plunge inversion. Same-launch
    comparison inverts under either rotation or a conformal factor; the
    genuine fixed-endpoint brachistochrone shows no inversion in every
    proven and tested regime (proven for $r_0\le R_\ast$ and, in the static
    case, as $r_0\to\infty$; machine-certified at representative $r_0$;
    numerically verified elsewhere).}
    \label{tab:protocol}
    \begin{tabular}{@{}lcc}
        \br
        protocol & spin $a$ & conformal $A$ \\
        \mr
        same launch & inverts (\S\ref{sec:inversion}) & inverts (\S\ref{sec:inversion}) \\
        fixed endpoints & no inversion$^\ast$ & no inversion$^\ast$ \\
        \mr
        \multicolumn{3}{@{}l}{\footnotesize $^\ast$Tiered status (see caption); the
        pointwise bound $n_t/n_\tau=E/f>1$ is \emph{necessary but not sufficient}.}\\
        \br
    \end{tabular}
\end{table}

\subsection{The breathing indicatrix: a unified adiabatic theorem}
\label{sec:unified}
The two closed-form corrections---Thakurta--Kerr, Eq.~\eqref{eq:psi-split}, and
Vaidya, Eq.~\eqref{eq:vaidya-deltaphi}---are one statement in disguise, and the
indicatrix ellipse of Sec.~\ref{sec:indicatrix} is the reason. Let $F(r;\lambda)$
be the frozen brachistochrone integrand of a one-parameter family of rail
indicatrices, $\lambda$ the slow datum (scale $a$, mass $m$, conformal $A$), on the
genus-two spectral curve $y^2=S(r;\lambda)$. Whenever $\lambda=\lambda(\text{clock})$
drifts slowly,
\begin{equation}
    \varphi(r)=\varphi_0\!\big(r;\lambda\big)+\dot\lambda\!\int\!\partial_\lambda
    F\,(\text{clock})\,\dd r+O(\dot\lambda^2)
    =\varphi_0+\dot\lambda\Big[\sum_{k<j} q_{kj}\,W_{kj}+(\text{elementary})\Big],
    \label{eq:unified}
\end{equation}
with $q_{kj}$ \emph{algebraic} (rational in the metric parameters, via the exact
reduction $\partial_\lambda F=N_\lambda/S^{3/2}$) and $W_{kj}=\int(U_j\dd U_k-U_k
\dd U_j)$ length-two iterated Abelian integrals of the frozen curve. Equation
\eqref{eq:unified} is the \emph{on-shell} component of the first-order response
(the parameter drift of the frozen family); the \emph{complete} first-order term is
the extended-Hamiltonian expression Eq.~\eqref{eq:adiab-exact}, whose off-shell
(costate $p_\eta$) piece is itself a length-two iterated integral of the same class
(a $\sim2\%$ correction here). Within the on-shell component the transcendence
structure is \emph{universal} across the three geometries. The three geometries are distinguished \emph{not} by which datum
of the indicatrix varies---in each, \emph{both} the centre $c$ and the semi-axes $R$
breathe, through $\partial_\lambda F=\partial_c F\,\partial_\lambda c+\partial_R
F\,\partial_\lambda R$ with \emph{both} terms nonzero (for Vaidya, explicitly,
$\partial_m c_r=-2/r$, $\partial_m R_r=E\,[r(2m+(E^2-1)r)]^{-1/2}$,
$\partial_m R_\varphi=[r^3(2m+(E^2-1)r)]^{-1/2}$; likewise every $E_{\rm
eff}=\Ehat/A$-dependent datum of Thakurta--Kerr varies via $\partial_A E_{\rm
eff}=-\Ehat/A^2$)---but by two structural features. First, the \emph{orientation} of
the wind: Vaidya's is the ingoing \emph{radial} centre-wind $c_r=f-\Ehat^2$;
Thakurta--Kerr's is the \emph{azimuthal} frame-dragging shift. Second, and
decisively for the weight-two letter, the \emph{clock} that weights the source,
through the location of its logarithmic singularity: the advanced time $v$ diverges
at the horizon $r=2m$, giving Vaidya a $\dd\log$ letter \emph{there}; the $t,\tau$
clocks grow logarithmically at large $r$, giving Thakurta--Kerr a third-kind dipole
\emph{at infinity}. FLRW is the degenerate circle---$c=0$, isotropic $R$---where the
three branches coincide and the polylogs collapse to genus zero: no wind, no
dilogarithm. The organizing thread ``circle $\to$ radial-wind ellipse $\to$
angular-wind ellipse'' is thus literally the ladder of weight-two letters the
breathing switches on, on one and the same hyperelliptic backbone.

\paragraph{Domain of validity.} The statement is a \emph{first-order} adiabatic
theorem, and it rests on two \emph{exact} lemmas---the algebraic reduction
$\partial_\lambda F=N_\lambda/S^{3/2}$ with the unique rational $c_k$ from the
polynomial identity~\eqref{eq:reduction-identity} (nonsingular whenever $S$ has
distinct roots, i.e.\ away from the separatrix), and the integration-by-parts
assembly~\eqref{eq:psi-split} using the shuffle relation---so the closed form is
rigorous, not fitted, wherever those hypotheses hold. Two caveats bound it honestly.
First, it is $O(\dot\lambda)$: the second-order term needs the next reduction.
Second, and more importantly, the expansion \emph{fails at the turning points}
$S(r)=0$: there $\partial_\lambda F\sim N_\lambda/S^{3/2}$ diverges and the na\"ive
$\dot\lambda\,\delta\varphi$ is not uniformly valid, so the closed forms hold on the
smooth arcs between turning points and must be joined across them by the standard WKB
connection formulae (matched asymptotics). On the separatrix the \emph{mathematics}
is a controlled limit---the genus drops and the reduction degenerates to the elliptic
(Bloch--Wigner) case---but the \emph{physics} there is delicate for a separate reason:
the radial period diverges (the orbit asymptotes an unstable circular orbit), so the
adiabatic hypothesis itself---``$\lambda$ varies slowly compared with the orbital
period''---degrades, and the separatrix correction is properly read as an
\emph{instantaneous, local} first-order statement about the neighbourhood of the
separatrix rather than a long-time separatrix-following orbit. Moreover the degeneration
locus is not fixed under the evolution it responds to: in Vaidya
$J_{\rm deg}=m\,j(E)$ scales linearly with the mass (verified exactly proportional,
$J_{\rm deg}/m=7.0266$ at $E=1.4$, independent of $m$). This is the \emph{algebraic}
degeneration at $r_d<0$, not an external separatrix: the physical accretion$/$evaporation
capture--escape divide is the dynamical penetration threshold, not $J_{\rm deg}$. A
particle of conserved $J$ nonetheless \emph{crosses} the moving locus as $m(v)$ drifts;
that crossing---the same phenomenon that makes the Thakurta--Kerr off-shell separatrix
limit a Neishtadt separatrix-crossing problem---is a genuinely non-perturbative regime,
beyond first order, and the accretion--evaporation asymmetry of
Eq.~\eqref{eq:vaidya-asymmetry} is the first-order fingerprint that survives on either
side of it. Within these bounds, which
are the usual bounds of first-order WKB, this closed form is the \emph{on-shell
component} of the first-order response---the parameter drift along the frozen
family---and is exact as such; the \emph{complete} first-order term is the
extended-Hamiltonian expression Eq.~\eqref{eq:adiab-exact}, which adds the
off-shell (costate $p_\eta$) displacement of the momentum shell (a $\sim2\%$
correction here, Fig.~\ref{fig:phi-true-dynamic}). The separatrix
results in particular quantify the response \emph{at} the instantaneous marginal
configuration $J=\Jc(m)$, not a trajectory pinned to it. The complementary,
separatrix-\emph{following} response---imposing $J=\Jc(m(v))$ so that the orbit
tracks the moving marginal curve, with $\dd J/\dd v=(\Jc/m)\dot m$---closes by the
\emph{identical} reduction, the only change being the source $\partial_m F\to
\partial_m F+(\Jc/m)\,\partial_J F=(N_m+(\Jc/m)N_J)/S^{3/2}$ with $N_J$ again
polynomial; the same pole-adapted Weierstrass assembly then yields its closed form
(verified to $10^{-8}$), all coefficients again algebraic in $m,E$. The
advanced-time branch closes on the same separatrix by the identical scheme:
$\delta\varphi_v|_{\rm sep}=\tilde G(z)\,v-\int\tilde G\,v_z\,\dd z$ with
$v_z=\dd v/\dd z=E\,r^3/(r-r_d)+r\sqrt{Q_4}/(r-2m)$, the first term reproducing the
proper-time dilogarithms and the second---odd on the curve, with a simple pole at
the horizon two-torsion image $z=i\omega_{\rm im}$ (where $r=2m=e_3$)---furnishing the
horizon dilogarithm; the whole is explicit in $\sigma,\zeta,\wp$ and elliptic
dilogarithms with closed-form coefficients (verified to $10^{-8}$).


% ======================================================================
\section{Conclusions}
\label{sec:concl}
% ----------------------------------------------------------------------
\paragraph{Proved.} The non-stationary brachistochrone is a well-posed controlled-rail
problem: the invariant $-u\cdot W=\Ehat$ is legitimate via Pontryagin's principle, the
selector $W$ fixed canonically by the Killing--CKV--Kodama hierarchy, with the Kodama
energy conserved along the rail even when free fall does not. Existence, normality, and
the (non-autonomous) Hamilton--Jacobi verification criterion hold on the
timelike-selector domain. On the equatorial sector we obtained closed-form separatrices
in Weierstrass functions, continued through the ergosphere in Doran coordinates, and an
analytic cusp theorem for the $|J|<\Jc$ reflection; the on-shell first-order reduction
to length-two iterated Abelian integrals is a proved weight-one closed form. FLRW is the
degenerate base---freezing without branch splitting---while the mass breaks the $t/\tau$
symmetry through its radial gradient and the spin drives the rotational plunge inversion.

\paragraph{Conditional and numerical.} The plunge inversion resolves into a three-way
picture: rotation and the conformal factor invert the $t/\tau$ ordering under a
same-launch comparison, while in the ingoing model the negative-rate sign change is
absent at the frozen-clock level. The fixed-endpoint no-inversion is a \emph{theorem}
for $r_0\le R_\ast(E,a)$ [Eq.~\eqref{eq:Rstar}] and asymptotically as $r_0\to\infty$ in
the static case, machine-certified (interval arithmetic) at representative $r_0$, and a
numerically supported \emph{conjecture} in general. The complete first-order
(on-shell $+$ off-shell) correction is assembled on all four branches and verified
against the true non-autonomous flow to $O(\varepsilon^2)$; the weight-two letters
migrate from a third-kind dipole at infinity to a $\dd\log$ at the horizon and collapse
on the separatrix to an elliptic dilogarithm, evaluated by a convergent theta-nome
series with irreducibility and single-valued completion left conjectural.

\paragraph{Open problems.} The
exact Sultana--Dyer source
\cite{SultanaDyer2005}, the McVittie/Schwarzschild--de Sitter limits
\cite{McVittie1933,Kottler1918,Kaloper2010} with their photon surfaces
\cite{ClaudelViracVirbhadra2001}, escape maps, and freezing-surface dynamics.

\paragraph{Future directions.} Three threads are left open. First, a general theory of
\emph{moving turning points}: the first-order correction is uniform only on sub-arcs
with $H_{p_r}\neq0$, and a matched-asymptotic (WKB connection) treatment across the
turning point would render the orbit globally uniform. Second, a general theory of
\emph{adiabatic separatrices}: the instantaneous separatrix response closes in elliptic
dilogarithms, but a genuine crossing of a \emph{moving} separatrix is a Neishtadt inner
problem---jump/capture laws and matched asymptotics beyond first order. Third, in Vaidya
a genuine \emph{third clock} would require a distinguished observer family (e.g.\ Kodama
observers) rather than a coordinate analogy: whether a Kodama-observer proper time
defines a brachistochrone distinct from the advanced-null $v$ and the proper time $\tau$
remains open. By contrast, the conformal-time branch needs no separate
development---$t$ and $\eta$ are the \emph{same} arrival family---and the off-shell
source is now known in universal form on each deformation class
[Eqs.~\eqref{eq:SD-conformal},~\eqref{eq:SD-vaidya}].

\ack
The author thanks the University of Pavia for providing a stimulating academic
environment. The author received no external funding for this work. The AI
assistants GPT~5.5 (OpenAI) and Claude Opus~4.8 (Anthropic) were used to assist with
symbolic derivations, numerical implementation, and manuscript preparation; all
results were verified independently by the author, who takes full responsibility
for the content of this work. The closed-form and hyperelliptic computations relied
on open-source software, and the author gratefully acknowledges their developers:
\textsc{SageMath}~\cite{SageMath} (The Sage Developers) and the
\textsc{abelfunctions} library~\cite{abelfunctions} (C.~Swierczewski and
collaborators) for Riemann surfaces, Abel maps, and Riemann theta functions with
derivatives; and \textsc{SymPy}~\cite{SymPy}, \textsc{NumPy}~\cite{NumPy},
\textsc{SciPy}~\cite{SciPy}, and \textsc{Matplotlib}~\cite{Matplotlib} for the
symbolic reductions, numerical integration, and figures.

\section*{Data availability}
The complete numerical and symbolic codebase underlying this work is openly
available under the MIT license, archived at the persistent
DOI~\cite{Rosignoli2026Code}
\begin{center}
\url{https://doi.org/10.5281/zenodo.21707378}
\end{center}
(tagged release \texttt{v1.0}), and developed in the GitHub repository
\url{https://github.com/ImanetorX98/NonStationaryBrachistochrone}.
The repository includes the trajectory integrators, the Pontryagin Hamiltonian
constructions, the Poisson-bracket closure and quasi-invariant scans
(Sec.~\ref{sec:quasiconst}), the Weierstrass separatrix and Doran-continuation
scripts, the adiabatic-reduction and hyperelliptic-dilogarithm pipelines mapped in
Table~\ref{tab:script-map}, the fixed-endpoint BVP and colormap pipelines, and the
figure-generation code. In particular the self-contained script
\texttt{reproduce\_reductions.py} regenerates the reduction coefficients of
Table~\ref{tab:reduction-ck} from the polynomial
identity~\eqref{eq:reduction-identity} alone (one linear solve, residual zero),
letting a reader reproduce the algebraic heart of the closed forms in isolation. It
also retains exploratory scripts and negative results,
documenting which ansatz classes failed the held-out trajectory and phase-space
residual tests reported in Sec.~\ref{sec:quasiconst} and
\ref{app:validation}.



% ======================================================================
\appendix
\section{Validation and consistency checks}
\label{app:validation}
We collect the numerical checks that underpin the analytic results. The
minimum principle is verified directly by perturbing the computed
brachistochrones and confirming that both $T_t$ and $T_\tau$ are minimized
(Fig.~\ref{fig:flrw-var} for the FLRW branch degeneracy, and
Fig.~\ref{fig:verifmin} for the two distinct Vaidya branches). Stationary-limit
residuals confirm that the non-stationary
Hamiltonian flow reproduces the closed-form Kerr expressions as $m'\to0$ and
$A\to1$ (Fig.~\ref{fig:vaidya-kerra0}); the Kodama energy is conserved to
$\sim10^{-16}$ along the rail (Fig.~\ref{fig:kodama}). The equatorial closed
forms are validated end-to-end by the three-method superposition
(ODE flow, quadrature, evaluated Weierstrass) through the ergosphere
(Fig.~\ref{fig:sep3}), and the fixed-endpoint no-inversion theorem is
checked over the asymmetric-endpoint grids (Figs.~\ref{fig:asym-spin}--%
\ref{fig:asym-conf}).

\subsection{Explicit closed forms}
\label{app:explicit}

All expressions below are at \emph{constant (frozen) conformal factor} $A$:
stationary Kerr with $E\to E_{\rm eff}=\Ehat/A$, $J\to J/A$ (spin $a$ retained). They give
the instantaneous orbit shape on a Thakurta--Kerr slice; the time dependence
enters only through $E_{\rm eff}(\eta)$ (Sec.~\ref{sec:quasiconst}).

\paragraph{Curve conventions (fixed once for all letters).} Every letter
below---the second-kind $U_k=\int r^k/\sqrt S\,\dd r$, the weight-two
$W_{jk}=\int(r^jU_k-r^kU_j)/\sqrt S\,\dd r$, and the third-kind letters $L,\mathcal
D_k$---is an integral on the hyperelliptic curve $Y^2=\mathcal S(r)$ read with the
following fixed data. \emph{(i) Sheet:} $\sqrt S>0$ on the physical outer arc
$r>r_{\min}$, the sign flipping across each branch point, which fixes $Y$ globally by
analytic continuation. \emph{(ii) Cuts:} along the real segments between consecutive
real roots of $\mathcal S$, so the path stays on the physical sheet between turning
points. \emph{(iii) Base point and orientation:} all integrals share the lower limit
$r_0$ (the launch radius) and run with increasing $r$ toward arrival; reversing the
orientation flips the overall sign. \emph{(iv) Additive constants:} each
antiderivative is pinned at a regular reference---$U_k$ at $r_0$, the third-kind
letters at the regular point $z_\infty$ (where $\mathcal S$ vanishes as $1/r$), giving
the constants $C_{U_0},C_0$ used below---so that no spurious constant survives; these
constants multiply the common holomorphic differential and \emph{cancel} in every
assembled $\delta\varphi$. The same conventions apply verbatim to the separatrix
degeneration, with $\mathcal S\to Q_4$ and the Kleinian $\sigma,\zeta$ replaced by the
Weierstrass $\sigma,\zeta,\wp$.

\paragraph{Equatorial $\tau$-branch separatrix in Weierstrass functions.} On the
$\tau$-separatrix $J=\Jc$ the doranTau factorization $\Delta-\Jc^2 w=f(r^2+c^2)$,
$c=a/E$, makes the
turning quartic $Q_4=r\big((E^2-1)r+2M\big)(r^2+c^2)$, and the factor $r^2+c^2$
divides $Q_4$ so that only the horizon poles survive. Writing
\begin{gather}
    A=\frac{Ma^2}{2E^2},\qquad B=\frac{a^2(E^2-1)}{12E^2},\notag\\
    g_2=\frac{a^4(E^2-1)^2}{12E^4}-\frac{a^2M^2}{E^2},\notag\\
    g_3=\frac{a^4\big[36E^2M^2(1-E^2)-a^2(E^2-1)^3\big]}{216E^6},
    \label{eq:wp-invariants}
\end{gather}
the uniformizing variable is $z(r)=\wp^{-1}\!\big(A/r+B;g_2,g_3\big)
=\int_0^r\dd r'/\sqrt{Q_4}$, the Abel map of the turning quartic $Q_4$: it is the
elliptic argument in which the radial motion linearizes, so that $r$ is recovered
by $A/r+B=\wp(z)$ and the physical range $r\in[r_{\min},\infty)$ maps to a
segment of $z$. With $r_\pm=M\pm\sqrt{M^2-a^2}$ the horizons,
\begin{equation}
\begin{aligned}
    \alpha_\pm&=\frac{r_\pm\big((E^2-1)r_\pm+2M\big)}{r_\pm-r_\mp},\\
    \lambda_\pm&=\frac{\alpha_\pm}{\sqrt{Q_4(r_\pm)}},\qquad v_\pm=z(r_\pm),
\end{aligned}
\end{equation}
$\wp(v_\pm)=A/r_\pm+B$, and
$\Lambda_0=(E^2-1)-\alpha_+/r_+-\alpha_-/r_-+2\lambda_+\zeta(v_+)+2\lambda_-\zeta(v_-)$,
the angle is linear in the uniformizing variable $z\equiv z(r)$ plus two
Weierstrass-$\sigma$ ratios,
\begin{equation}
\begin{split}
    \varphi(r)=\frac{a}{E}\Big[&\Lambda_0\,z(r)
      +\lambda_+\ln\frac{\sigma(z-v_+)}{\sigma(z+v_+)}\\
      &+\lambda_-\ln\frac{\sigma(z-v_-)}{\sigma(z+v_-)}\Big]+\text{const}.
\end{split}
    \label{eq:sep-phi}
\end{equation}
Because $Q_4(2M)=4M^2(4E^2M^2+a^2)>0$, Eq.~\eqref{eq:sep-phi} is regular
\emph{through} the ergosphere $\re=2M$; the only singularities are the
third-kind logarithms at the horizons. The conformal case is
$E\to E_{\rm eff}=\Ehat/A$ in every parameter.

\paragraph{Doran continuation through the horizon.} The Boyer--Lindquist form
$\varphi(r)$ winds logarithmically at $\rp$ (the standard BL artifact); adding the
Doran shift $\dd\varphi_{\rm shift}=2Mar\,\dd r/(\Delta\sqrt{C_3})$ with the cubic
$C_3=2Mr(r^2+a^2)$ regularizes it. This is a \emph{second} elliptic curve, in
clean Weierstrass form $r(z)=(2/M)\wp(z;g_2,g_3)$ with $g_2=-M^2a^2$, $g_3=0$
(quasi-lemniscatic), giving
$\varphi_{\rm shift}=\tfrac{a}{r_+-r_-}\big(B_+-B_-\big)$,
$B_\pm=2\zeta(v_\pm)z+\ln[\sigma(z-v_\pm)/\sigma(z+v_\pm)]$. With
$\dd\varphi_{\rm shift}$ defined as above (a positive residue at $\rp$), the
log residues of $\varphi_{\rm BL}$ and $\varphi_{\rm shift}$ are \emph{equal}
(numerically $1.0324=1.0324$ at $a=0.9,E=1.2,J=3$), so they cancel under
\emph{subtraction}: the horizon-regular Doran azimuth is
\begin{equation}
    \varphi_D=\varphi_{\rm BL}-\varphi_{\rm shift},
    \label{eq:doran-phi}
\end{equation}
closed on two elliptic curves. (The single consistent sign is fixed by this
residue cancellation; the same convention is used throughout,
Eq.~\eqref{eq:doran-phi}.)

\paragraph{General $\tau$-branch $\varphi(r)$: genus-2 Kleinian form.} Off the
separatrix, for \emph{arbitrary} $J$, the $\tau$ integrand
$\dd\varphi/\dd r=\mathcal K_\tau(r)/\sqrt{\mathcal S(r)}$ is genus~2, with the
sextic
\begin{equation}
    \mathcal S(r)=r(r-2M)\big[(E^2-1)r+2M\big]
      \big[r\Delta-J^2\big((E^2-1)r+2M\big)\big]
    \label{eq:tau-sextic}
\end{equation}
(six generically simple branch points, hyperelliptic) and the rational numerator
$\mathcal K_\tau=Jr(r-2M)\big[(E^2-1)r+2M\big]/\Delta$. Partial fractions split it
into a first-kind (holomorphic) part and third-kind poles at the horizons
$r_\pm=M\pm\sqrt{M^2-a^2}$,
\begin{equation}
    \varphi(r)=c_1\!\int\!\frac{r\,\dd r}{\sqrt{\mathcal S}}
      +c_0\!\int\!\frac{\dd r}{\sqrt{\mathcal S}}
      +\sum_{\pm}\alpha_\pm\!\int\!\frac{\dd r}{(r-r_\pm)\sqrt{\mathcal S}},
    \label{eq:tau-genus2}
\end{equation}
with closed-form coefficients (in $M,a,E,J$)
\begin{equation}
    c_1=J(E^2-1),\quad c_0=2MJ,\quad
    \alpha_\pm=\mp\frac{Ja^2\big[M(E^2+1)\pm(E^2-1)\sqrt{M^2-a^2}\big]}
      {2\sqrt{M^2-a^2}}.
    \label{eq:tau-coeffs}
\end{equation}
The two holomorphic integrals are the Abel-map coordinates (arguments of the
Kleinian hyperelliptic $\sigma$), the horizon terms the third-kind
$\ln\sigma$-quotients---the exact genus-2 analogue of the $t$-branch $J_+^t$
[Eq.~\eqref{eq:t-genus2}]. We built $Y^2=\mathcal S$ in SageMath at $J=-0.9\Jc$
and computed its Riemann period matrix $\tau$; both Riemann relations hold
($\tau=\tau^\top$ to $7\times10^{-12}$, $\operatorname{Im}\tau\succ0$), and
Eq.~\eqref{eq:tau-genus2} reproduces the direct integral to $10^{-15}$. At
$|J|=\Jc$ the factor $f^2$ (from $\Delta-\Jc^2w=f(r^2+c^2)$) makes $r_e$ a
\emph{double} root of $\mathcal S$; the $(r-2M)$ in $\mathcal K_\tau$ cancels one
factor, dropping the radical to the quartic $Q_4$---genus~1, the Weierstrass
separatrix~\eqref{eq:sep-phi}. This double root is also why the retrograde
marginal $J=-\Jc$ reflects with a \emph{corner} (finite $\dd\varphi/\dd r$)
rather than a smooth grazing: only the prograde $+\Jc$, aided by co-rotation,
threads through.

\paragraph{$t$-branch closed form: generic integrand.} The coordinate-time
brachistochrone has $\dd\varphi/\dd r=\mathcal K(r)/\sqrt{\mathcal R(r)}$ with the
\emph{sextic} radical $\mathcal R=r\,Q_2(r)\,[(E^2-1)r+2M]$, genus~2, the quartic
being (no cubic term)
\begin{equation}
    Q_2(r)=E^2r^4+\big[E^2(a^2-J^2)+J^2\big]r^2
      +2M(J-a)\big[E^2(J-a)-2J\big]r+4M^2(J-a)^2,
    \label{eq:t-Q2}
\end{equation}
and a \emph{closed-form} numerator in the problem quantities,
\begin{equation}
    \mathcal K(r)=\frac{r\,[(E^2-1)r+2M]\,\big(J(r-2M)+2Ma\big)}{\Delta(r)},
    \qquad \Delta=r^2-2Mr+a^2,
    \label{eq:t-K}
\end{equation}
holding for \emph{any} $(M,a,E,J)$: the third-kind poles sit at the two horizons
$r_\pm$ ($\Delta=0$) and the polynomial part carries the first-/second-kind
pieces (there is \emph{no} pole at $r=0$, cancelled by the leading factor $r$).
By the conformal transfer $n^2_{\rm TK}=n^2_{\rm Kerr}(E_{\rm eff})$
(Sec.~\ref{sec:thakurta}), Eq.~\eqref{eq:t-K} is also the Thakurta--Kerr
$t$-integrand under $E\to E_{\rm eff}=\Ehat/A$, $J\to J/A$, so the conformal
factor is \emph{not} a separate variable of $\varphi(r)$.

\paragraph{$t$-branch separatrix in Weierstrass functions.} At the saddle-node
$J_c^-$ [Eq.~\eqref{eq:Jc-minus}] the radical acquires a
\emph{double root} at the marginal-capture radius $r_*$,
$\mathcal R=(r-r_*)^2 Q_4(r)$, and
$\dd\varphi/\dd r=\mathcal G(r)/\sqrt{Q_4}$ becomes a third-kind elliptic
differential on the quartic $Q_4$ (which retains the root $r=0$). With the same
uniformization $z(r)=\wp^{-1}(A/r+B;g_2,g_3)$ (invariants of $Q_4$) the closed
form is
\begin{equation}
    \varphi(r)=\Lambda_0\,z(r)+\sum_{k}\lambda_k\,
      \ln\frac{\sigma(z-v_k)}{\sigma(z+v_k)},
    \label{eq:t-sep-phi}
\end{equation}
now with third-kind poles only at $c_k\in\{\rp,\rm,r_*\}$---three real ones, with
\emph{no} pole at $r=0$ and no complex pair, unlike the $\tau$-separatrix. Here
$v_k=\wp^{-1}(A/c_k+B)$, $\lambda_k=-A\alpha_k/(c_k^2\wp'(v_k))$, and
$\Lambda_0=\Lambda_{\rm poly}-\sum_k\alpha_k/c_k+\sum_k 2\zeta(v_k)\lambda_k$,
with $\alpha_k$ the residues of $\mathcal G$. For $M=1$, $a=0.9$, $E=1.2$
($r_*=3.514$, $J_c^-=-8.05$) Eq.~\eqref{eq:t-sep-phi} reproduces the Hamiltonian
flow to $2\times10^{-14}$ (Fig.~\ref{fig:t-sep}). So both separatrices are
genus~1 by different mechanisms; only non-separatrix $t$-orbits need genus-2
Kleinian functions.

\begin{figure}[ht]
    \centering
    \includegraphics[width=\textwidth,height=0.68\textheight,keepaspectratio]{Immagini/fig_tbranch_separatrix.png}
    \caption{$t$-branch separatrix (the retrograde saddle-node $J_c^-$): the
    orbit spirals onto the marginal circle $r_*$. Top: Hamilton ODE vs the
    genus-1 closed form (the sextic radical has a double root, dropping to the
    quartic $Q_4$). Bottom: they agree to $\sim10^{-13}$; inset shows the double
    root of $R_6=(r-r_*)^2Q_4$.}
    \label{fig:t-sep}
\end{figure}

\paragraph{Genus-2 Kleinian form and ergosphere penetration.} Away from the
separatrices the sextic $\mathcal R$ keeps six \emph{simple} branch points
(here $\{-4.545,\,-1.722\pm1.793i,\,0,\,1.444,\,2.000\}$), so the curve
$y^2=\mathcal R(r)$ is a genuine genus-2 hyperelliptic Riemann surface and
$\varphi(r)$ closes only in \emph{Kleinian} (hyperelliptic $\sigma$) functions.
With the numerator~\eqref{eq:t-K} the partial-fraction split
$\mathcal K=c_1 r+c_0+\alpha_+/(r-\rp)+\alpha_-/(r-\rm)$ gives
\begin{equation}
    \varphi(r)=\underbrace{c_1\!\!\int\!\frac{r\,\dd r}{\sqrt{\mathcal R}}
      +c_0\!\!\int\!\frac{\dd r}{\sqrt{\mathcal R}}}_{\text{first kind}}
      +\underbrace{\sum_{\pm}\alpha_\pm\!\!\int\!
      \frac{\dd r}{(r-r_\pm)\sqrt{\mathcal R}}}_{\text{third kind}},
    \label{eq:t-genus2}
\end{equation}
the holomorphic pair being the Abel-map coordinates (arguments of the Kleinian
$\sigma$) and the horizon terms the third-kind $\ln\sigma$-quotients---the exact
genus-2 analogue of Eq.~\eqref{eq:t-sep-phi}. For $M=1$, $a=0.9$, $E=1.2$ we
built $y^2=\mathcal R$ in SageMath and computed its Riemann period matrix
$\tau$; both Riemann relations hold ($\tau=\tau^\top$ to $4\times10^{-12}$,
$\operatorname{Im}\tau\succ0$ with eigenvalues $\{0.61,1.08\}$), certifying the
homology basis. The single Abelian integral $\varphi(r)$ does not invert to a
one-$\theta$ quotient $r(\varphi)$ (that needs the 2D Jacobi inversion,
$r=\wp_{22}$ along the trajectory line in the Jacobian, à la
Hackmann--L\"ammerzahl); but Eq.~\eqref{eq:t-genus2} reproduces the Hamiltonian
flow to $2\times10^{-13}$. Physically these genus-2 orbits are exactly the ones
that \emph{penetrate the ergosphere}: the prograde threshold $J_+^t$
[Eq.~\eqref{eq:Jplus-t}] has its turning point \emph{on} $r_e=2M$ (simple root at
$r=2$), while for $J$ just inside the capture window $[J_c^-,J_+^t]$ the turning
point moves \emph{inside} the ergoregion and the closed form stays accurate
across and below $r_e$ (Fig.~\ref{fig:jpt-g2}). This penetration is generic in a
sharp sense that distinguishes the two branches (Table~\ref{tab:penetration}). The
$t$-branch obeys a \emph{dichotomy}: it penetrates over the whole finite interval
$J\in(J_c^-,J_+^t)=(-8.05,3.43)$---every $J$ inside crosses, every $J$ outside
reflects at a smooth periapsis $r_{\min}>r_e$, with \emph{no} cusp---because the
frame-dragging term auto-tunes the effective momentum
$\mathcal K(r_e)=fJ_t/E+2Ma/(rE)\to a/E$ to the critical value for \emph{all}
$J_t$, so a whole family (retrograde included) is carried across. The
$\tau$-branch obeys the opposite \emph{trichotomy}: its proper-time cost carries
the reality wall $\sqrt{wf}$ ($w=\Ehat^2-f$), imaginary for $f<0$ inside $r_e$,
which confines $\tau$-orbits to $r\ge r_e$ \emph{except} at the single fine-tuned
prograde value $J=+a/E$, where $\sqrt f$ cancels [Eq.~\eqref{eq:sep-phi}] and the
orbit crosses and spirals onto $r_+$. Off that value the $\tau$-orbit reflects:
smoothly for $|J|>a/E$ (periapsis outside $r_e$), but with a semicubic \emph{cusp}
at $r_e$ for $|J|<a/E$---the retrograde marginal $J=-a/E$ cusps too, counter-%
rotation being forbidden inside $r_e$. The $\tau$ penetrating set is thus a single
point, the $t$ one a finite interval: at coordinate time ergosphere penetration
is the rule, at proper time the fine-tuned exception. Only the retrograde
$t$-separatrix $J_c^-$ (the saddle-node) never reaches the ergosphere---it spirals
onto $r_*=3.514>r_e$.

\begin{figure}[ht]
    \centering
    \includegraphics[width=\textwidth,height=0.68\textheight,keepaspectratio]{Immagini/fig_jpt_genus2.png}
    \caption{Genus-2 prograde $t$-orbits and ergosphere penetration. Top: the
    threshold $J_+^t=3.43$ (grey) turns \emph{on} $r_e=2M$; a momentum just
    inside the capture window ($J=3.0$) penetrates the ergoregion (shaded),
    Hamilton ODE (black) and the genus-2 closed form~\eqref{eq:t-genus2}
    (dotted) coinciding across $r_e$. Bottom: the deviation stays
    $\sim2\times10^{-13}$ including \emph{inside} the ergosphere ($r<r_e$,
    shaded).}
    \label{fig:jpt-g2}
\end{figure}

The two $t$-window boundaries are themselves of different kinds
(Table~\ref{tab:penetration}): the prograde $J_+^t$ is a \emph{simple} root on
$r_e$ (the orbit reaches the ergosphere and crosses), whereas the retrograde
$J_c^-$ is a \emph{double} root at $r_*=3.514>r_e$---a logarithmic
$\varphi$-barrier onto which the orbit spirals with infinite winding, never
reaching $r_e$ even though the region $(r_e,r_*)$ is itself allowed
($\mathcal R>0$). At $J=J_c^-$ the closed form~\eqref{eq:t-genus2} accordingly
diverges as $r\to r_*$. Just inside the
window the mechanism is a \emph{saddle-node}: the double root splits into a
complex-conjugate pair ($r_*=3.514\to3.393\pm0.601i$ at $J=0.95\,J_c^-$), the
$\varphi$-barrier vanishes, and the retrograde orbit is captured---plunging
through $r_e$ to the horizon, with the genus-2 closed form tracking it to
$\sim10^{-9}$ (Fig.~\ref{fig:jcm-cap}).

\begin{figure}[ht]
    \centering
    \includegraphics[width=\textwidth,height=0.68\textheight,keepaspectratio]{Immagini/fig_jcm_capture.png}
    \caption{Retrograde capture just inside the window ($J=0.95\,J_c^-=-7.65$).
    Top: the orbit plunges from $r_0=6$ through the ergosphere (shaded) to the
    horizon $r_+$ (Hamilton ODE, black; genus-2 closed form, dotted), whereas
    the marginal $J_c^-$ (grey) only spirals onto $r_*$. Bottom: the saddle-node
    in the complex-$r$ plane---the double root of $\mathcal R$ at $r_*$ (square)
    splits into a complex-conjugate pair (circles) as $J$ moves off $J_c^-$,
    removing the barrier and opening the plunge.}
    \label{fig:jcm-cap}
\end{figure}

\begin{table}[ht]
    \centering
    \small
    \caption{Ergosphere behaviour of the two branches. The $\tau$-branch is a
    \emph{trichotomy} that penetrates only at the single prograde $J=+J_c$
    ($J_c=a/E$): the reality wall $\sqrt{wf}$ confines it to $r\ge r_e$, and off
    $+J_c$ it reflects---smoothly for $|J|>J_c$, with a semicubic cusp at $r_e$
    for $|J|<J_c$ (the retrograde marginal $-J_c$ cusps too, counter-rotation
    being forbidden inside $r_e$). The $t$-branch is a \emph{dichotomy} that
    penetrates the whole interval $(J_c^-,J_+^t)$, frame dragging auto-tuning the
    crossing; it never cusps. Values for $M=1$, $a=0.9$, $E=1.2$; orbit shapes in
    Figs.~\ref{fig:atlas-tau} and~\ref{fig:atlas-t}.}
    \label{tab:penetration}
    \begin{tabular}{lccc}
    \hline
    branch & penetrates & smooth bounce (outside $r_e$) & cusp bounce (at $r_e$) \\
    \hline
    $\tau$ & only $J=+J_c=+0.75$ & $|J|>J_c$
      & $|J|<J_c$ \ (and $J=-J_c$) \\
    $t$ & $(J_c^-,J_+^t)=(-8.05,3.43)$ & $J\notin(J_c^-,J_+^t)$ & --- (none) \\
    \hline
    \end{tabular}
\end{table}

\begin{figure}[ht]
    \centering
    \includegraphics[width=\textwidth,height=0.40\textheight,keepaspectratio]{Immagini/fig_atlas_tau.png}
    \caption{$\tau$-brachistochrones by regime (trichotomy). Each panel is one
    momentum relative to $\pm J_c$ (integrated by adaptive Hamiltonian flow):
    smooth bounce outside $r_e$ ($|J|>J_c$), the single penetrating $J=+J_c$
    (crosses and spirals onto $r_+$), cusp bounce at $r_e$ ($|J|<J_c$), the
    retrograde marginal $J=-J_c$ (also a cusp), and a smooth bounce for
    $J<-J_c$. Grey disc: ergoregion; dashed circles $r_e,r_+$.}
    \label{fig:atlas-tau}
\end{figure}

\begin{figure}[ht]
    \centering
    \includegraphics[width=\textwidth,height=0.40\textheight,keepaspectratio]{Immagini/fig_atlas_t.png}
    \caption{$t$-brachistochrones by regime (dichotomy). Relative to
    $J_c^\pm=(-8.05,3.43)$: smooth bounce ($J>J_c^+$), tangent contact at $r_e$
    ($J=J_c^+$), penetration ($J_c^-<J<J_c^+$), retrograde capture just inside
    the edge ($J=0.99\,J_c^-$, plunging to $r_+$), and a smooth bounce
    ($J<J_c^-$). Unlike $\tau$, a whole interval penetrates and no orbit cusps.}
    \label{fig:atlas-t}
\end{figure}

\paragraph{Stationary-limit conserved quantities.} At $A=1$ the branch flows admit
the closed first integrals used for validation,
\begin{equation}
    \mathcal K_\tau=J,\qquad
    \mathcal K_t=\frac{fJ+2Ma/r}{E},\qquad
    \mathcal K_\eta=\frac{fJ}{E},
    \label{eq:kerr-conserved}
\end{equation}
the equatorial Kerr doranTau/doranT constants.

\paragraph{Rotational inversion function.} The same-launch $t/\tau$ ordering
inverts where $\rho(r_*)^2=H_{t0}/H_{\tau0}$ with the explicit ratio of
centrifugal potentials
\begin{equation}
    \rho=\frac{A^2}{\Ehat}\Big(f+\frac{2Ma}{rJ}\Big),
    \label{eq:rho-explicit}
\end{equation}
$\rho<1$ giving $t$ deeper, $\rho>1$ giving $\tau$ deeper.

\paragraph{Conformal trichotomy.} With constant $A$ the separatrix momentum is
$\Jc=sA^2/\Ehat$, and above it the smooth periapsis is the outer root of
\begin{equation}
    \Delta(r)-\Big(\tfrac{J}{A}\Big)^2\Big[\big(\tfrac{\Ehat}{A}\big)^2-f(r)\Big]=0,
    \label{eq:periapsis-root}
\end{equation}
while $|J|<\Jc$ reflects at a cusp $r_{\min}=\re$ (as does the retrograde
marginal $-\Jc$, forced counter-rotation forbidding it inside $\re$), and only
the prograde $+\Jc$ crosses to $\rp$.

\paragraph{Closed toroidal averages.} The secular source $\langle S_1\rangle$ of
Eq.~\eqref{eq:S1} follows from the orbit-plane averages ($\cos\theta=\sin i\,
\sin\psi$, $k=1-J^2/L^2$)
\begin{equation}
    \Big\langle\cos^2\theta\Big\rangle=\tfrac{k}{2},\quad
    \Big\langle\tfrac{1}{\sin^2\theta}\Big\rangle=\tfrac{L}{|J|},\quad
    \big\langle p_\theta^2\big\rangle=L(L-|J|),
    \label{eq:averages}
\end{equation}
with $\langle\cos2\theta\,p_\theta^2\rangle=|J|(L-|J|)$.

\subsection{Reproducing the reductions}
\label{app:reproduce}
Every coefficient in the adiabatic corrections is fixed by a \emph{single}
polynomial identity that the reader can solve directly; nothing is fitted. The
depth-one reduction $\partial_\lambda\varphi_0=\mathcal{A}/\sqrt S+\sum_{k=0}^4 c_k
U_k$ [Eqs.~\eqref{eq:dEphi0},~\eqref{eq:dm-reduction}], differentiated in $r$ and
cleared of the radical using $\partial_\lambda F=N_\lambda/S^{3/2}$, is equivalent to
\begin{equation}
    2\,\mathcal{A}'(r)\,S(r)-\mathcal{A}(r)\,S'(r)
    +2\,S(r)\sum_{k=0}^4 c_k\,r^k \;=\; 2\,N_\lambda(r),
    \qquad \mathcal{A}=\sum_{i=0}^5 a_i r^i,
    \label{eq:reduction-identity}
\end{equation}
a polynomial identity of degree ten. Matching the eleven coefficients of
$r^0,\dots,r^{10}$ gives a nonsingular $11\times11$ linear system in the eleven
unknowns $\{a_0,\dots,a_5,c_0,\dots,c_4\}$; solving it (one call to a linear solver)
returns the $c_k$ \emph{rationally} in the metric parameters (here $S$ stands for the
branch radical: $S$ on the $\tau$-branch, $R_6=r\,Q_2\,DE$ on the $t$-branch). The
numerator is $N_\lambda=EJ\,r^4(r-2M)^2 DE$ for the $\tau$-branch energy reduction
(the horizon factor $\Delta$ cancels), $N_\lambda=R_6\,\partial_E\mathcal K
-\tfrac12\mathcal K\,\partial_E R_6$ for the $t$-branch (with $Q_2$, $\mathcal K$ of
Eqs.~\eqref{eq:t-Q2},~\eqref{eq:t-K}, the $\Delta$ in $\mathcal K$ again cancelling to
leave a degree-seven polynomial), and $N_m=S\,\partial_m K-\tfrac12K\,\partial_m S$
for the Vaidya mass reduction. The $t$-branch clock is entirely on-curve
[Eq.~\eqref{eq:clock-t}]: its second-kind vector $b$ feeds the $W_{kj}$ of $R_6$ and
its horizon part $R_\Delta/\Delta$ the third-kind $\ln\sigma$-quotients, with no
double cover. Three reference solutions to check against (polynomial residual $0$
exactly):
\begin{table}[htbp]
    \centering
    \caption{Reference reduction coefficients, obtained by solving
    Eq.~\eqref{eq:reduction-identity}, at $M{=}1$ and (Thakurta--Kerr rows)
    $a{=}0.9,E{=}1.2,J{=}2.5$, (Vaidya row) $a{=}0,E{=}1.4,J{=}2.5$. For
    Thakurta--Kerr the conformal drift acts through \emph{both} charges, so the source
    is the Euler operator and the combined weight is $C_k=\Ehat\,c_k^{E}+J\,c_k^{J}$:
    we list the energy ($\partial_E$) and angular ($\partial_J$) reductions separately;
    Vaidya has only $\partial_m$. Decimals are shown for readability; each $c_k$ is the
    exact rational solution of a single linear solve (residual $0$ exactly), regenerated
    by \texttt{reproduce\_reductions.py} (exact coefficients in the reproducibility
    package).}
    \label{tab:reduction-ck}
    \small
    \begin{tabular}{lrrrrr}
        \hline\hline
        reduction & $c_0$ & $c_1$ & $c_2$ & $c_3$ & $c_4$\\
        \hline
        TK $\tau$, $\partial_E$ & $-1.1575$ & $1.86406$ & $-0.57499$ & $0.02429$ & $0.07117$\\
        TK $\tau$, $\partial_J$ & $0.18027$ & $0.06996$ & $-0.29682$ & $0.03840$ & $0.03760$\\
        TK $t$, $\partial_E$ & $-0.39084$ & $1.72654$ & $-0.19058$ & $-0.14503$ & $-0.01676$\\
        TK $t$, $\partial_J$ & $4.33922$ & $4.03122$ & $-0.64349$ & $-2.63600$ & $-0.64444$\\
        Vaidya, $\partial_m$ & $-0.87153$ & $-0.35970$ & $1.03538$ & $0.21292$ & $-0.18939$\\
        \hline\hline
    \end{tabular}
\end{table}

Given the combined $C_k=\Ehat c_k^{E}+J c_k^{J}$, the weight-two
assembly~\eqref{eq:psi-split} is mechanical: the clock supplies the vector $b$ (on
the $\tau$-branch $\eta=U_3-2M\,U_2$, i.e.\ $b=(0,0,-2M,1,0)$), the algebraic
weights are $Q_{kj}=C_kb_j-C_jb_k$, the weight-two
letters are $W_{kj}=\int(U_k\dd U_j-U_j\dd U_k)$, and the elementary remainders
$\mathcal{I}_j=\int\mathcal{A}r^j/S\,\dd r$ integrate to logarithms. For Vaidya the
same steps with the advanced-time clock give the fully assembled
form~\eqref{eq:vaidya-full}; the horizon dilogarithm closes by one integration by
parts, Eq.~\eqref{eq:horizon-dilog}. Table~\ref{tab:script-map} maps each result to
the script in the repository~\cite{Rosignoli2026Code} that generates and verifies it;
the first entry, \texttt{pipeline\_completa\_deltaphi.py}, runs the entire chain in one
file---from $(M,a,E,J)$ through the $11\times11$ solve, the algebraic $c_k,b,Q_{kj}$,
and the building blocks $U_k,W_{kj},\mathcal{I}_j$, to the assembled $\delta\varphi$,
matching the direct double integral to $4\times10^{-15}$.
\begin{table}[htbp]
    \centering
    \caption{Result-to-script map with the achieved numerical residual, so each
    closed form can be regenerated and re-checked independently.}
    \label{tab:script-map}
    \small
    \begin{tabular}{@{}p{0.47\textwidth}p{0.34\textwidth}l@{}}
        \hline\hline
        result & script & residual\\
        \hline
        \emph{end-to-end example} $(M,a,E,J)\!\to\!\delta\varphi$
            & \texttt{pipeline\_completa\_deltaphi.py} & $4\times10^{-15}$\\
        TK $\tau$ reduction $c_k$, $\psi=\tfrac12\Ehat\sum Q_{kj}W_{kj}$
            & \texttt{kerr\_psi\_explicit\_verified.py} & $5\times10^{-14}$\\
        TK $t$-branch on-curve clock \& full assembly
            & \texttt{kerr\_tbranch\_psi\_assembly.py} & $9\times10^{-13}$\\
        quasi-periods $\eta$; $L$ (Fay third-kind)
            & \texttt{kerr\_quasiperiods\_bel.sage} & $10^{-12}$; $5\times10^{-8}$\\
        Vaidya reduction $c_k^m$
            & \texttt{vaidya\_dMF\_reduction.py} & $10^{-15}$\\
        Vaidya clock $v$; assembled $\delta\varphi$~\eqref{eq:vaidya-full}
            & \texttt{vaidya\_fully\_explicit.py} & $2\times10^{-13}$\\
        Vaidya $\tau$-branch assembly
            & \texttt{vaidya\_tau\_assembly.py} & $10^{-14}$\\
        horizon dilogarithm $\mathcal{D}_k=U_k\ln-\mathcal{G}_k$
            & \texttt{vaidya\_horizon\_dilog.py} & $10^{-12}$\\
        separatrix elliptic form $U_0(\sigma,\zeta)$; $\tau$
            & \texttt{vaidya\_ell\_dilog\_match.py} & $10^{-14}$\\
        accretion/evaporation split $\mathcal{A}_\infty,\mathcal{B}_{\rm hor}$
            & \texttt{vaidya\_asymmetry.py} & $10^{-15}$\\
        \hline
        \multicolumn{3}{@{}p{\textwidth}@{}}{\emph{First-order
            extended-Hamiltonian correction, Eq.~\eqref{eq:adiab-exact} (on-shell
            $+$ off-shell); paths under \texttt{ThakurtaMetric/},
            \texttt{VaidyaMetric/}:}}\\
        TK $t$ exact PT, $\delta\varphi$ vs true flow
            & \texttt{adiabatic\_first\_order\_exact.py} & slope $1.95$\\
        Vaidya $v$ exact PT $+$ Fig.~\ref{fig:vaidya-offshell}
            & \texttt{vaidya\_first\_order\_offshell.py} & slope $2.00$\\
        true-dynamics validation ($t,\tau$, 2 sets), Fig.~\ref{fig:phi-true-dynamic}
            & \texttt{fig\_phi\_validation\_corrected.py} & slopes $1.95$--$2.16$\\
        \hline\hline
    \end{tabular}
    \\[2pt]
    {\footnotesize SHA-256 (first 12 hex):
    \texttt{adiabatic\_first\_order\_exact.py} $=$\,\texttt{78ab24b5373c};
    \texttt{vaidya\_first\_order\_offshell.py} $=$\,\texttt{5dc9fe177f2f};
    \texttt{fig\_phi\_validation\_corrected.py} $=$\,\texttt{a086a0692e8c}.}
\end{table}

\subsection{Consistency residuals}
\label{app:residuals}
The closed forms and reductions above were checked symbolically and numerically.
The Weierstrass separatrix~\eqref{eq:sep-phi} reproduces the direct quadrature to
$|\Delta\varphi|<2.8\times10^{-10}$ through the ergosphere, with the $\wp$-identity
satisfied to $9\times10^{-28}$ and the antiderivative to $30$ digits; the Doran
shift reduces to $1.1\times10^{-16}$. The non-autonomous Randers
data~\eqref{eq:randers-data} match the Kerr null Fermat metric to
$\sim10^{-14}$. The equatorial Thakurta--Kerr flow reproduces the $A=1$ Kerr
constants~\eqref{eq:kerr-conserved} to ten digits at $r=4,6,9$ ($a=0.9$, $E=1.2$),
and the $a\to0$ limit the Thakurta--Schwarzschild flow to $10^{-12}$; the $3+1$
Hamiltonian~\eqref{eq:H3d} reduces to the equatorial one exactly, and conserves
$L^2$ to $6\times10^{-13}$ at $a=0$. The Kodama energy is held to
$|{-}u\cdot K-\Ehat|<7\times10^{-16}$ along the rail, and the static
inversion/trichotomy colormap agrees with the analytic
contours~\eqref{eq:periapsis-root} to a relative $7\times10^{-12}$.

\begin{figure}[tbp]
    \centering
    \includegraphics[width=\textwidth,height=0.68\textheight,keepaspectratio]{Immagini/fig_colormap_conforme_estremi_fissi.png}
    \caption{Conformal scan at fixed symmetric endpoints (via
    $E_{\rm eff}=\hat E/A$): $\Delta r>0$ throughout the sampled $(J,A)$ domain,
    consistent with the conjectured ordering and with the closed-form frozen-clock
    bound $n_t/n_\tau=E/f>1$.}
    \label{fig:cm-conf-fix}
\end{figure}


\begin{figure}[tbp]
    \centering
    \includegraphics[width=\textwidth,height=0.68\textheight,keepaspectratio]{Immagini/fig_colormap_spin_asimmetrico.png}
    \caption{Numerical robustness scan (not a theorem) of the fixed-endpoint
    no-inversion result under \emph{asymmetric} endpoints ($r_A=10,r_B=6$), spin
    scan (Kerr flow): $\Delta r>0$ over the sampled $(J,a)$ domain.}
    \label{fig:asym-spin}
\end{figure}

\begin{figure}[tbp]
    \centering
    \includegraphics[width=\textwidth,height=0.68\textheight,keepaspectratio]{Immagini/fig_colormap_conforme_asimmetrico.png}
    \caption{Same numerical robustness scan (not a theorem), conformal
    (via $E_{\rm eff}=\hat E/A$): $\Delta r>0$ over the sampled $(J,A)$ domain for
    asymmetric endpoints. The
    blank wedge is the region where the $\tau$-turning point rises above
    $r_B$ (no dipping orbit reaches $B$).}
    \label{fig:asym-conf}
\end{figure}

\clearpage

\begin{figure}[tbp]
    \centering
    \includegraphics[width=\textwidth,height=0.68\textheight,keepaspectratio]{Immagini/fig_vaidya_kerr_a0.png}
    \caption{Stationary-limit check ($m'=0$): numerical orbit vs closed-form
    Kerr $a=0$.}
    \label{fig:vaidya-kerra0}
\end{figure}


\begin{figure}[tbp]
    \centering
    \includegraphics[width=\textwidth,height=0.68\textheight,keepaspectratio]{Immagini/fig_verifica_minimo_brachi.png}
    \caption{Each branch minimizes its own travel time; $t$ and $\tau$ are
    distinct curves.}
    \label{fig:verifmin}
\end{figure}

\begin{figure}[tbp]
    \centering
    \includegraphics[width=\textwidth,height=0.68\textheight,keepaspectratio]{Immagini/fig_vaidya_bounce.png}
    \caption{Bounce through periapsis in the $v$-parametrization: the
    Euler--Lagrange flow is regular at $r'=0$ and integrates through the turning
    point, reaching $r_{\min}$ exactly where the $r$-parametrization diverges.}
    \label{fig:vaidya-bounce}
\end{figure}

\begin{figure}[tbp]
    \centering
    \includegraphics[width=\textwidth,height=0.68\textheight,keepaspectratio]{Immagini/fig_vaidya_timing.png}
    \caption{Timing window: $r_{\min}$ depends on the arrival epoch $v_1$. The
    orbit--horizon race has opposite winners in absolute ($r_{\min}$) and relative
    ($r_{\min}/2m$) terms for accretion vs evaporation.}
    \label{fig:vaidya-timing}
\end{figure}



\section{Explicit adiabatic forms on the genus-degeneration loci}
\label{app:separatrix}
This appendix collects in one place the fully explicit first-order adiabatic
brachistochrone on the loci where the branch sextic acquires a double root and drops to
genus one. For \emph{spherical} Vaidya the relevant value $J=\Jc$ is the \emph{algebraic}
genus-degeneration $J_{\rm deg}$---the double root sits at $r_d<0$
(Sec.~\ref{sec:trichotomy}), so it is \emph{not} a physical, accessible separatrix,
though the genus-two$\,\to\,$genus-one reduction is valid there all the same; we give both
the proper-time ($\tau$) and advanced-time ($v$) branches with every coefficient in closed
form. We then treat the rotating Thakurta--Kerr $\tau$ and $t$ branches, whose degeneration
values $\Jc^\pm$ include the \emph{physical} external separatrix (accessible double root
$r_d>2M$) alongside further algebraic $J_{\rm deg}$ roots at $r_d<0$. All objects live on
the elliptic curve $E:\,w^2=Q_4(r)$ into which the branch sextic collapses. The adiabatic expansion is
$\varphi(r)=\varphi_0(r)+\dot\lambda\,\delta\varphi(r)+O(\dot\lambda^2)$
($\lambda=m$ for Vaidya, $\lambda=A$ via $E_{\rm eff}=\Ehat/A$ for Thakurta--Kerr).

\paragraph{Curve and special functions.}
$\Jc=5\sqrt{3011/3072+581\sqrt{249}/9216}$ (at $E=7/5,\,m=1$; in general $\Jc=m\,j(E)$),
$S=(r-r_d)^2Q_4(r)$ with double root $r_d$ and quartic $Q_4$ (roots
$e_1<e_2<e_3=2m<e_4$, leading coefficient $a_4$). The Abel coordinate is
$z(r)=\int_{r_0}^r\dd r'/\sqrt{Q_4}$; $\sigma,\zeta,\wp$ are the Weierstrass functions
of $E$ (periods $2\omega_1,2i\omega_{\rm im}$, built from $\theta_1$). Marked points:
$\pm z_d$ (images of $r_d$), $\pm z_\infty$ (of $r=\infty$), the two-torsion turning
image $z=0$ ($r=e_4$) and horizon image $z=i\omega_{\rm im}$ ($r=e_3=2m$).

\paragraph{Zeroth order (frozen orbit), verified to $10^{-12}$.}
\begin{equation}
\varphi_0(r)=\Jc\big[(E^2-1)z+((E^2-1)r_d+2m)\,U_0(z)\big],\;\;
U_0=\rho[\ln\sigma(z{-}z_d)-\ln\sigma(z{+}z_d)]+C_{U_0}z,
\end{equation}
$\rho=1/\sqrt{Q_4(r_d)}$, $U_k=\int_{r_0}^r r^k/\sqrt S\,\dd r$ (explicit in
$\sigma,\zeta$ by Hermite recursion).

\paragraph{First order: source and antiderivative.}
$\partial_m F=N_m/S^{3/2}$, $N_m=S\,\partial_mK-\tfrac12K\,\partial_mS$, $K=\Jc\,DE$,
$DE=(E^2-1)r+2m$; on the curve $\partial_m F\,\dd r=R(z)\,\dd z$ with
$R=N_m(r(z))/[(r(z)-r_d)^3Q_4(r(z))]$ and $r(z)=c_r-a_4^{-1/2}[\zeta(z-z_\infty)-
\zeta(z+z_\infty)]$. Its antiderivative is Eq.~\eqref{eq:vaidya-sep-G},
$\tilde G(z)=C_0 z+\sum_{a}[b_1^a\ln\sigma(z-a)-b_2^a\zeta(z-a)-\tfrac12 b_3^a
\wp(z-a)]$, $a\in\{\pm z_d,0,i\omega_{\rm im}\}$, with closed-form residues: at the
triple pole $z_d$, writing $s=\sqrt{Q_4(r_d)}$, $F=N_m/Q_4$, $a_1=Q_4'(r_d)/4s$,
$a_2=Q_4''(r_d)/12$, $h_0=F(r_d)$, $h_1=F'(r_d)s$,
$h_2=\tfrac12F''(r_d)s^2+\tfrac14F'(r_d)Q_4'(r_d)$,
\begin{equation}
b_3^{z_d}=\frac{h_0}{s^3},\quad b_2^{z_d}=\frac{h_1-3a_1h_0}{s^3},\quad
b_1^{z_d}=\frac{h_2-3a_1h_1+(6a_1^2-3a_2)h_0}{s^3},\quad
b_2^{e_i}=\frac{4N_m(e_i)}{(e_i-r_d)^3Q_4'(e_i)^2},
\end{equation}
($e_i=e_4,e_3$ the two-torsion poles) and the constant closing at $z_\infty$, where
$R$ vanishes: $C_0=-\sum_a[b_1^a\zeta(z_\infty-a)+b_2^a\wp(z_\infty-a)-\tfrac12
b_3^a\wp'(z_\infty-a)]$. This example is Vaidya (source $\partial_m F$, hence
$N_m$); for the Thakurta--Kerr separatrices the identical Weierstrass assembly runs
with the combined conformal numerator $N=\Ehat N_E+J N_J$ (Sec.~\ref{sec:thakurta},
referee point on $J_{\rm eff}=J/A$), giving $b_i=\Ehat b_i^{E}+J b_i^{J}$; the
angular part shifts $b_1$ by $\sim\!30\%$ and is tabulated in
\texttt{ThakurtaMetric/conformal\_source\_Jterm.py}.

\paragraph{First order, proper-time branch (verified to $10^{-9}$).}
Clock $\eta=U_3-2mU_2$; block form
\begin{equation}
\delta\varphi_\tau=\tfrac12\,\tilde G(z)\,\eta(r)
+\tfrac12\sum_{i<j}(c_i d_j-c_j d_i)\,A[e_i,e_j],
\end{equation}
$A[e_i,e_j]=\int(e_i\,\mathrm{Pe}_j-e_j\,\mathrm{Pe}_i)\dd z$ the antisymmetric
elliptic dilogarithms (Brown--Levin $\tilde\Gamma$), over the eight even blocks
$\{1,Z_{z_d},P_{z_d},P'_{z_d},\wp(z),\wp(z{-}i\omega_{\rm im}),P_{z_\infty},
Z_{z_\infty}\}$ with $Z_a=\zeta(z{-}a)-\zeta(z{+}a)$,
$P_a=\wp(z{-}a)+\wp(z{+}a)$, $P'_a=\wp'(z{-}a)-\wp'(z{+}a)$. Coefficients: $c_i$ (the
$\tilde G$ residues), $d_j$ (clock residues $e_1^{z_d}=(r_d^3-2r_d^2)/s$,
$e_2^{z_\infty}=1/a_4$, $e_1^{z_\infty}=-(2B+r_d-2)/\sqrt{a_4}$) and the constant
$C_e=(e_4^3-2e_4^2)/(e_4-r_d)+2e_1^{z_d}\zeta(z_d)-2e_2^{z_\infty}\wp(z_\infty)
+2e_1^{z_\infty}\zeta(z_\infty)$---all closed form.

\paragraph{First order, advanced-time branch (verified to $10^{-8}$).}
Clock $v=E\,U_3+r+2m\ln(r-2m)$;
$\delta\varphi_v=\tilde G(z)\,v-\int_{z_0}^z\tilde G\,v_z\,\dd z$ with
$v_z=E\,r^3/(r-r_d)+r\sqrt{Q_4}/(r-2m)$. The first term reproduces the proper-time
dilogarithms; the second is odd on $E$ with a simple pole at $z=i\omega_{\rm im}$ and
supplies the horizon dilogarithm.

\subsection{Adiabatic derivative of the degeneration family}
\label{app:tracking}
As the slow parameter drifts, the genus-degeneration value itself moves,
$J=\Jc(\lambda)$. This is \emph{not} the evolution of a single torque-free extremal:
the axial costate $J$ is conserved along every extremal
(Sec.~\ref{sec:principle}), so one trajectory cannot follow the moving locus by
changing $J$ unless an external azimuthal torque---a modified control law with
$\dd J/\dd\lambda\neq0$---is supplied. The object below is instead the total
$\lambda$-derivative of the \emph{one-parameter family of marginal (degeneration)
solutions} along $J=\Jc(\lambda)$: the adiabatic derivative of the family, which
packages how the closed form deforms as the locus drifts,
\begin{equation}
N_{\rm tot}=N+\frac{\dd\Jc}{\dd\lambda}\,N_J,\qquad
N_J=S\,\partial_JK-\tfrac12K\,\partial_JS,
\label{eq:Ntot-track}
\end{equation}
with the block form and residue formulae above holding verbatim under $N\to N_{\rm tot}$;
$\varphi_0$, the clocks and the dilogarithm structure are unchanged. The three branches
differ only in $\dd\Jc/\dd\lambda$ and $N_J$.

\paragraph{Vaidya ($\lambda=m$).}
The linear scaling $\Jc=m\,j(E)$ gives $\dd\Jc/\dd m=\Jc/m$, with
$N_J=S\,DE+\Jc^2\,r(r-2m)DE^3$ (verified to $10^{-8}$).

\paragraph{Thakurta--Kerr $\tau$ and $t$ branches.}
The rotating (conformal) case follows by the transfer $E\to E_{\rm eff}=\Ehat/A$,
$J\to J/A$ (spin $a$ retained); the control parameter is now $A$ (equivalently $E_{\rm eff}$), and
the separatrix $\Jc(E_{\rm eff})$ drifts as $A$ runs. Two curves carry the two clocks.
\emph{Proper-time ($\tau$) branch:} curve $S=r(r-2M)DE\,(r\Delta-J^2DE)$, source
$N_\tau=\partial_E(K/\sqrt S)S^{3/2}=E\,J\,r^4(r-2M)^2DE$, clock
$\eta=U_3-2MU_2$; the block form and residue formulae are those above with
$N_m\to N_\tau$. Tracking uses $N_\tau\to N_\tau+(\dd\Jc/\dd E)N_J$ with
$N_J=r^3(r-2M)^2DE^2$ ($\Delta$ cancelling identically) and
$\dd\Jc/\dd E=-E\,\Jc\,r_d/DE(r_d)=-152.96$ (at $M=1,\,a=9/10,\,E_{\rm eff}=6/5$;
verified against finite differences, block $=$ direct to $10^{-7}$).
\emph{Coordinate-time ($t$) branch:} curve $R_6=r\,Q_2\,DE$ with $Q_2$
\emph{quadratic} in $J$ (it carries $J^2$ and $(J-a)^2$ terms---not linear), so the
discriminant $\mathrm{Res}_r(Q_2,\partial_r Q_2)=0$ has \emph{six} real roots that
must be classified by the sign and location of the associated double root $r_d$
(script \texttt{separatrix\_classification.py}). Only an \emph{external} double
root $r_d>r_+$, at which two real turning points \emph{merge}, is a physical
separatrix (capture$/$scatter boundary): this is the retrograde $\Jc^-=-8.054$
($r_d=+3.514>2M$) of Eq.~\eqref{eq:Jc-minus}, consistent with the body, plus a
prograde \emph{ergospheric} branch $\Jc=+2.936$ at $r_+<r_d=1.512<2M$. The large
roots $\Jc=+19.089,\,-18.671$ instead sit at $r_d=-6.62,\,-6.59<0$: these are
\emph{algebraic genus-degenerations} $J_{\rm deg}$ (the sextic drops to genus one
and the Weierstrass reduction still closes) but \emph{not} physical separatrices---
no external turning-point merger, no divergent radial period, and they must not be
conflated with the dynamical penetration threshold $J_c(v_0)$. The clock is the
mixed second/third kind $\dd t/\dd r=\rho_t/\sqrt{R_6}$, $\rho_t=P_3+R_\Delta/\Delta$
[Eq.~\eqref{eq:clock-t}], adding a third-kind horizon letter. On the physical branch
$\Jc^-$ the corrected conformal source (Sec.~\ref{sec:thakurta}, with the
$J\,\partial_J F$ term) gives $b=(b_1,b_2,b_3)=(-0.685,-0.074,-0.0035)$ and the
\emph{tame} separatrix drift
\begin{equation}
\frac{\dd\Jc^-}{\dd E}=-\frac{\partial_E Q_2}{\partial_J Q_2}\bigg|_{(r_d,\Jc^-)}=+2.44
\end{equation}
(the earlier $\mp148$--$150$ magnitudes belonged to the unphysical $J_{\rm deg}$
roots at $r_d<0$). All tracking coefficients are closed-form residues, symbolic in
$(M,a,E,r_d,\Jc)$; the triple-pole cancellation below is an algebraic identity at
\emph{any} double root, hence holds on the physical branch as well.

\paragraph{Tracking cancels the triple pole.}
A clean structural consequence closes the logic: at the double root the tracked
source \emph{vanishes},
\begin{equation}
N_{\rm tot}(r_d)=-\tfrac12K(r_d)\Big[\partial_\lambda S
+\tfrac{\dd\Jc}{\dd\lambda}\,\partial_JS\Big](r_d)
=\tfrac12K(r_d)\,S'(r_d)\,\frac{\dd r_d}{\dd\lambda}=0,
\label{eq:track-cancel}
\end{equation}
because $S'(r_d)=0$ (the separatrix double-root condition). Hence the triple-pole
residue drops, $b_3^{\rm track}=0$: following the moving separatrix softens the
strongest singularity, precisely where the double root travels. This is verified
symbolically in Mathematica, $N_{\rm tot}(r_d)\equiv0\pmod{S,\,S'}$, for both Vaidya
and Thakurta--Kerr.

\paragraph{Fully symbolic coefficients, off-separatrix $q$-series, and independent checks.}
Off the separatrix (generic $J$, genus two) every block coefficient is a
\emph{rational function of $(M,a,E,J)$}: the source coefficients $c_k$,
$Q_{kj}=c_kb_j-c_jb_k$, and the canonical pole coefficients $g_i$ (Taylor
coefficients of $q_6^{-1/2}$ at the odd-model infinity), with
$g_0=1/\sqrt{E^2-1}$ and $g_1=(2E^2-3)/(E^2-1)^{3/2}$ set by the leading behaviour
($E$-only), while $g_2$ and the $c_k$ carry the full $(a,J)$ dependence. The
algebraic remainder is elementary in closed form,
$T_{\rm alg}=P(r)+\sum_i\mathrm{res}_i\ln(r-r_i)$ over the six roots of $S$
(three residues vanish, at $r=0,\,2M,\,2M/(1-E^2)$). The genus-two dilogarithm is a
convergent \emph{nome ($q$-)series}: after Siegel reduction $\tau'=-\tau^{-1}$ (all
nomes ${<}1$), $\Lambda=\int Q\,\dd\mathcal A+\int L\,\dd\mathcal A$ with $Q$
elementary and $L=\log[\theta[\delta](w-e_+)/\theta[\delta](w-e_-)]$ a genus-two
Kronecker--Eisenstein series that converges geometrically (order $N{=}4$ already
matches to $10^{-13}$) and does \emph{not} reduce to classical dilogarithms (there is
no genus-two Jacobi product formula). All these symbolic statements are
\emph{independently reproduced in Mathematica} (native $\wp$, Gröbner reduction,
Laurent extraction): the separatrix residues to $10^{-16}$--$10^{-38}$, the generic-$J$
reductions to an exact polynomial identity, and the tracking
theorem~\eqref{eq:track-cancel}.

\paragraph{Transcendence.}
None of these is elementary. The $U_k$ are third-kind elliptic integrals, $\tilde G$
lives in $\sigma,\zeta,\wp$, and the weight-two content is the elliptic dilogarithm,
conjecturally irreducible to weight one (theta-divisor obstruction $\Theta=W_{g-1}$). The
advanced-time horizon logarithm is elementary only as a $\dd\log$ letter; integrated
against the elliptic $\partial_m F$ it yields the horizon elliptic dilogarithm, of the
same transcendental weight as the proper-time case---the horizon merely relocates the
dilogarithm's anchor from infinity to the two-torsion point $r=2m$. The sole
numerical ingredient anywhere is the optional compression of the weight-two content to
an apparent minimal five-dimensional dilogarithm basis (a numerically stable rank five,
singular values falling to $10^{-79}$---strong evidence, \emph{not} a proof, and
curve-specific), which the block forms above do not require.

\clearpage
\section{Complete first-order extended-Hamiltonian correction}
\label{app:firstorder}
This appendix derives Eq.~\eqref{eq:adiab-exact}---the complete first-order (on-shell
\emph{plus} off-shell) adiabatic correction to the orbit shape---from the extended
Hamiltonian, and specializes it to the $t$-, $\tau$- and Vaidya branches so that a
reader can reconstruct it without the accompanying code.

\paragraph{Notation and the slow datum.} $\Ehat$ and $J\equiv J_{\rm TK}=p_\varphi$ are
the \emph{conserved} control charges (Kodama energy and axial costate); the effective
unit-mass Kerr charges on the frozen slice are $E_{\rm eff}=\Ehat/A$ and
$J_{\rm eff}=J/A$, both of conformal weight $A^{-1}$. The slow datum is $s$ ($s=\eta$ for
Thakurta--Kerr, $s=v$ for Vaidya); we write $A(s)=A_0(1+\varepsilon s+\cdots)$ with
$A_0\equiv1$ and the \emph{dimensionless} rate $\varepsilon=M\,A_{,\eta}/A$ (for Vaidya
$\varepsilon=M\dot m/m$). The operator generating the drift of the frozen family is the
Euler operator
\begin{equation}
    \Theta\equiv E_{\rm eff}\,\partial_{E_{\rm eff}}+J_{\rm eff}\,\partial_{J_{\rm eff}}
    \quad(\text{Thakurta--Kerr}),\qquad
    \Theta\equiv m\,\partial_m\quad(\text{Vaidya}),
    \label{eq:theta-op}
\end{equation}
since $\dd E_{\rm eff}/\dd s=-\varepsilon E_{\rm eff}$, $\dd J_{\rm eff}/\dd s=-\varepsilon
J_{\rm eff}$ at leading order (Vaidya: $\dd m/\dd s=+\varepsilon m$).

\paragraph{Frozen data (per branch).} At fixed $s$ the branch Hamiltonian
$H_0(r,p_r;E_{\rm eff},J_{\rm eff})=0$ fixes: the \emph{ingoing} frozen momentum
$p_{r0}(r)$, the root of $H_0=0$ with $H_{0,p_r}<0$ (the outgoing arc is its mirror
across the turning point); the shape integrand $G_0(r)=\dd\varphi/\dd r=H_{0,J}/H_{0,p_r}$;
and the clock density $\dd\lambda/\dd r=1/H_{0,p_r}$, with $\lambda$ the elapsed clock
oriented \emph{inward from the free-arrival endpoint}.

\paragraph{The conformal rescaling: canonical derivation of the dilation term.} The
Thakurta--Kerr Hamiltonians factorize (Eqs.~\eqref{eq:Heta}--\eqref{eq:Htau-tk}) as
$H_{\rm br}(r,p_r,p_\varphi;A)=A\,\bar H_{\rm br}(r,P_r,J_{\rm eff};E_{\rm eff})$ with the
\emph{normalized} radial momentum $P_r=p_r/A$ (and $J_{\rm eff}=p_\varphi/A$,
$E_{\rm eff}=\Ehat/A$). We derive the flow in the \emph{original} canonical variables
$(r,p_r)$ and only then pass to $P_r$. Adjoining the slow clock $s$ with conjugate
$p_s$, the extended Hamiltonian $\mathcal H=p_s+H_{\rm br}(r,p_r;A(s))=0$ generates the
\emph{canonical} equations $\dd r/\dd s=\partial_{p_r}H_{\rm br}$,
$\dd p_r/\dd s=-\partial_r H_{\rm br}$. The rescaling $P_r=p_r/A$ is \emph{not} itself
canonical---at fixed $A$, $\dd r\wedge\dd p_r=A\,\dd r\wedge\dd P_r$, so it is a
time-dependent (conformally symplectic) \emph{normalization}, not a canonical
transformation. Its $s$-derivative along the canonical flow follows by the chain rule
(with $\alpha\equiv A_{,s}/A=\varepsilon+O(\varepsilon^2)$ and the identity
$\partial_r H_{\rm br}|_{p_r}=A\,\bar H_r|_{P_r}$):
\begin{equation}
    \frac{\dd r}{\dd s}=\bar H_{P_r},\quad
    \frac{\dd P_r}{\dd s}=\frac1A\frac{\dd p_r}{\dd s}-\frac{A_{,s}}{A}P_r
      =-\bar H_r-\alpha P_r,\quad
    \frac{\dd E_{\rm eff}}{\dd s}=-\alpha E_{\rm eff},\quad
    \frac{\dd J_{\rm eff}}{\dd s}=-\alpha J_{\rm eff},
    \label{eq:normflow}
\end{equation}
an identity verified symbolically (\texttt{ThakurtaMetric/canonical\_dilation\_check.py}).
The $-\alpha P_r$ term is thus a \emph{consequence} of the canonical flow in $(r,p_r)$
together with the normalization---not an added piece, and not the generator of a
canonical transformation; dropping it would misrepresent $\dd(p_r/A)/\dd s$. The
normalized system~\eqref{eq:normflow} is non-canonical (conformally symplectic) but
\emph{equivalent} to the canonical derivation. Consequently the shell-departure rate is
the \emph{full} Euler operator
\begin{equation}
    \frac{\dd\bar H}{\dd s}=-\alpha\,\mathcal D\bar H,\qquad
    \mathcal D\equiv\Theta+P_r\partial_{P_r}
    =E_{\rm eff}\partial_{E_{\rm eff}}+J_{\rm eff}\partial_{J_{\rm eff}}+P_r\partial_{P_r}.
    \label{eq:Dop}
\end{equation}

\paragraph{Extended system and first order.} Adjoin $s$ with conjugate $p_s$:
$\mathcal H=p_s+\bar H(r,P_r;E_{\rm eff}(s),J_{\rm eff}(s))=0$. Free-arrival transversality
sets $\bar H=0$ (hence $p_s=0$) at $\lambda=0$. Along the flow
$\dd p_s/\dd\lambda=\varepsilon\,\mathcal D\bar H$, so the \emph{terminally-anchored}
source is
\begin{equation}
    p_s^{(1)}(\lambda)\equiv S_{\mathcal D}(\lambda)
    =\int_0^{\lambda}\mathcal D\bar H\,\dd\lambda'
    =\int_0^{\lambda}\big(\Theta\bar H+P_r\bar H_{P_r}\big)\,\dd\lambda',
    \label{eq:pS-app}
\end{equation}
and, linearizing $\bar H$ at fixed $r$, the momentum-shell shift
\begin{equation}
    \delta P_r=\frac{\lambda\,\Theta\bar H-S_{\mathcal D}}{\bar H_{P_r}},
    \label{eq:dpr-app}
\end{equation}
the axial costate remaining conserved ($\delta J=0$; it is $J_{\rm eff}=J/A$ that
breathes). The shape varies as $\dd\varphi/\dd r=G_0+\varepsilon[\partial_{P_r}G_0\,\delta
P_r-\lambda\,\Theta G_0]+O(\varepsilon^2)$; integrating in $r$ gives
Eq.~\eqref{eq:adiab-exact} \emph{with $S\to S_{\mathcal D}$}. The earlier $S=\int\Theta\bar
H\,\dd\lambda$ (omitting $P_r\bar H_{P_r}$) closes only against a surrogate flow lacking
$-\alpha P_r$; against the true canonical flow it leaves an $O(\varepsilon)$ error. The
reduction is non-uniform at the turning point ($\bar H_{P_r}\to0$): it holds on compact
sub-arcs, and all displayed residuals are taken there.

\paragraph{The dilation letter in closed form (so the reader can reconstruct it).}
Since $\dd\lambda=\dd r/\bar H_{P_r}$, the extra source piece is the radial action,
\begin{equation}
    \int_0^{\lambda}P_r\bar H_{P_r}\,\dd\lambda'=\int p_{r0}\,\dd r,
    \label{eq:dilation-action}
\end{equation}
and this is elementary to reduce. On the frozen $t$-branch shell the radial momentum is
\begin{equation}
    p_{r0}=\frac{\sqrt{S}}{\Delta\,[(E^2-1)r+2M]},\qquad
    \Delta=r^2-2Mr+a^2,\qquad S=r\,[(E^2-1)r+2M]\,Q_2(r),
    \label{eq:pr0-tbranch}
\end{equation}
with $E\equiv E_{\rm eff}$, $J\equiv J_{\rm eff}$, $Q_2$ the branch quartic (Eq.~\eqref{eq:t-K})
and $S$ the on-shell sextic. Polynomial division of $S$ by $\Delta\,[(E^2-1)r+2M]$ gives,
in the standard basis $U_k=\int r^k/\sqrt S\,\dd r$,
\begin{equation}
    \int p_{r0}\,\dd r=\sum_{k=0}^{3}a_k\,U_k+(\text{third kind at }\Delta=0),
    \label{eq:action-reduction}
\end{equation}
\begin{equation}
    a_3=E^2,\quad a_2=2E^2M,\quad a_1=4E^2M^2-(E^2-1)J^2,\quad
    a_0=2M\big[4E^2M^2-J^2+2(1-E^2)Ja\big],
    \label{eq:ak}
\end{equation}
with $[a_k]=\text{length}^{3-k}$. The kernel is likewise algebraic on the shell,
$\partial_{P_r}G_0/\bar H_{P_r}=A(r)/\sqrt S$ with $A$ rational (the inner radical
$\sqrt{\Delta v/\bar P}$ cancels on shell), so the complete off-shell correction is a
length-two iterated integral on the single spectral curve $\sqrt S$, with the symbolic
coefficients~\eqref{eq:ak}. Every step is verified to machine precision against the
\emph{true} canonical flow~\eqref{eq:normflow}
(\texttt{ThakurtaMetric/adiabatic\_offshell\_reduction.py},
\texttt{adiabatic\_first\_order\_exact.py}, \texttt{tests/test\_adiabatic\_noreg.py}).

\paragraph{Special-function class.} The building blocks form a strict transcendence
hierarchy. The second-kind integrals $U_0,\dots,U_3=\int r^k/\sqrt S\,\dd r$ are
\emph{ordinary Abelian integrals} (weight one), reducible to the classical Kleinian
$\zeta,\sigma$ basis (Baker; \ref{app:explicit}), the same class as $\varphi_0$. The
letters $W_{jk}=\int(r^jU_k-r^kU_j)/\sqrt S\,\dd r$ are \emph{length-two iterated} Abelian
integrals with algebraic coefficients---one transcendence level above the $U_k$ and
\emph{not} reducible to the weight-one Kleinian class (their irreducibility to weight one
is \emph{conjectural}). The third-kind letter at the seed Kerr null surfaces
$\Delta=0$---whose residue is nonzero, unlike the on-shell $\psi$ where it
vanishes---adds a genus-two dilogarithm (descriptively, a hyperelliptic polylogarithm
\cite{BauneEtal2024,DHokerHiddingSchlotterer2025}), evaluated by the geometrically
convergent theta-nome series in the reproducibility package; its irreducibility and
single-valued completion remain open. Thus the on-shell reduction is a proved weight-one
closed form, while the off-shell term additionally carries the length-two iterated
integrals $W_{jk}$ and one genus-two dilogarithm (conjecturally irreducible).

\paragraph{What is reconstructible from the text alone.} Everything algebraic is fixed
by the manuscript without the repository: the frozen integrands
$\dd\varphi/\dd r=\mathcal K/\sqrt S$ with their spectral radicals and numerators;
the on-shell reduction (the polynomial identity $2A'S-AS'+2S\sum_k c_kr^k=2N_\lambda$
and its $11\times11$ linear system---no fit); all block coefficients $c_k,a_k,g_i,Q_{kj}$;
the Kleinian $\zeta,\sigma$ closed forms of $U_k,W_{jk}$ (Hermite recursion); the elementary
remainders; and the \emph{form} of the genus-two letter
$L=\log[\theta[\delta](w-e_+)/\theta[\delta](w-e_-)]$ with its odd half-characteristic
$\delta$ and the genus-one nome series $g^{(1)}$ printed in
Sec.~\ref{sec:closed}. What the text does \emph{not} reproduce by hand is the
\emph{numerical evaluation} of the genus-two $\theta$: constructing the period matrix
$\tau$ from $Y^2=\mathcal S$, its Siegel reduction, and summing the geometrically
convergent nome series (order $N{=}4$ already to $10^{-13}$) are performed by
\texttt{Sage}$+$\texttt{abelfunctions} (\texttt{RiemannTheta}) in the reproducibility
package~\cite{Rosignoli2026Code}. In short: the symbolic reduction and the closed-form
\emph{representation} are reconstructible from the PDF; the genus-two \emph{numbers}
require the archived routine.

\paragraph{Branch forms.} For the \emph{Vaidya} $v$-branch the slow datum $m$ is a metric
modulus, \emph{not} a conformal rescaling of momenta: there is no $p_r/A$ transformation
and hence no dilation term, so $\mathcal D=\Theta=m\partial_m$ is complete and
\begin{equation}
    \Theta H_v=m\,\partial_m H_v
    =m\!\left[-\frac{2p_r}{r}+\frac{\sqrt{\Ehat^2p_r^2+J^2/r^2}}{r\sqrt w}\right],
    \qquad w=\Ehat^2-f;
    \label{eq:ThetaHv}
\end{equation}
since $m$ \emph{increases}, the Vaidya correction is Eq.~\eqref{eq:adiab-exact} with
$(\lambda\Theta H_0-S_{\mathcal D})\to(S_{\mathcal D}-\lambda\Theta H_0)$ and
$-\lambda\Theta G_0\to+\lambda\Theta G_0$; the proper-time branch uses $H_\tau$ of
Eq.~\eqref{eq:Htau-vaidya}. All four cases close to $O(\varepsilon^2)$ against the true
flow: least-squares log--log slope $\approx1$ for the on-shell term alone versus
$2.05,\,2.07$ ($t,\tau$) and a second parameter set $2.01$
(Fig.~\ref{fig:phi-true-dynamic}), and $2.00$ (Vaidya, Fig.~\ref{fig:vaidya-offshell}).

\paragraph{Separatrix limit ($J\to\Jc$).} The off-shell separatrix correction is simply
the generic term at $J=\Jc$: the coefficients~\eqref{eq:ak} are polynomials in $J$, hence
finite there, while the spectral curve \emph{degenerates}, $S=(r-r_d)^2Q_4$ (double root
$r_d$, genus two $\to$ genus one). Each hyperelliptic letter collapses,
$U_k=\int r^k/\sqrt S\to V_k+r_d^{\,k}\,\Pi_{r_d}$ with $V_k=\int(\text{poly})/\sqrt{Q_4}$
(elliptic second kind, Weierstrass $\wp,\zeta$) and $\Pi_{r_d}=\int\dd r/[(r-r_d)\sqrt{Q_4}]$
(third kind); the spurious poles at $r_d$ cancel in the physical combination because
$\sum_k a_k(\Jc)r_d^{\,k}+R(\Jc,r_d)/[\Delta\,((E^2-1)r+2M)]|_{r_d}=[S/\text{deng}](r_d)=0$
(as $S(r_d)=0$), leaving the weight-two block as a genuine \emph{elliptic}
dilogarithm~\cite{BeilinsonLevin1994,BrownLevin2011} plus third kind at $\Delta=0$. Thus
the separatrix off-shell form is one degree \emph{lower} than the generic genus-two case;
verified letter by letter to machine precision (\texttt{ThakurtaMetric/tk\_sep\_genus2\_to\_elliptic.py},
\texttt{tk\_sep\_offshell\_elliptic.py}; Vaidya \texttt{VaidyaMetric/vaidya\_sep\_offshell\_elliptic.py}).
For Thakurta--Kerr the double root sits inside the ergosphere, on the physical arc, and
the dilation letter injects there a power pole $1/(r-r_d)^2$ (residue
$\propto\Delta S(r_d)$, verified against the true flow) that is \emph{not} removed by
the charge tracking $\dd J_c/\dd E$: the pole comes from the kernel $A\sim1/Q_2^2$, not
from the source. Crucially this is not a missing algebraic counterterm. Under the
adiabatic evolution the dilation ray $(\Ehat/A,\,J/A)$ leaves the separatrix locus
$J=J_c(E)$, so the double root does not merely move---it \emph{dissolves}, and the orbit
\emph{crosses} the separatrix. The off-shell non-uniformity at $r_d$ is therefore a
separatrix-crossing phenomenon \cite{Neishtadt1987,Neishtadt2013}: the first-order
adiabatic term is non-uniform there and its uniform completion requires inner/outer
matched asymptotics, outside the symbolic-coefficient reduction. For Vaidya the double
root is a branch point off the physical arc, $\mathcal D=\Theta$ carries no dilation,
and no such crossing occurs.

\begin{figure}[htbp]
    \centering
    \includegraphics[width=0.62\textwidth]{Immagini/fig_vaidya_offshell.png}
    \caption{Vaidya $v$-branch first-order correction, residual against the true
    non-autonomous flow versus $\varepsilon=M\dot m/m$ (log--log). The on-shell term
    alone (dropping the costate integral $S$) is $O(\varepsilon)$ (slope $\approx1.00$);
    the complete extended-Hamiltonian term of Eq.~\eqref{eq:adiab-exact}, adding the
    off-shell piece, closes the correction to $O(\varepsilon^2)$ (slope $\approx2.00$).
    Generated by \texttt{VaidyaMetric/vaidya\_first\_order\_offshell.py}.}
    \label{fig:vaidya-offshell}
\end{figure}

\clearpage
\bibliographystyle{iopart-num}
\bibliography{refs}

\end{document}
